% Lesson 61 — General revision — Anglais Tle A — Cours signé OnBuch+
% Programme officiel MINESEC. Compiler avec : tectonic EN61-general-revision-61.tex
\newcommand{\DOCMATIERE}{Anglais}
\newcommand{\DOCNIVEAU}{Tle A}
\newcommand{\DOCMODULE}{Chapter 5 : Media and Communication}
\newcommand{\DOCLECON}{EN61}
\newcommand{\DOCTITRE}{Lesson 61 — General revision}
% OnBuch+ — préambule « cours signés » ENRICHI (compilé avec tectonic / XeLaTeX)
% Système visuel « Pop » d'OnBuch+ : fond crème (jamais blanc), bordures encre
% épaisses, ombres « dures » décalées, titres Archivo Black en capitales.
% Version MATHÉMATIQUES : graphes (pgfplots/tikz), boîtes pédagogiques
% (définition, propriété, méthode, attention, le savais-tu, exercice résolu,
% exercice, TP, bilan), page de garde et signature OnBuch+ ancrée en bas.
%
% Macros attendues dans main.tex AVANT \input{preamble} :
%   \DOCMATIERE  \DOCNIVEAU  \DOCMODULE  \DOCLECON (numéro)  \DOCTITRE
\documentclass[11pt]{article}

\usepackage{fontspec}
\usepackage[english,french]{babel} % français = langue principale ; l'anglais via \eng{..} / textebox
\usepackage[a4paper,margin=2cm,top=3.4cm,bottom=2.8cm,headheight=1.4cm,footskip=1.3cm]{geometry}
\usepackage{amsmath,amssymb,amsfonts}
\usepackage{mathtools} % \xmapsto, \coloneqq... souvent utilisés par les modèles
\usepackage{cancel} % simplifications barrées (bilans de forces)
% \abs{z} (module, valeur absolue) et \norm{u} (norme), utilisés par les modèles
\providecommand{\abs}{}\let\abs\relax\DeclarePairedDelimiter{\abs}{\lvert}{\rvert}
\providecommand{\norm}{}\let\norm\relax\DeclarePairedDelimiter{\norm}{\lVert}{\rVert}
\usepackage[table]{xcolor} % [table] : fournit aussi \rowcolors (alternance de lignes)
\usepackage{pagecolor}
\usepackage{tikz}
\usetikzlibrary{arrows.meta,positioning,calc,shapes.geometric,shapes.misc,shapes.arrows,decorations.pathmorphing,decorations.pathreplacing,decorations.markings,patterns,shadows,mindmap,trees,fit,backgrounds,matrix,intersections,angles,quotes}
\usepackage{pgfplots}
\usepackage[version=4]{mhchem} % équations nucléaires \ce{^{235}_{92}U}
\usepackage[european,siunitx]{circuitikz} % schémas électriques
\pgfplotsset{compat=1.18}
\usepgfplotslibrary{fillbetween,groupplots} % aires entre courbes (calcul intégral)
\usepackage{enumitem}
\usepackage{fancyhdr}
% [most] charge toutes les bibliothèques tcolorbox (listings, minted, raster,
% documentation...) et gonfle inutilement la mémoire TeX sur un document long
% (~300 boîtes) : on ne charge que ce qui est réellement utilisé.
\usepackage[skins,breakable]{tcolorbox}
\usepackage{graphicx}
\usepackage{colortbl}
\usepackage{booktabs}
\usepackage{tabularx}
\usepackage{diagbox} % case d'en-tête diagonale des tableaux à double entrée
% Colonnes extensibles centrée / gauche / droite (C, L, R), souvent utilisées par les modèles
\newcolumntype{C}{>{\centering\arraybackslash}X}
\newcolumntype{L}{>{\raggedright\arraybackslash}X}
\newcolumntype{R}{>{\raggedleft\arraybackslash}X}
\usepackage{array}
\usepackage{multicol}
\usepackage{multirow}
\usepackage{siunitx}
% parse-numbers=false : accepte aussi \SI{2,5 \times 10^{-3}}{...}
\sisetup{locale=FR,output-decimal-marker={,},group-minimum-digits=4,parse-numbers=false}
\DeclareSIUnit\ohm{\ensuremath{\symup{\Omega}}} % U+2126 absent de la police
\DeclareSIUnit\division{div} % réglages d'oscilloscope (ms/div, V/div)
\DeclareSIUnit\tr{tr} % tours (vitesse de rotation, ex. \SI{1500}{\tr\per\minute})
\DeclareSIUnit\molar{M} % concentration molaire (ex. \SI{3}{\molar}, \SI{1}{\micro\molar})
\DeclareSIUnit\an{an} % année (le modèle confond parfois avec \year, un compteur TeX) — ex. \SI{5}{\kilo\meter\per\an}
\DeclareSIUnit\FCFA{FCFA} % franc CFA (ex. \SI{500}{\FCFA\per\kilo\gram})
\DeclareSIUnit\degreeBrix{\textdegree Brix} % teneur en sucre d'un sirop/jus (réfractométrie)
\DeclareSIUnit\month{mois} % le modèle utilise parfois \month, un compteur TeX, comme unité de durée
\usepackage{qrcode}
% \degree : commande fréquemment utilisée par les modèles pour un angle ou une
% température en dehors de \SI{}{} (non fournie par les paquets chargés).
\providecommand{\degree}{^{\circ}}
% \qed (fin de démonstration), utilisé par les modèles sans amsthm
\providecommand{\qed}{\hfill$\square$}
% \barre : double barre des valeurs interdites dans un tableau de variations (tkz-tab)
\providecommand{\barre}{\ensuremath{\|}}
% environnement proof (amsthm non chargé) utilisé par les modèles
\ifdefined\proof\else\newenvironment{proof}[1][Démonstration]{\par\noindent\textit{#1.}\ }{\qed\par}\fi
\usepackage{lastpage}
\usepackage{hyperref}
\hypersetup{hidelinks}

% ============================================================
% Palette « Pop » OnBuch+ — mêmes hex que la classe P (plus_ui.dart)
% ============================================================
\definecolor{poporange}{HTML}{F59321}
\definecolor{popdark}{HTML}{E07A0C}
\definecolor{popcream}{HTML}{FFF6E8}
\definecolor{popink}{HTML}{1C1714}
\definecolor{poppurple}{HTML}{7A5AE0}
\definecolor{popgreen}{HTML}{2E9E6A}
\definecolor{poppink}{HTML}{E0457B}
\definecolor{popblue}{HTML}{2F7DE1}
\definecolor{popgold}{HTML}{C98A12}
\definecolor{popmuted}{HTML}{8A7F76}
\definecolor{popline}{HTML}{EADFD2}
\definecolor{popgray}{HTML}{8A7F76}
\colorlet{popgrey}{popgray}
% Alias tolérants (le modèle nomme parfois les couleurs différemment)
\colorlet{popwhite}{white}
\colorlet{popblack}{popink}
\colorlet{popred}{poppink}
\colorlet{popyellow}{popgold}
% Teintes claires (fonds de tableaux/encadrés) — le modèle nomme parfois
% les couleurs avec un suffixe L (ex. popblueL) sans les avoir définies.
\colorlet{poporangeL}{poporange!20}
\colorlet{popdarkL}{popdark!20}
\colorlet{poppurpleL}{poppurple!20}
\colorlet{popgreenL}{popgreen!20}
\colorlet{poppinkL}{poppink!20}
\colorlet{popblueL}{popblue!20}
\colorlet{popgoldL}{popgold!20}
\colorlet{popredL}{popred!20}
\colorlet{popgrayL}{popgray!20}
% Teintes claires pour fonds de tableaux / zones
\colorlet{popblueL}{popblue!10!white}
\colorlet{poporangeL}{poporange!14!white}
\colorlet{popgreenL}{popgreen!12!white}
\colorlet{poppurpleL}{poppurple!12!white}
\colorlet{poppinkL}{poppink!12!white}
\colorlet{popgoldL}{popgold!14!white}

\pagecolor{popcream}

% ============================================================
% Typographie
% ============================================================
\defaultfontfeatures{Ligatures=TeX}
\setmainfont{PlusJakartaSans-Medium.ttf}[Path=../../../fonts/,BoldFont=PlusJakartaSans-Bold.ttf]
% Maths en sans-serif (Fira Math) avec chiffres et lettres droites de Plus
% Jakarta, pour que les formules se fondent dans le texte.
\usepackage[math-style=ISO,bold-style=ISO]{unicode-math}
\setmathfont{FiraMath-Regular.otf}[Path=../../../fonts/]
\setmathfont{PlusJakartaSans-Medium.ttf}[Path=../../../fonts/,range={up/num,up/{Latin,latin},\mathup}]
\usepackage{icomma} % virgule décimale sans espace en mode math
% Caractères mathématiques Unicode absents des polices de texte (Plus Jakarta,
% Archivo Black) : rendus par leur équivalent mathématique, en texte comme en
% formule — sinon ils s'affichent en carré vide (ex. « Arithmétique dans ℤ »).
\usepackage{newunicodechar}
\newunicodechar{ℝ}{\ensuremath{\mathbb{R}}}
\newunicodechar{ℤ}{\ensuremath{\mathbb{Z}}}
\newunicodechar{ℂ}{\ensuremath{\mathbb{C}}}
\newunicodechar{ℕ}{\ensuremath{\mathbb{N}}}
\newunicodechar{ℚ}{\ensuremath{\mathbb{Q}}}
\newunicodechar{𝔻}{\ensuremath{\mathbb{D}}}
\newunicodechar{α}{\ensuremath{\alpha}}
\newunicodechar{β}{\ensuremath{\beta}}
\newunicodechar{γ}{\ensuremath{\gamma}}
\newunicodechar{δ}{\ensuremath{\delta}}
\newunicodechar{ε}{\ensuremath{\varepsilon}}
\newunicodechar{θ}{\ensuremath{\theta}}
\newunicodechar{λ}{\ensuremath{\lambda}}
\newunicodechar{μ}{\ensuremath{\mu}}
\newunicodechar{σ}{\ensuremath{\sigma}}
\newunicodechar{τ}{\ensuremath{\tau}}
\newunicodechar{φ}{\ensuremath{\varphi}}
\newunicodechar{ψ}{\ensuremath{\psi}}
\newunicodechar{ω}{\ensuremath{\omega}}
\newunicodechar{Σ}{\ensuremath{\Sigma}}
% majuscules produites par \MakeUppercase dans les titres (α-aminés -> Α)
\newunicodechar{Α}{\ensuremath{\alpha}}
\newunicodechar{Β}{\ensuremath{\beta}}
\newunicodechar{Γ}{\ensuremath{\Gamma}}
\newunicodechar{Δ}{\ensuremath{\Delta}}
\newunicodechar{Ε}{\ensuremath{\varepsilon}}
\newunicodechar{Θ}{\ensuremath{\Theta}}
\newunicodechar{Λ}{\ensuremath{\Lambda}}
\newunicodechar{Μ}{\ensuremath{\mu}}
\newunicodechar{Τ}{\ensuremath{\tau}}
\newunicodechar{Φ}{\ensuremath{\Phi}}
\newunicodechar{Ψ}{\ensuremath{\Psi}}
\newunicodechar{Ω}{\ensuremath{\Omega}}
\newunicodechar{↦}{\ensuremath{\mapsto}}
\newunicodechar{∈}{\ensuremath{\in}}
\newunicodechar{∉}{\ensuremath{\notin}}
\newunicodechar{∧}{\ensuremath{\wedge}}
\newunicodechar{⇔}{\ensuremath{\Leftrightarrow}}
\newunicodechar{⟺}{\ensuremath{\Longleftrightarrow}}
\newunicodechar{⇒}{\ensuremath{\Rightarrow}}
\newunicodechar{⟹}{\ensuremath{\Longrightarrow}}
\newunicodechar{∘}{\ensuremath{\circ}}
\newunicodechar{∪}{\ensuremath{\cup}}
\newunicodechar{∩}{\ensuremath{\cap}}
\newunicodechar{≡}{\ensuremath{\equiv}}
\newunicodechar{⊂}{\ensuremath{\subset}}
\newunicodechar{⊥}{\ensuremath{\perp}}
\newunicodechar{∃}{\ensuremath{\exists}}
\newunicodechar{∀}{\ensuremath{\forall}}
\newunicodechar{∛}{\ensuremath{\sqrt[3]{\,}}}
\newunicodechar{∜}{\ensuremath{\sqrt[4]{\,}}}
\newunicodechar{⃗}{\ensuremath{{}^{\to}}}
\newunicodechar{ⁿ}{\ifmmode^{n}\else\textsuperscript{\ensuremath{n}}\fi}
\newunicodechar{ˣ}{\ifmmode^{x}\else\textsuperscript{\ensuremath{x}}\fi}
\newunicodechar{ᵇ}{\ifmmode^{b}\else\textsuperscript{\ensuremath{b}}\fi}
\newunicodechar{ʳ}{\ifmmode^{r}\else\textsuperscript{\ensuremath{r}}\fi}
\newunicodechar{ᵘ}{\ifmmode^{u}\else\textsuperscript{\ensuremath{u}}\fi}
\newunicodechar{ʸ}{\ifmmode^{y}\else\textsuperscript{\ensuremath{y}}\fi}
\newunicodechar{ᵃ}{\ifmmode^{a}\else\textsuperscript{\ensuremath{a}}\fi}
\newunicodechar{ᵉ}{\ifmmode^{e}\else\textsuperscript{\ensuremath{e}}\fi}
\newunicodechar{ᵏ}{\ifmmode^{k}\else\textsuperscript{\ensuremath{k}}\fi}
\newunicodechar{ᵖ}{\ifmmode^{p}\else\textsuperscript{\ensuremath{p}}\fi}
\newunicodechar{ᴺ}{\ifmmode^{N}\else\textsuperscript{\ensuremath{N}}\fi}
\newunicodechar{ᶜ}{\ifmmode^{c}\else\textsuperscript{\ensuremath{c}}\fi}
\newunicodechar{ʰ}{\ifmmode^{h}\else\textsuperscript{\ensuremath{h}}\fi}
\newunicodechar{ᵠ}{\ifmmode^{q}\else\textsuperscript{\ensuremath{q}}\fi}
\newunicodechar{ₙ}{\ifmmode_{n}\else\textsubscript{\ensuremath{n}}\fi}
\newunicodechar{ₐ}{\ifmmode_{a}\else\textsubscript{\ensuremath{a}}\fi}
\newunicodechar{ᵢ}{\ifmmode_{i}\else\textsubscript{\ensuremath{i}}\fi}
\newunicodechar{ᵣ}{\ifmmode_{r}\else\textsubscript{\ensuremath{r}}\fi}
\newunicodechar{ₖ}{\ifmmode_{k}\else\textsubscript{\ensuremath{k}}\fi}
\newunicodechar{ᵦ}{\ifmmode_{\beta}\else\textsubscript{\ensuremath{\beta}}\fi}
\newunicodechar{ₓ}{\ifmmode_{x}\else\textsubscript{\ensuremath{x}}\fi}
\newfontfamily\popdisplay{ArchivoBlack-Regular.ttf}[Path=../../../fonts/]
\newfontfamily\poplogo{SpaceGrotesk-Bold.ttf}[Path=../../../fonts/]
\newfontfamily\popsemi{PlusJakartaSans-SemiBold.ttf}[Path=../../../fonts/]
\newfontfamily\popxbold{PlusJakartaSans-ExtraBold.ttf}[Path=../../../fonts/]
\newfontfamily\popxboldls{PlusJakartaSans-ExtraBold.ttf}[Path=../../../fonts/,LetterSpace=3.0]

\setlist[enumerate]{leftmargin=1.7em,itemsep=5pt,topsep=4pt}
\setlist[itemize]{leftmargin=1.7em,itemsep=4pt,topsep=4pt,label={\color{poporange}\textbullet}}
\setlength{\parindent}{0pt}
\setlength{\parskip}{5pt}
\renewcommand{\arraystretch}{1.25}

% pgfplots : style Pop par défaut
\pgfplotsset{
  every axis/.append style={axis line style={popink,line width=1pt},tick style={popink},
    grid style={popline,line width=0.5pt},label style={font=\small},tick label style={font=\footnotesize}},
  popcurve/.style={poporange,line width=1.8pt,no markers,smooth},
}
% Même style disponible dans un \draw TikZ simple (hors pgfplots)
\tikzset{popcurve/.style={poporange,line width=1.8pt}}
\tikzset{popgrid/.style={popline,very thin}}
\colorlet{popcurve}{poporange} % le nom du style est parfois utilisé comme couleur

% ============================================================
% Pill / squiggle
% ============================================================
\newcommand{\poppill}[3]{%
  \tikz[baseline=-0.55ex]\node[draw=popink,line width=1.2pt,rounded corners=2.6pt,%
    fill=#2,inner xsep=6.5pt,inner ysep=3pt]{%
    \popxboldls\fontsize{7}{7}\selectfont\color{#3}\MakeUppercase{#1}};%
}
\newcommand{\popsquiggle}[1]{%
  \tikz[baseline=-0.4ex]{
    \draw[#1,line width=2.6pt,line cap=round]
      plot [smooth] coordinates {(0,0.09) (0.34,0.01) (0.68,0.21) (1.02,0.01) (1.36,0.10)};
  }%
}
% Mot-clé surligné (marqueur orange sous le texte)
\newcommand{\cle}[1]{{\popxbold\color{popdark}#1}}

% ============================================================
% En-tête / pied
% ============================================================
\pagestyle{fancy}
\fancyhf{}
\renewcommand{\headrulewidth}{0pt}
\renewcommand{\footrulewidth}{0pt}
\lhead{\raisebox{-0.15ex}{\poplogo\fontsize{13}{13}\selectfont\color{popink}OnBuch\color{poporange}+}}
\rhead{\poppill{\DOCMATIERE\ \textbullet\ \DOCNIVEAU\ \textbullet\ Leçon \DOCLECON}{white}{popink}}
\lfoot{\popxbold\fontsize{7.6}{7.6}\selectfont\color{popmuted}\MakeUppercase{Cours signé OnBuch+}\ \textbullet\ {\popsemi\DOCTITRE}}
\rfoot{\tikz[baseline=-0.6ex]\node[rounded corners=8.5pt,fill=popink,minimum height=17pt,minimum width=17pt,inner xsep=5pt,inner sep=0]{\popxbold\fontsize{8}{8}\selectfont\color{popcream}\thepage\,/\,\pageref*{LastPage}};}

% ============================================================
% Titres
% ============================================================
\newcommand{\chaptitre}[2]{%
  \begin{tcolorbox}[enhanced,colback=poporange,colframe=popink,boxrule=2.6pt,arc=11pt,
      drop shadow=popink,left=15pt,right=15pt,top=12pt,bottom=11pt,before skip=3pt,after skip=18pt]
    \poppill{#2}{popink}{popcream}\par\vspace{9pt}
    {\raggedright\hyphenpenalty=10000\exhyphenpenalty=10000\popdisplay\fontsize{22}{24}\selectfont\color{popcream}\MakeUppercase{#1}\par}
  \end{tcolorbox}
}
% Section numérotée : pastille orange avec le numéro + titre Archivo
\newcounter{popsec}
\newcommand{\coursec}[1]{%
  \stepcounter{popsec}\setcounter{popsub}{0}%
  \par\vspace{16pt}\noindent
  \begin{minipage}{\linewidth}
  \tikz[baseline=-0.7ex]\node[draw=popink,line width=1.6pt,fill=poporange,rounded corners=4pt,minimum size=22pt,inner sep=0]{\popdisplay\fontsize{12}{12}\selectfont\color{popcream}\thepopsec};%
  \hspace{8pt}{\popdisplay\fontsize{15}{17}\selectfont\color{popink}\MakeUppercase{#1}}\par
  \vspace{4pt}\noindent{\color{poporange}\rule{36pt}{3.6pt}}
  \end{minipage}\par\nopagebreak\vspace{8pt}
}
\newcounter{popsub}
\newcommand{\courssub}[1]{%
  \stepcounter{popsub}%
  \par\vspace{9pt}\noindent{\popxbold\fontsize{11.5}{13}\selectfont\color{popink}{\color{poporange}\thepopsec.\thepopsub}\hspace{6pt}#1}\par\nopagebreak\vspace{4pt}
}

% ============================================================
% Boîtes pédagogiques — même recette : fond blanc, bordure épaisse colorée,
% ombre dure encre, badge plein. Titre optionnel [#1].
% ============================================================
\tcbset{popbase/.style={breakable,enhanced,colback=white,boxrule=2.3pt,arc=9pt,
  shadow={3pt}{-3pt}{0pt}{popink},left=14pt,right=14pt,top=11pt,bottom=11pt,before skip=11pt,after skip=15pt,
  fontupper=\popsemi\fontsize{10.6}{14.5}\selectfont\color{popink},
  fontlower=\popsemi\fontsize{10.6}{14.5}\selectfont\color{popink}}}
% Coche dessinée (le glyphe ✓ n'existe pas dans les polices du document)
\newcommand{\popcheck}[1]{\tikz[baseline=-0.5ex]\draw[#1,line width=1.6pt,line cap=round,line join=round](-0.13,0.0)--(-0.04,-0.09)--(0.13,0.1);}
\providecommand{\coche}{\popcheck{popgreen}}
\providecommand{\resultat}[1]{\par\smallskip\noindent\fcolorbox{popgreen}{popgreenL}{\parbox{\dimexpr\linewidth-2\fboxsep-2\fboxrule\relax}{\textbf{Résultat.} #1}}\par\smallskip}
\newcommand{\popboxhead}[3]{\poppill{#1}{#2}{white}\ifx\relax#3\relax\else\hspace{6pt}{\popxbold\fontsize{10.6}{12}\selectfont\color{popink}#3}\fi\par\vspace{7pt}}

\newtcolorbox{aretenir}[1][]{popbase,colframe=popgreen,before upper={\popboxhead{À retenir}{popgreen}{#1}}}
% Texte en anglais (document de lecture, dialogue, extrait) : cadre
% d'étude. Un titre optionnel (source, auteur) s'affiche dans l'en-tête.
\newtcolorbox{textebox}[1][]{popbase,colframe=popblue,before upper={\popboxhead{Text}{popblue}{#1}\begin{otherlanguage}{english}},after upper={\end{otherlanguage}}}
% Marques ✓ / ✗ (forme correcte / fautive) souvent utilisées par les modèles
\providecommand{\cross}{\textcolor{poppink}{\ensuremath{\times}}}
\providecommand{\xmark}{\cross}\providecommand{\crossmark}{\cross}\providecommand{\cmark}{\ensuremath{\checkmark}}
% Ligne de réponse / justification à compléter (inventée par les modèles)
\providecommand{\justification}[1]{\par\noindent\textit{Justification~:} \dotfill\par}
% \textipa (package tipa non chargé) : la police ne couvre pas l'API -> simple passage
\providecommand{\textipa}[1]{#1}
% Soulignement (ulem non chargé) : \uline, \ul
\providecommand{\uline}[1]{\underline{#1}}\providecommand{\ul}[1]{\underline{#1}}
% \glossaire{\item ...} inventée par les modèles : liste de glossaire
\providecommand{\glossaire}[1]{\begin{itemize}#1\end{itemize}}
% \alert (beamer) inventée par les modèles : texte mis en évidence
\providecommand{\alert}[1]{\textbf{\textcolor{poppink}{#1}}}
% variantes ASCII d'ordinaux écrites en macro
\providecommand{\iere}{\textsuperscript{re}}\providecommand{\ieme}{\textsuperscript{e}}\providecommand{\ere}{\textsuperscript{re}}\providecommand{\eme}{\textsuperscript{e}}\providecommand{\er}{\textsuperscript{er}}\providecommand{\ie}{\textsuperscript{e}}
% Ordinaux français écrits en macro par les modèles : 1\ière, 3\ième
\newcommand{\ière}{\textsuperscript{re}}\newcommand{\ième}{\textsuperscript{e}}
% \exemple (inventée par les modèles) : étiquette « Exemple : »
\providecommand{\exemple}{\textbf{Exemple~:}\ }
% \square utilisable aussi hors mode math (cases à cocher)
\AtBeginDocument{\let\squareMath\square\renewcommand{\square}{\ensuremath{\squareMath}}}
% Mot ou phrase en anglais à l'intérieur d'un paragraphe français
\newcommand{\eng}[1]{\ifmmode\text{\textit{#1}}\else\begingroup\selectlanguage{english}\itshape #1\endgroup\fi}
\let\esp\eng % alias toléré
% flèches d'intonation inventées par les modèles
\providecommand{\textsearrow}{}\renewcommand{\textsearrow}{\ensuremath{\searrow}}\providecommand{\textnearrow}{}\renewcommand{\textnearrow}{\ensuremath{\nearrow}}
% \fr{..} : mot français (inventé par les modèles)
\providecommand{\fr}[1]{\textit{#1}}
\newtcolorbox{exemplebox}[1][]{popbase,colframe=poporange,before upper={\popboxhead{Exemple}{poporange}{#1}}}
\newtcolorbox{definition}[1][]{popbase,colframe=popblue,before upper={\popboxhead{Définition}{popblue}{#1}}}
\newtcolorbox{propriete}[1][]{popbase,colframe=poppurple,before upper={\popboxhead{Propriété}{poppurple}{#1}}}
\newtcolorbox{methode}[1][]{popbase,colframe=popgold,before upper={\popboxhead{Méthode}{popgold}{#1}}}
\newtcolorbox{attention}[1][]{popbase,colframe=poppink,before upper={\popboxhead{Attention}{poppink}{#1}}}
\newtcolorbox{experience}[1][]{popbase,colframe=popblue,colback=popblueL,before upper={\popboxhead{Expérience}{popblue}{#1}}}
% « Le savais-tu ? » : carte violette pleine (contexte, culture, Cameroun)
\newtcolorbox{savaistu}[1][]{popbase,colback=poppurple,colframe=popink,
  fontupper=\popsemi\fontsize{10.4}{14.2}\selectfont\color{white},
  before upper={\poppill{Le savais-tu\,?}{white}{poppurple}\ifx\relax#1\relax\else\hspace{6pt}{\popxbold\color{white}#1}\fi\par\vspace{7pt}}}
% Exercice résolu : énoncé en haut, \tcblower puis la solution sur fond crème
\newcounter{popexo}\newcounter{popexr}
\newtcolorbox{exoresolu}[1][]{popbase,colframe=poporange,colbacklower=popcream,
  segmentation style={popink,line width=1.2pt,dashed},
  before upper={\stepcounter{popexr}\popboxhead{Exercice résolu \thepopexr}{poporange}{#1}},
  before lower={\poppill{Solution}{popink}{popcream}\par\vspace{6pt}}}
% Exercice (non résolu en place) — la correction est dans la partie Corrigés
\newtcolorbox{exercice}[1][]{popbase,colframe=popink,
  before upper={\stepcounter{popexo}\popboxhead{Exercice \thepopexo}{popink}{#1}}}
% Corrigé
\newtcolorbox{corrige}[1][]{popbase,colframe=popgreen,colback=popgreenL,
  before upper={\popboxhead{Corrigé}{popgreen}{#1}}}
% Objectifs / prérequis / bilan
\newtcolorbox{objectifs}{popbase,colframe=popink,colback=poporangeL,
  before upper={\popboxhead{Objectifs de la leçon}{popink}{}}}
\newtcolorbox{prerequis}{popbase,colframe=popink,colback=popblueL,
  before upper={\popboxhead{Prérequis}{popblue}{}}}
\newtcolorbox{bilan}{popbase,colframe=popink,colback=white,boxrule=2.8pt,
  before upper={\popboxhead{Fiche bilan}{poporange}{L'essentiel à connaître pour l'examen}}}
% Figure encadrée avec légende
\newtcolorbox{popfigure}[1][]{enhanced,breakable=false,colback=white,colframe=popline,boxrule=1.4pt,arc=8pt,
  left=10pt,right=10pt,top=10pt,bottom=8pt,before skip=10pt,after skip=12pt,halign=center,
  lower separated=false,halign lower=center,
  fontlower=\popsemi\fontsize{9}{11.5}\selectfont\color{popmuted},
  #1}
\newcommand{\legende}[1]{\par\vspace{6pt}{\popsemi\fontsize{9}{11.5}\selectfont\color{popmuted}#1}}

% ============================================================
% Signature OnBuch+
% ============================================================
\newcommand{\poplogoinline}[1]{{\poplogo\fontsize{#1}{#1}\selectfont\color{popink}OnBuch\color{poporange}+}}
\newcommand{\signatureonbuch}{%
  \begin{tcolorbox}[enhanced,colback=popink,colframe=popink,boxrule=2.2pt,arc=10pt,
      drop shadow=poporange,left=14pt,right=14pt,top=10pt,bottom=10pt,before skip=0pt,after skip=0pt]
    \tikz[baseline=-0.7ex]\node[circle,fill=popgreen,draw=popcream,line width=1.3pt,minimum size=22pt,inner sep=0]{\popcheck{white}};%
    \hspace{9pt}\begin{minipage}[c]{0.62\linewidth}
      {\popxbold\fontsize{12}{13}\selectfont\color{popcream}Signé\ \textbullet\ {\poplogo OnBuch\color{poporange}+}}\par\vspace{2pt}
      {\raggedright\popsemi\fontsize{8}{10.5}\selectfont\color{popline}\MakeUppercase{Équipe pédagogique OnBuch+}\\\MakeUppercase{Conforme au programme officiel MINESEC}\par}
    \end{minipage}\hfill
    {\poplogo\fontsize{22}{22}\selectfont\color{popcream}OnBuch\color{poporange}+}
  \end{tcolorbox}%
}

% ============================================================
% Page de garde
% #1 titre · #2 matière · #3 niveau · #4 module · #5 n° leçon · #6 durée ·
% #7 description / objectifs (texte ou itemize)
% ============================================================
\newcommand{\pagegarde}[7]{%
  \thispagestyle{empty}
  \newgeometry{a4paper,margin=1.7cm,top=1.6cm,bottom=1.4cm}%
  \begin{tikzpicture}[remember picture,overlay]
    \fill[poporange!18] ([xshift=-3.2cm,yshift=-3.6cm]current page.north east) circle (4.6cm);
    \fill[poppurple!14] ([xshift=2.4cm,yshift=7.6cm]current page.south west) circle (3.8cm);
  \end{tikzpicture}%
  {\poplogo\fontsize{30}{30}\selectfont\color{popink}OnBuch\color{poporange}+}\hfill
  \raisebox{0.6ex}{\poppill{Cours signé \textbullet\ Premium}{popink}{popcream}}\par
  \vspace{1pt}\popsquiggle{poppurple}\par
  \vspace{8pt}
  \begin{tcolorbox}[enhanced,colback=poporange,colframe=popink,boxrule=2.8pt,arc=13pt,
      drop shadow=popink,left=18pt,right=18pt,top=12pt,bottom=13pt,after skip=10pt]
    \poppill{#2 \textbullet\ #3}{popink}{popcream}\hspace{5pt}\poppill{#4}{white}{popink}\par\vspace{8pt}
    {\popdisplay\fontsize{14}{14}\selectfont\color{popink}LEÇON #5}\par\vspace{4pt}
    {\raggedright\hyphenpenalty=10000\exhyphenpenalty=10000\popdisplay\fontsize{25}{27}\selectfont\color{popcream}\MakeUppercase{#1}\par}
    \vspace{8pt}
    {\popxbold\fontsize{9.5}{9.5}\selectfont\color{popink}\MakeUppercase{Durée conseillée : #6}}
  \end{tcolorbox}
  \begin{tcolorbox}[enhanced,colback=white,colframe=popink,boxrule=2.2pt,arc=10pt,
      drop shadow=popink,left=16pt,right=16pt,top=10pt,bottom=10pt,after skip=8pt]
    \poppill{Dans cette leçon}{poppurple}{white}\par\vspace{5pt}
    \popsemi\fontsize{9.3}{12.6}\selectfont\color{popink}#7
  \end{tcolorbox}
  \noindent\begin{minipage}[t]{0.487\linewidth}
    \begin{tcolorbox}[enhanced,colback=popgreen,colframe=popink,boxrule=2.2pt,arc=10pt,
        drop shadow=popink,left=11pt,right=11pt,top=10pt,bottom=10pt,height=3.1cm,valign=center]
      \tcbox[colback=white,colframe=popink,boxrule=1.3pt,arc=5pt,boxsep=0pt,nobeforeafter,
        left=3pt,right=3pt,top=3pt,bottom=3pt,valign=center]{\qrcode[height=1.9cm]{https://play.google.com/store/apps/details?id=com.luvvix.onbuch&hl=fr}}%
      \hspace{8pt}\begin{minipage}[c]{0.5\linewidth}
        {\popdisplay\fontsize{11}{12.5}\selectfont\color{white}\MakeUppercase{Télécharge OnBuch}}\par\vspace{4pt}
        {\raggedright\popsemi\fontsize{8.4}{11}\selectfont\color{white}Gratuit \textbullet\ Google Play\\Tuteur IA, annales, concours, résultats\par}
      \end{minipage}
    \end{tcolorbox}
  \end{minipage}\hfill
  \begin{minipage}[t]{0.487\linewidth}
    \begin{tcolorbox}[enhanced,colback=poppurple,colframe=popink,boxrule=2.2pt,arc=10pt,
        drop shadow=popink,left=11pt,right=11pt,top=10pt,bottom=10pt,height=3.1cm,valign=center]
      \tikz[baseline=-0.7ex]\node[circle,fill=white,draw=popink,line width=1.3pt,minimum size=1.6cm,inner sep=0]{\popdisplay\fontsize{24}{24}\selectfont\color{poppurple}+};%
      \hspace{8pt}\begin{minipage}[c]{0.6\linewidth}
        {\popdisplay\fontsize{11}{12.5}\selectfont\color{white}\MakeUppercase{Passe à OnBuch+}}\par\vspace{4pt}
        {\raggedright\popsemi\fontsize{8.4}{11}\selectfont\color{white}Tous les cours signés, Tuteur IA illimité, fiches PDF — onglet OnBuch+ de l'app\par}
      \end{minipage}
    \end{tcolorbox}
  \end{minipage}
  \vspace{8pt}
  \popcontenu
  \vfill
  \signatureonbuch
  \restoregeometry
  \newpage
}

% Rangée « Ce cours contient » (page de garde)
\newcommand{\poptile}[3]{% #1 couleur #2 gros texte #3 légende
  \begin{tcolorbox}[enhanced,colback=#1,colframe=popink,boxrule=2pt,arc=9pt,shadow={2.5pt}{-2.5pt}{0pt}{popink},
      left=6pt,right=6pt,top=9pt,bottom=9pt,halign=center,valign=center,height=1.75cm,width=\linewidth,nobeforeafter]
    {\popdisplay\fontsize{12.5}{13}\selectfont\color{white}\MakeUppercase{#2}}\par\vspace{3pt}
    {\centering\popsemi\fontsize{7.8}{9.5}\selectfont\color{white}#3\par}
  \end{tcolorbox}}
\newcommand{\popcontenu}{%
  \par\noindent{\popxboldls\fontsize{7.5}{7.5}\selectfont\color{popmuted}CE COURS SIGNÉ CONTIENT}\par\vspace{7pt}
  \noindent\begin{minipage}[t]{0.235\linewidth}\poptile{popblue}{Cours}{complet, illustré, pas à pas}\end{minipage}\hfill
  \begin{minipage}[t]{0.235\linewidth}\poptile{poporange}{Méthodes}{et exercices résolus}\end{minipage}\hfill
  \begin{minipage}[t]{0.235\linewidth}\poptile{poppink}{Exercices}{type examen + corrigés}\end{minipage}\hfill
  \begin{minipage}[t]{0.235\linewidth}\poptile{popgreen}{Fiche bilan}{l'essentiel à retenir}\end{minipage}\par}

% Sommaire
\newtcolorbox{sommaire}{enhanced,colback=white,colframe=popink,boxrule=2.2pt,arc=10pt,
  drop shadow=popink,left=18pt,right=18pt,top=13pt,bottom=13pt,before skip=6pt,after skip=20pt,
  before upper={\poppill{Sommaire}{poppurple}{white}\par\vspace{9pt}\popsemi\fontsize{10.6}{14.5}\selectfont\color{popink}}}

% Signature de fin de cours — poussée tout en bas de la dernière page
\newcommand{\findecours}{%
  \par\vfill\nopagebreak
  \begin{center}\popsquiggle{poporange}\end{center}\vspace{4pt}
  \signatureonbuch
}
\AtBeginDocument{\def\llbracket{\mathopen{[\mkern-3.5mu[}}\def\rrbracket{\mathclose{]\mkern-3.5mu]}}} % intervalles d'entiers (glyphe ⟦ absent de la police)
% Glyphes absents de Fira Math : redéfinis à partir de glyphes disponibles
\AtBeginDocument{%
  \def\setminus{\mathbin{\backslash}}\let\smallsetminus\setminus
  \def\vdots{\mathord{\vbox{\baselineskip3.2pt\lineskiplimit0pt\kern4pt\hbox{.}\hbox{.}\hbox{.}}}}%
  \def\ddots{\mathinner{\mkern1mu\raise6.5pt\hbox{.}\mkern2mu\raise3.6pt\hbox{.}\mkern2mu\raise0.7pt\hbox{.}\mkern1mu}}%
  \def\triangle{\mathord{\scalebox{0.9}{$\symup{\Delta}$}}}%
  \def\nabla{\mathord{\raisebox{\depth}{\scalebox{1}[-1]{$\symup{\Delta}$}}}}%
  \def\checkmark{\popcheck{popdark}}%
  \def\star{\mathbin{\ast}}\def\bigstar{\mathord{\ast}}%
  \def\diamond{\mathbin{\scalebox{0.6}{$\Diamond$}}}%
  \def\bigcup{\mathop{\mathchoice{\raisebox{-0.3em}{\scalebox{1.7}{$\cup$}}}{\scalebox{1.25}{$\cup$}}{\cup}{\cup}}\displaylimits}%
  \def\bigcap{\mathop{\mathchoice{\raisebox{-0.3em}{\scalebox{1.7}{$\cap$}}}{\scalebox{1.25}{$\cap$}}{\cap}{\cap}}\displaylimits}%
  \def\lesseqqgtr{\mathrel{\lessgtr}}\def\bigoplus{\mathop{\oplus}\displaylimits}%
}
\providecommand{\ssi}{si et seulement si} % abréviation fréquente
\AtBeginDocument{\providecommand{\wideparen}{\overparen}} % arc de cercle

%% FIGFIX-BEGIN v4 — correctifs globaux des figures (docs/onbuch-generation/tools/figfix)
\usepackage{adjustbox}\usepackage{etoolbox}
% toute figure TikZ/pgfplots plus large que la ligne est réduite à la largeur disponible
\BeforeBeginEnvironment{tikzpicture}{\begin{adjustbox}{max width=\linewidth}}
\AfterEndEnvironment{tikzpicture}{\end{adjustbox}}
% légendes pgfplots lisibles par-dessus les courbes
\pgfplotsset{every axis/.append style={legend style={fill=white,fill opacity=0.92,text opacity=1,draw=black!15,inner sep=3pt}}}
% étiquettes d'arêtes (syntaxe edge["..."]) avec fond blanc
\tikzset{every edge quotes/.append style={fill=white,inner sep=1.2pt}}
%% FIGFIX-END
\begin{document}

\pagegarde{Lesson 61 — General revision}{Anglais}{Tle A}{Chapter 5 : Media and Communication}{EN61}{2 h}{Leçon de révision générale de Terminale : synthèse des temps verbaux, des structures clés et du lexique de l'année autour du thème des médias, avec entraînement progressif de type Bac (compréhension, grammaire, vocabulaire, composition).
\vspace{4pt}
\begin{multicols}{2}\small
\begin{itemize}
\item À la fin de cette leçon, je sais mobiliser les principaux temps verbaux anglais (present, past, future, perfect) dans un texte de type examen
\item À la fin de cette leçon, je sais transformer des phrases (passif, discours rapporté, comparatifs, conditionnel) sans erreur
\item À la fin de cette leçon, je sais utiliser le lexique des médias et de la communication dans un contexte approprié
\item À la fin de cette leçon, je sais repérer et corriger les erreurs fréquentes des francophones (faux amis, concordance des temps, a/an, -ing ou infinitif)
\item À la fin de cette leçon, je sais répondre efficacement à des questions de compréhension écrite (vrai/faux, QCM, réponses courtes)
\item À la fin de cette leçon, je sais rédiger une composition structurée (introduction, développement, conclusion) avec des connecteurs variés
\end{itemize}\end{multicols}}

\begin{sommaire}
\begin{multicols}{2}
\begin{enumerate}
\item Reading comprehension — exam practice
\item Grammar synthesis — verb tenses
\item Grammar synthesis — key structures
\item Vocabulary consolidation — Media and Communication
\item Writing — composition de type Bac
\item Exam strategy and final training
\item Activité d'intégration
\item Méthodes et exercices résolus
\item Exercices
\item Corrigés des exercices
\item Fiche bilan
\end{enumerate}
\end{multicols}
\end{sommaire}

\begin{objectifs}
\begin{itemize}
\item À la fin de cette leçon, je sais mobiliser les principaux temps verbaux anglais (present, past, future, perfect) dans un texte de type examen
\item À la fin de cette leçon, je sais transformer des phrases (passif, discours rapporté, comparatifs, conditionnel) sans erreur
\item À la fin de cette leçon, je sais utiliser le lexique des médias et de la communication dans un contexte approprié
\item À la fin de cette leçon, je sais repérer et corriger les erreurs fréquentes des francophones (faux amis, concordance des temps, a/an, -ing ou infinitif)
\item À la fin de cette leçon, je sais répondre efficacement à des questions de compréhension écrite (vrai/faux, QCM, réponses courtes)
\item À la fin de cette leçon, je sais rédiger une composition structurée (introduction, développement, conclusion) avec des connecteurs variés
\item À la fin de cette leçon, je sais gérer mon temps et ma stratégie le jour de l'examen
\end{itemize}
\end{objectifs}

\begin{prerequis}
\begin{itemize}
\item Les temps verbaux étudiés en Première et Terminale (present simple/continuous, past simple/continuous, present perfect, future forms)
\item Le discours rapporté et la voix passive
\item Le lexique des médias (press, broadcast, social media) vu dans le chapitre 5
\item Les connecteurs logiques et la structure d'une composition argumentative
\item Les comparatifs, superlatifs et le conditionnel
\end{itemize}
\end{prerequis}

\newpage

\coursec{Reading comprehension — exam practice}

The BEPC and the Probatoire are behind you; now the Baccalauréat is only weeks away. Imagine a student in Bamenda who listens to the BBC every morning, reads \emph{The Guardian} online, and chats with friends on WhatsApp in English — without realising it, she is revising everything she learnt this year: tenses, reported speech, the passive voice, media vocabulary. This final lesson is your last training session before the exam: we will pull together grammar, vocabulary, comprehension and composition skills through exam-style practice on the theme of Media and Communication.

\textbf{Problématique :} comment lire un texte de presse en anglais de manière rapide et stratégique, et répondre à toutes les questions de compréhension avec la précision exigée au Baccalauréat ?

\courssub{Reading the text — first contact}

\begin{textebox}[Social Media and the Cameroonian Youth]
In Yaoundé and Douala, most teenagers now get their news from Facebook and WhatsApp rather than from television or newspapers. According to a recent survey, 70\% of students aged 15 to 19 spend more than three hours a day online. While social media help young people to communicate and to learn, they also spread fake news very quickly. A photo of a “free scholarship” or a voice message about a “miracle cure” can travel from Buea to Bamenda in a few minutes, and many users share such messages without checking them.

Teachers and parents are worried, but many schools have started media-literacy lessons to teach students how to check information before sharing it. In these lessons, students learn to ask simple questions: Who wrote this? When was it published? Can I find the same story on a reliable website? Some schools have even created “press clubs” where teenagers write short articles for a school magazine.

Experts believe that critical thinking is the best protection against disinformation. “Young people are not the problem,” says one teacher in Douala. “They are the solution — if we give them the right tools.” The Ministry of Secondary Education now encourages every school to include media education in its programme. For the Cameroonian youth, learning to read critically has become as important as learning to read.
\end{textebox}

\begin{definition}[Glossaire du texte]
\begin{center}
\begin{tabularx}{\linewidth}{|l|l|X|}
\hline
\rowcolor{popblueL} \textbf{English} & \textbf{Nature} & \textbf{Français} \\
\hline
a survey & noun & une enquête, un sondage \\
\hline
reliable & adjective & fiable, digne de confiance \\
\hline
fake news & noun & fausses informations, infox \\
\hline
to spread & verb & répandre, diffuser (irrégulier : spread–spread–spread) \\
\hline
to check & verb & vérifier, contrôler (une information) \\
\hline
disinformation & noun & désinformation \\
\hline
critical thinking & noun & esprit critique \\
\hline
\end{tabularx}
\end{center}
\end{definition}

\begin{exemplebox}[Observation grammaticale dans le texte]
\eng{Many schools have started media-literacy lessons.} = « Beaucoup d'écoles ont commencé des cours d'éducation aux médias. »

On reconnaît le \cle{present perfect} (\emph{have/has} + participe passé) : une action passée dont le bilan est visible dans le présent (les cours existent maintenant). Le français utilise ici un passé composé, mais l'anglais insiste sur le lien avec le présent : on ne peut pas ajouter une date passée précise. Comparez :

\eng{Many schools started media-literacy lessons in 2020.} = « Beaucoup d'écoles ont commencé ces cours en 2020. » (past simple + date précise)

Même opposition dans la dernière phrase du texte : \eng{Learning to read critically has become as important as learning to read.} = « Apprendre à lire avec un esprit critique est devenu aussi important qu'apprendre à lire. »
\end{exemplebox}

\courssub{The reading method — la méthode en cinq étapes}

\begin{methode}[Réussir la compréhension écrite au Bac]
\begin{enumerate}
\item \textbf{Lire d'abord les questions} (avant le texte) : elles indiquent ce qu'il faut chercher.
\item \textbf{Repérer les mots-clés} dans chaque question : noms propres, chiffres, verbes d'action.
\item \textbf{Première lecture} du texte : viser le sens global, sans dictionnaire, sans bloquer sur un mot inconnu.
\item \textbf{Deuxième lecture} : repérer les détails et souligner la ligne qui contient chaque réponse.
\item \textbf{Rédiger les réponses} : \textbf{ne jamais répondre de mémoire}, toujours justifier en citant une ligne du texte, et \textbf{reprendre la formulation de la question} dans la réponse (\eng{Why are teachers worried?} $\rightarrow$ \eng{Teachers are worried because...}).
\end{enumerate}
\end{methode}

\begin{popfigure}
\begin{tikzpicture}[node distance=0.55cm,
  etape/.style={rectangle, rounded corners, draw=popblue, fill=popblueL, text width=4.6cm, align=center, font=\small, inner sep=5pt},
  fleche/.style={-{Stealth[length=3mm]}, thick, poporange}]
\node[etape, align=center] (q){\textbf{1. Questions}\\ lire les questions d'abord};
\node[etape, below=of q, align=center] (kw){\textbf{2. Mots-clés}\\ souligner noms, chiffres, verbes};
\node[etape, below=of kw, align=center] (r1){\textbf{3. Première lecture}\\ sens global};
\node[etape, below=of r1, align=center] (r2){\textbf{4. Deuxième lecture}\\ détails, souligner les preuves};
\node[etape, below=of r2, align=center] (ans){\textbf{5. Réponses justifiées}\\ citer une ligne du texte};
\draw[fleche] (q) -- (kw);
\draw[fleche] (kw) -- (r1);
\draw[fleche] (r1) -- (r2);
\draw[fleche] (r2) -- (ans);
\end{tikzpicture}
\legende{Organigramme de la méthode de lecture : cinq étapes, dans l'ordre, sans exception.}
\end{popfigure}

\begin{savaistu}[Le saviez-vous ?]
Au Baccalauréat, la section \emph{Reading comprehension} vaut généralement 8 à 10 points sur 20. C'est la partie la plus « rentable » de l'épreuve : le texte contient toutes les réponses, il suffit de savoir les trouver. Un candidat méthodique peut y gagner la quasi-totalité des points, même avec un niveau d'anglais moyen.
\end{savaistu}

\courssub{The four question types — les quatre types de questions}

\textbf{Type 1 — True or False, with justification.} Vous devez écrire \eng{True} ou \eng{False} \emph{et} recopier la phrase du texte qui le prouve. Sans justification, la réponse vaut zéro, même si elle est correcte.

\textbf{Type 2 — Multiple choice (QCM).} Quatre options (A, B, C, D). Éliminez d'abord les deux options absurdes, puis choisissez entre les deux restantes en relisant le passage concerné.

\textbf{Type 3 — Short answers (who, what, why, how).} Répondez par une phrase complète qui reprend les mots de la question. \eng{Who} interroge sur une personne, \eng{what} sur une chose, \eng{why} sur une cause (réponse avec \eng{because}), \eng{how} sur une manière ou une quantité.

\textbf{Type 4 — Vocabulary in context.} On vous demande de trouver \emph{dans le texte} un mot synonyme d'un mot donné, ou son contraire. La réponse doit être un mot du texte, pas un mot de votre mémoire.

\begin{exoresolu}[Un exemple corrigé de chaque type — True/False]
\textbf{Énoncé 1 (True/False) :} \eng{Teenagers in Yaoundé get their news mainly from television. Justify.}

\tcblower

\textbf{Solution :} \eng{False. “Most teenagers now get their news from Facebook and WhatsApp rather than from television or newspapers.”} — La justification est recopiée du texte entre guillemets.
\end{exoresolu}

\begin{exoresolu}[QCM corrigé]
\textbf{Énoncé 2 (QCM) :} \eng{According to the survey, students aged 15 to 19 spend...\\
A. less than one hour online. \quad B. exactly three hours online.\\
C. more than three hours online. \quad D. no time online.}

\tcblower

\textbf{Solution :} \eng{C. more than three hours online.} — Le texte dit : \eng{“70\% of students aged 15 to 19 spend more than three hours a day online.”} Les options A et D sont absurdes ; B est un piège (\eng{exactly} contredit \eng{more than}).
\end{exoresolu}

\begin{exoresolu}[Réponse courte corrigée]
\textbf{Énoncé 3 (short answer) :} \eng{Why are teachers and parents worried?}

\tcblower

\textbf{Solution :} \eng{Teachers and parents are worried because social media spread fake news very quickly.} — La réponse reprend la formulation de la question et commence par \eng{because}, puisque la question commence par \eng{why}.
\end{exoresolu}

\begin{exoresolu}[Vocabulaire en contexte corrigé]
\textbf{Énoncé 4 (vocabulary) :} \eng{Find in the text words or expressions meaning: (a) to verify, (b) false, (c) young people.}

\tcblower

\textbf{Solution :} (a) \eng{to verify} = \eng{to check} (dans \eng{“to check information”}) ; (b) \eng{false} = \eng{fake} (dans \eng{fake news}) ; (c) \eng{young people} = \eng{teenagers} / \eng{the youth}. Chaque réponse est un mot réellement présent dans le texte.
\end{exoresolu}

\courssub{Common traps — les pièges des francophones}

\begin{attention}[Faux amis et traduction mot à mot]
\begin{itemize}
\item \eng{actually} signifie « en fait », \textbf{jamais} « actuellement ». Pour « actuellement », dites \eng{currently} ou \eng{at present}.
\item \eng{a survey} = « une enquête, un sondage », pas « une surveillance ».
\item \eng{to assist} = « aider » ; « assister à » (un match, une cérémonie) se dit \eng{to attend}.
\item \eng{to check} = « vérifier » ; \eng{to control} signifie « contrôler, commander » (une machine, une situation).
\item Ne traduisez \textbf{jamais} le texte mot à mot pour répondre : lisez pour comprendre les idées, pas pour traduire chaque phrase. Le correcteur évalue votre compréhension en anglais, pas votre traduction en français.
\item Attention aux négations et aux mots restrictifs dans les QCM : \eng{rather than}, \eng{without}, \eng{hardly}, \eng{instead of} inversent ou limitent le sens de la phrase.
\end{itemize}
\end{attention}

\begin{center}
\begin{tabularx}{\linewidth}{|X|X|X|}
\hline
\rowcolor{popblueL} \textbf{Français} & \textbf{English correct} & \textbf{Piège à éviter} \\
\hline
en fait & actually & pas « actually » pour « actuellement » \\
\hline
actuellement & currently, at present & pas « actually » \\
\hline
une enquête / un sondage & a survey & pas « a surveillance » \\
\hline
vérifier une information & to check information & pas « to control » (= commander) \\
\hline
de fausses informations & fake news, disinformation & pas « false news » \\
\hline
sensibiliser les jeunes & to raise young people's awareness & pas « to sensibilize » \\
\hline
\end{tabularx}
\end{center}

\courssub{Your turn — entraînement de type examen}

\begin{exercice}[Reading comprehension — exam practice (10 marks)]
Read the text \emph{Social Media and the Cameroonian Youth} again and answer the following questions.

\textbf{A. True or False? Justify with a line from the text. (4 marks)}
\begin{enumerate}
\item \eng{All students in Cameroon spend three hours a day online.}
\item \eng{Some schools have created press clubs.}
\end{enumerate}

\textbf{B. Choose the correct answer. (2 marks)}
\begin{enumerate}
\item \eng{Media-literacy lessons teach students how to...}\\\eng{A. use Facebook. \quad B. check information before sharing it. \quad C. write fake news. \quad D. watch television.}
\item \eng{According to experts, the best protection against disinformation is...}\\\eng{A. WhatsApp. \quad B. television. \quad C. critical thinking. \quad D. the Ministry.}
\end{enumerate}

\textbf{C. Answer in complete sentences. (2 marks)}
\begin{enumerate}
\item \eng{What three questions do students learn to ask in media-literacy lessons?}
\item \eng{How does the Ministry of Secondary Education react to the problem?}
\end{enumerate}

\textbf{D. Vocabulary in context. (2 marks)}
\begin{enumerate}
\item \eng{Find in the text a synonym of “false”.}
\item \eng{Find in the text the opposite of “unreliable”.}
\end{enumerate}
\end{exercice}

\begin{corrige}[Exercice — Reading comprehension]
\textbf{A.} 1. \eng{False. “70\% of students aged 15 to 19 spend more than three hours a day online.”} (70 \%, pas tous ; et \emph{more than}, pas exactement trois.) 2. \eng{True. “Some schools have even created ‘press clubs’ where teenagers write short articles for a school magazine.”}

\textbf{B.} 1. \eng{B. check information before sharing it.} 2. \eng{C. critical thinking.}

\textbf{C.} 1. \eng{Students learn to ask: “Who wrote this? When was it published? Can I find the same story on a reliable website?”} 2. \eng{The Ministry of Secondary Education encourages every school to include media education in its programme.}

\textbf{D.} 1. \eng{fake} (dans \eng{fake news}). 2. \eng{reliable} (dans \eng{“a reliable website”}).
\end{corrige}

\begin{aretenir}[L'essentiel de la section]
\begin{itemize}
\item Lisez les \textbf{questions avant le texte}, puis lisez le texte \textbf{deux fois}.
\item Chaque réponse doit être \textbf{justifiée par une ligne citée du texte} — jamais de mémoire.
\item Reprenez la \textbf{formulation de la question} dans votre réponse.
\item Méfiez-vous des \textbf{faux amis} (\eng{actually}, \eng{survey}, \eng{to control}) et des \textbf{négations} dans les QCM.
\item La compréhension écrite est la section la plus rentable du Bac : la méthode compte autant que le niveau de langue.
\end{itemize}
\end{aretenir}

\coursec{Grammar synthesis — verb tenses}

Après l'entraînement à la compréhension de l'écrit de la section précédente, nous revenons maintenant sur le cœur de la grammaire anglaise : les \cle{temps verbaux}. Au Baccalauréat, les épreuves de grammaire et de composition exigent que vous choisissiez spontanément le bon temps, sans hésitation. Cette synthèse reprend, temps par temps, tout ce que nous avons étudié cette année, avec les repères qui permettent de ne jamais se tromper : la forme, l'emploi, les marqueurs temporels et les pièges classiques du francophone.

\courssub{Observing the tenses in context}

Lisons d'abord ce court texte original, qui mélange volontairement plusieurs temps. Repérez les verbes conjugués et les marqueurs temporels avant de lire l'observation qui suit.

\begin{textebox}[A journalist's week — original text]
Every morning, Amina checks her e-mails before she goes to the newsroom in Yaoundé. She has worked for the same newspaper since 2019, and she has already published more than two hundred articles. Last Tuesday, while she was interviewing a minister, her phone rang three times. She has just finished a long report about social media in Cameroon, and her editor says it will be published next week. Amina is going to travel to Douala next month because the newspaper is opening a new office there. “I love this job,” she said yesterday. She added that she had never imagined such a career when she was a student. Right now, she is preparing her questions for tomorrow's press conference.
\end{textebox}

\textbf{Glossaire.} \emph{newsroom} (n.) : salle de rédaction ; \emph{editor} (n.) : rédacteur en chef ; \emph{report} (n.) : reportage ; \emph{press conference} (n.) : conférence de presse.

\textbf{Observation.} \eng{checks} (habitude, present simple), \eng{has worked / has published / has just finished} (lien passé-présent, present perfect), \eng{was interviewing / rang} (action longue interrompue par une action courte), \eng{will be published / is going to travel / is opening} (trois futurs différents), \eng{she had never imagined} (recul d'un temps au discours rapporté). Toute la grammaire de l'année est concentrée dans ce paragraphe.

\courssub{The present tenses}

\begin{propriete}[Present simple et present continuous]
\textbf{Present simple} : \emph{base verbale} (+ \emph{s} à la 3\up{e} personne du singulier). Emplois : habitudes, vérités générales, faits permanents.
\par \eng{Journalists check their sources.} « Les journalistes vérifient leurs sources. »
\par \eng{Water boils at 100 degrees.} « L'eau bout à 100 degrés. » (vérité générale)
\par \textbf{Present continuous} : \emph{be + V-ing}. Emplois : action en cours au moment où l'on parle, ou futur déjà programmé (rendez-vous fixé avec date et lieu).
\par \eng{She is preparing her questions right now.} « Elle prépare ses questions en ce moment. »
\par \eng{We are meeting the editor at 4 p.m. tomorrow.} « Nous rencontrons le rédacteur en chef demain à 16 h. » (futur programmé)
\end{propriete}

\begin{attention}[Verbes sans forme continue]
Les verbes d'état (\emph{know, believe, like, love, want, need, understand, own, seem}) ne se mettent pas à la forme continue. On dit \eng{I know the answer}, jamais « I am knowing » ; \eng{She wants a new phone}, jamais « She is wanting ». Attention aussi à l'orthographe du \emph{-ing} : \emph{write} $\to$ \emph{writing} (on supprime le e final), \emph{run} $\to$ \emph{running} (on double la consonne après une voyelle courte accentuée), \emph{lie} $\to$ \emph{lying}.
\end{attention}

\courssub{Past simple vs present perfect}

C'est la distinction la plus délicate pour un francophone, car le français utilise souvent le passé composé pour les deux situations.

\begin{propriete}[Le present perfect relie le passé au présent]
\textbf{Past simple} (\emph{V-ed} ou 2\up{e} colonne des verbes irréguliers) : action \textbf{terminée}, située dans un passé \textbf{daté et coupé du présent}. Marqueurs : \emph{yesterday, ago, last week, in 2015, when}.
\par \textbf{Present perfect} (\emph{have/has + participe passé}) : action passée \textbf{en lien avec le présent} (résultat visible, expérience de vie, action qui continue encore). Marqueurs : \emph{already, just, never, ever, since, for, yet, so far}.
\par \eng{I have seen him yesterday} est \textbf{FAUX} : \emph{yesterday} date l'action, il faut le past simple : \eng{I saw him yesterday.} « Je l'ai vu hier. »
\end{propriete}

\begin{exemplebox}[Exemples traduits]
\eng{She has just published her article.} « Elle vient de publier son article. » (résultat présent : l'article est disponible)
\par \eng{I have never read this newspaper.} « Je n'ai jamais lu ce journal. » (expérience de vie)
\par \eng{He bought a smartphone two years ago.} « Il a acheté un smartphone il y a deux ans. » (passé daté)
\par \eng{Have you ever listened to the BBC?} « As-tu déjà écouté la BBC ? » (question sur l'expérience)
\end{exemplebox}

\begin{attention}[Since ou for ?]
\cle{since} + \textbf{point de départ} : \eng{since 2020, since Monday, since I was a child} « depuis 2020, depuis lundi ».
\par \cle{for} + \textbf{durée} : \eng{for three years, for a long time} « depuis trois ans, pendant longtemps ».
\par Erreur classique du francophone : « since three years » est faux ; dites \eng{for three years} ou \eng{since 2022}. Autre piège majeur : le français dit « je vis ici depuis dix ans » au présent ; l'anglais exige le present perfect : \eng{I have lived here for ten years.} — jamais « I live here since ten years ».
\end{attention}

\courssub{Past continuous : the interrupted action}

\begin{propriete}[While + action longue, when + interruption]
\textbf{Past continuous} (\emph{was/were + V-ing}) : action longue en cours dans le passé, souvent interrompue par une action courte au past simple.
\par Structure : \emph{while + past continuous} (l'action longue) ; \emph{when + past simple} (l'interruption).
\par \eng{While I was listening to the radio, the electricity went off.} « Pendant que j'écoutais la radio, l'électricité a été coupée. »
\par \eng{She was driving to Buea when her editor called.} « Elle conduisait vers Buea quand son rédacteur en chef a appelé. »
\par Deux actions longues simultanées : \eng{While she was typing, he was taking photos.} « Pendant qu'elle tapait, il prenait des photos. »
\end{propriete}

\courssub{Past perfect : the earlier past}

\begin{propriete}[Had + participe passé : l'antériorité]
Le \textbf{past perfect} (\emph{had + participe passé}) situe une action \textbf{avant} une autre action passée. C'est l'équivalent du plus-que-parfait français.
\par \eng{When the guests arrived, the programme had already started.} « Quand les invités sont arrivés, l'émission avait déjà commencé. »
\par \eng{She had never imagined such a career.} « Elle n'avait jamais imaginé une telle carrière. »
\par Marqueurs fréquents : \emph{before, after, by the time, already, never}.
\end{propriete}

\courssub{Expressing the future : three choices}

\begin{propriete}[Will, going to, present continuous]
\begin{itemize}
\item \cle{will + base verbale} : décision \textbf{spontanée} (prise au moment de parler), \textbf{prédiction} personnelle ou promesse. \eng{— The phone is ringing! — I will answer it.} « Le téléphone sonne ! Je vais répondre. » \eng{I think online media will replace paper newspapers.} « Je pense que les médias en ligne remplaceront les journaux papier. »
\item \cle{be going to + base verbale} : \textbf{projet} déjà décidé, ou conséquence \textbf{visible} (évidence sous les yeux). \eng{Amina is going to travel to Douala next month.} « Amina va se rendre à Douala le mois prochain. » \eng{Look at those clouds! It is going to rain.} « Regarde ces nuages ! Il va pleuvoir. »
\item \cle{present continuous} : rendez-vous \textbf{fixé}, avec date et lieu précis. \eng{We are meeting the minister at 10 a.m. on Friday.} « Nous rencontrons le ministre vendredi à 10 h. »
\end{itemize}
\end{propriete}

\courssub{Reported speech : the tense shift}

\begin{propriete}[Au discours rapporté, recul d'un temps]
Quand le verbe introducteur est au passé (\emph{said, told, asked}), le temps de la parole rapportée \textbf{recule d'un temps} :
\par present simple $\to$ past simple ; present continuous $\to$ past continuous ; past simple et present perfect $\to$ past perfect ; will $\to$ would ; can $\to$ could.
\par \eng{He said: “I read the newspaper.”} $\to$ \eng{He said (that) he read the newspaper.} « Il a dit qu'il lisait le journal. »
\par \eng{She said: “I have finished my report.”} $\to$ \eng{She said she had finished her report.} « Elle a dit qu'elle avait fini son reportage. »
\par \eng{They said: “We will broadcast the debate.”} $\to$ \eng{They said they would broadcast the debate.} « Ils ont dit qu'ils diffuseraient le débat. »
\par N'oubliez pas les changements de repères : \emph{now} $\to$ \emph{then}, \emph{today} $\to$ \emph{that day}, \emph{yesterday} $\to$ \emph{the day before}, \emph{tomorrow} $\to$ \emph{the next day}, \emph{here} $\to$ \emph{there}.
\end{propriete}

\courssub{Time markers : the summary table}

\begin{center}
\begin{tabularx}{\linewidth}{|l|X|X|}
\hline
\rowcolor{popblueL} \textbf{Temps} & \textbf{Marqueurs typiques} & \textbf{Exemple} \\
\hline
Present simple & always, usually, every day, often, never & \eng{She checks her e-mails every morning.} \\
\hline
Present continuous & now, right now, at the moment, Look! & \eng{She is preparing her questions now.} \\
\hline
Past simple & yesterday, ago, last week, in 2015, when & \eng{Her phone rang three times.} \\
\hline
Past continuous & while, when, at 8 p.m. yesterday & \eng{She was interviewing a minister.} \\
\hline
Present perfect & already, just, never, ever, since, for, yet & \eng{She has just finished her report.} \\
\hline
Past perfect & before, after, by the time, already & \eng{She had never imagined such a career.} \\
\hline
Future (will / going to) & tomorrow, next week, soon, in 2030 & \eng{It will be published next week.} \\
\hline
\end{tabularx}
\end{center}

\begin{popfigure}
\begin{tikzpicture}[scale=1]
\draw[-{Stealth[length=3mm]}, thick, popdark] (-5.2,0) -- (5.6,0);
\foreach \x/\lab in {-3.6/{past perfect}, -1.6/{past simple}, 0/{NOW}, 2.2/{present perfect}, 4.4/{future}} {
  \draw[popblue, thick] (\x,-0.12) -- (\x,0.12);
  \node[below, font=\small, text=popdark] at (\x,-0.15) {\lab};
}
\node[circle, fill=poporange, inner sep=2.5pt] at (0,0) {};
\node[above, font=\small\itshape, text=poppurple, align=center] at (-3.6,0.25) {She had finished\\before noon.};
\node[above, font=\small\itshape, text=popgreen, align=center] at (-1.6,0.25) {She published it\\yesterday.};
\node[above, font=\small\itshape, text=popblue, align=center] at (2.2,0.25) {She has just\\sent the article.};
\node[above, font=\small\itshape, text=poporange, align=center] at (4.4,0.25) {It will appear\\next week.};
\end{tikzpicture}
\legende{Frise des temps : le past perfect précède le past simple ; le present perfect touche le présent ; le futur s'ouvre après NOW.}
\end{popfigure}

\begin{popfigure}
\begin{tikzpicture}[mindmap, grow cyclic, every node/.style={concept, font=\small},
  root concept/.append style={concept color=popblue, font=\small\bfseries},
  level 1/.append style={level distance=3.1cm, sibling angle=90},
  level 2/.append style={level distance=2.3cm, sibling angle=50, font=\scriptsize}]
\node[root concept]{Tenses}
  child[concept color=popgreen] { node {Present}
    child { node[concept color=popgreenL] {simple : habits, every day} }
    child { node[concept color=popgreenL] {continuous : now, Look!} } }
  child[concept color=poporange] { node {Past}
    child { node[concept color=poporangeL] {simple : yesterday, ago} }
    child { node[concept color=poporangeL] {continuous : while, when} } }
  child[concept color=poppurple] { node {Perfect}
    child { node[concept color=poppurpleL] {present : just, since, for} }
    child { node[concept color=poppurpleL] {past : before, already} } }
  child[concept color=popgold] { node {Future}
    child { node[concept color=popgoldL] {will : prediction} }
    child { node[concept color=popgoldL] {going to : plan} } };
\end{tikzpicture}
\legende{Carte mentale des temps anglais et de leurs marqueurs.}
\end{popfigure}

\begin{methode}[Choisir le bon temps en trois étapes]
\begin{enumerate}
\item \textbf{Repérer le marqueur temporel} dans la phrase (\emph{yesterday, since, now, tomorrow}\ldots) : il oriente presque toujours vers un temps précis.
\item \textbf{Se demander si l'action est finie ou en lien avec le présent} : finie et datée $\to$ past simple ; résultat présent ou expérience $\to$ present perfect ; en cours $\to$ forme continue ; antérieure à une autre action passée $\to$ past perfect.
\item \textbf{Vérifier la concordance} : si la phrase rapporte une parole au passé, faire reculer le temps d'un cran (\emph{will} $\to$ \emph{would}, \emph{have done} $\to$ \emph{had done}).
\end{enumerate}
\end{methode}

\begin{exoresolu}[Choose the correct tense]
Énoncé — Complétez avec le temps correct du verbe entre parenthèses :
\par 1. I \ldots\ (never / visit) the CRTV studios, but I would love to.
\par 2. While we \ldots\ (watch) the debate, the connection \ldots\ (cut) off.
\par 3. She said she \ldots\ (send) the article the day before.
\par 4. Look! It \ldots\ (rain); take your umbrella.
\tcblower
Solution détaillée :
\par 1. \eng{have never visited} — marqueur \emph{never}, expérience de vie en lien avec le présent : present perfect.
\par 2. \eng{were watching ; was cut} — action longue (\emph{while} + past continuous) interrompue par une action courte. La connexion « subit » la coupure : past simple \textbf{passif} (\emph{was cut off}). \emph{cut} est irrégulier et invariable : cut--cut--cut.
\par 3. \eng{had sent} — discours rapporté avec \emph{said} au passé : le temps recule au past perfect ; \emph{the day before} confirme le recul des repères.
\par 4. \eng{is going to rain} — prédiction fondée sur une évidence visible (\emph{Look!}) : \emph{be going to}, pas \emph{will}.
\end{exoresolu}

\begin{exercice}[Correct the mistakes]
Chaque phrase contient une erreur de temps. Corrigez-la et justifiez.
\par 1. \eng{I have seen that documentary last night.}
\par 2. \eng{She works for this radio station since 2018.}
\par 3. \eng{When I arrived, they already broadcast the news.}
\par 4. \eng{He said he will call me tomorrow.}
\par 5. \eng{Look at the sky! It will rain.}
\end{exercice}

\begin{corrige}[Exercice « Correct the mistakes »]
1. \eng{I saw that documentary last night.} — \emph{last night} est un repère passé daté : past simple obligatoire.
\par 2. \eng{She has worked for this radio station since 2018.} — action commencée dans le passé qui continue encore : present perfect (le présent français trompe ici).
\par 3. \eng{When I arrived, they had already broadcast the news.} — action antérieure à une action passée : past perfect (\emph{broadcast} est invariable : broadcast--broadcast--broadcast).
\par 4. \eng{He said he would call me the next day.} — discours rapporté : \emph{will} $\to$ \emph{would}, \emph{tomorrow} $\to$ \emph{the next day}.
\par 5. \eng{Look at the sky! It is going to rain.} — prédiction fondée sur une évidence visible : \emph{be going to}.
\end{corrige}

\begin{savaistu}[Used to ou be used to ?]
\cle{used to + infinitif} exprime une habitude passée \textbf{révolue} : \eng{I used to read paper newspapers.} « Je lisais les journaux papier (avant, mais plus maintenant). »
\par \cle{to be used to + -ing} signifie « être habitué à » : \eng{I am used to reading the news on my phone.} « J'ai l'habitude de lire les informations sur mon téléphone. »
\par Deux structures, deux sens différents : l'une regarde le passé, l'autre décrit le présent. Les examinateurs adorent ce piège !
\end{savaistu}

\begin{aretenir}[L'essentiel de la synthèse]
Un marqueur temporel daté (\emph{yesterday, ago, last}) impose le past simple et interdit le present perfect. \emph{Since} + point de départ, \emph{for} + durée. \emph{While} + action longue, \emph{when} + interruption. Le past perfect exprime l'antériorité (\emph{had + participe passé}). Trois futurs pour trois situations : \emph{will} (spontané, prédiction), \emph{going to} (projet, évidence), present continuous (rendez-vous fixé). Au discours rapporté, tout recule d'un temps. Entraînez-vous avec ces réflexes, et la grammaire du Bac n'aura plus de secret pour vous.
\end{aretenir}

\coursec{Grammar synthesis — key structures}

Dans la section précédente, nous avons fait la synthèse des \cle{temps verbaux} anglais. Mais maîtriser les temps ne suffit pas à l'examen : le Bac teste aussi votre capacité à manipuler les \cle{structures} de la phrase — passer de l'actif au passif, rapporter les paroles de quelqu'un, comparer, exprimer une hypothèse ou un regret. C'est précisément l'objet de cette section : réviser, une par une, les grandes structures étudiées cette année, avec pour chacune la règle, des exemples traduits, les pièges classiques des francophones, et des exercices de transformation de type Bac.

\courssub{The passive voice — la voix passive}

\textbf{Observation.} Lisez ces deux phrases tirées d'un article sur les médias :

\eng{The press publishes the news every morning.} « La presse publie les informations chaque matin. »

\eng{The news is published every morning.} « Les informations sont publiées chaque matin. »

Les deux phrases décrivent le même fait, mais le point de vue change : dans la seconde, c'est \eng{the news} qui devient le sujet, et l'agent (\eng{the press}) disparaît ou passe en second plan. C'est la \cle{voix passive}.

\begin{propriete}[Formation du passif]
\textbf{Forme :} sujet + \cle{be} (conjugué au temps voulu) + \cle{participe passé} (+ \eng{by} + agent, facultatif).

\eng{Fake news is spread quickly.} « Les fausses informations se répandent vite. »

Le temps se porte sur \cle{be}, jamais sur le participe passé :
\begin{center}
\begin{tabular}{|l|l|}
\hline
\rowcolor{popblueL} Temps & Exemple au passif \\
\hline
Present simple & News \textbf{is broadcast} at 8 p.m. \\
\hline
Past simple & The journalist \textbf{was interviewed} yesterday. \\
\hline
Present perfect & The report \textbf{has been published}. \\
\hline
Future & The results \textbf{will be announced} tomorrow. \\
\hline
Modal + passif & The letter \textbf{must be sent} today. \\
\hline
\end{tabular}
\end{center}
\end{propriete}

\textbf{Emploi.} On utilise le passif quand l'\cle{agent} est inconnu, évident ou secondaire : \eng{My phone was stolen.} « Mon téléphone a été volé » (on ne sait pas par qui). Le complément d'agent est introduit par \eng{by} : \eng{The journalist was interviewed by the BBC.} « Le journaliste a été interviewé par la BBC. »

\begin{popfigure}
\begin{center}
\begin{tikzpicture}[node distance=0.6cm and 1.6cm]
\node[draw=popblue, fill=popblueL, rounded corners, align=center, text width=4.6cm] (active) {\textbf{Active}\\ The press\\ \textbf{publishes}\\ the news};
\node[draw=popgreen, fill=popgreenL, rounded corners, align=center, text width=4.6cm, right=of active] (passive) {\textbf{Passive}\\ The news\\ \textbf{is published}\\ (by the press)};
\draw[-{Stealth[length=3mm]}, very thick, poporange] (active) -- (passive) node[midway, above, align=center, text=popdark]{\footnotesize objet $\rightarrow$ sujet\\ \footnotesize be + p.p.};
\end{tikzpicture}
\end{center}
\legende{Schéma de la transformation passive : le complément d'objet devient sujet, le verbe prend la forme \eng{be} + participe passé.}
\end{popfigure}

\begin{attention}[Piège des francophones]
Le français utilise souvent « on » ou un pronominal là où l'anglais exige le passif : « Ici on parle anglais » se dit \eng{English is spoken here} (jamais \eng{Here is spoken English}). Et attention : seuls les verbes \cle{transitifs} (suivis d'un complément d'objet) peuvent se mettre au passif. On ne peut pas dire \eng{I was arrived}.
\end{attention}

\courssub{Reported speech — questions et ordres}

Nous avons vu le discours rapporté des affirmations ; révisons maintenant les \cle{questions} et les \cle{ordres} rapportés, très fréquents au Bac.

\begin{propriete}[Questions rapportées]
\textbf{Forme :} \eng{asked} + (if/whether ou mot interrogatif) + \textbf{sujet + verbe} (ordre de la phrase affirmative, sans inversion, sans auxiliaire \eng{do/did}).

\eng{“Where do you live?” $\rightarrow$ She asked me where I lived.} « Elle m'a demandé où j'habitais. »

\eng{“Do you watch the news?” $\rightarrow$ He asked if I watched the news.} « Il a demandé si je regardais les informations. »

Pour les questions fermées (oui/non), on introduit \eng{if} ou \eng{whether}. Le recul des temps s'applique comme pour les affirmations (present $\rightarrow$ past, etc.). D'autres verbes introducteurs sont possibles : \eng{wonder}, \eng{want to know}.
\end{propriete}

\begin{attention}[L'erreur n°1 au Bac]
Dans une question rapportée, il n'y a \textbf{plus d'inversion} : on écrit \eng{She asked me where I lived} et jamais \eng{She asked me where did I live}. De même, pas de point d'interrogation à la fin d'une question rapportée.
\end{attention}

\begin{propriete}[Ordres et conseils rapportés]
\textbf{Forme :} \eng{tell / order / advise / warn} + complément + \cle{to} + infinitif ; à la forme négative : \eng{not to} + infinitif.

\eng{“Turn off your phone,” the teacher said. $\rightarrow$ The teacher told me to turn off my phone.} « Le professeur m'a dit d'éteindre mon téléphone. »

\eng{“Don't be late.” $\rightarrow$ She told me not to be late.} « Elle m'a dit de ne pas être en retard. »
\end{propriete}

\courssub{Comparatives and superlatives — comparatifs et superlatifs}

\begin{propriete}[Formation]
\begin{itemize}
\item Adjectifs \textbf{courts} (1 syllabe) : \cle{-er} / \cle{-est} : \eng{fast $\rightarrow$ faster $\rightarrow$ the fastest}.
\item Adjectifs \textbf{longs} (2 syllabes et +) : \cle{more} / \cle{most} : \eng{important $\rightarrow$ more important $\rightarrow$ the most important}.
\item Adjectifs en \textbf{-y} : \eng{easy $\rightarrow$ easier $\rightarrow$ the easiest}.
\item \textbf{Irréguliers} : \eng{good $\rightarrow$ better $\rightarrow$ the best} ; \eng{bad $\rightarrow$ worse $\rightarrow$ the worst} ; \eng{far $\rightarrow$ further $\rightarrow$ the furthest}.
\end{itemize}
Le comparatif de supériorité est suivi de \cle{than} : \eng{Radio is cheaper than television.} « La radio est moins chère que la télévision. » Le superlatif prend \cle{the} : \eng{The internet is the most powerful medium.} « Internet est le média le plus puissant. »
Pour l'infériorité : \cle{less} + adjectif + \cle{than} (\eng{less influential than}). Pour le superlatif d'infériorité : \cle{the least} + adjectif (\eng{the least reliable source}).
\end{propriete}

\begin{attention}[Double marque interdite]
Ne jamais cumuler les deux marques : on dit \eng{better} ou \eng{more interesting}, jamais \eng{more better} ni \eng{the most easiest}. Et « que » de comparaison = \eng{than}, pas \eng{that}.
\end{attention}

\courssub{Conditionals — les conditionnels (révision synthétique)}

\begin{propriete}[Second conditional — hypothèse irréelle au présent/futur]
\textbf{Forme :} \cle{if} + past simple, \cle{would} + infinitif. Exprime une situation imaginaire ou peu probable maintenant ou dans le futur.

\eng{If I had a smartphone, I would follow the news online.} « Si j'avais un smartphone, je suivrais l'actualité en ligne. »

Remarque : avec \eng{be}, la forme correcte est \eng{were} pour toutes les personnes dans le style soigné : \eng{If I were you, I would read the article.} « Si j'étais toi, je lirais l'article. »
\end{propriete}

\begin{propriete}[Third conditional — regret / hypothèse passée irréelle]
\textbf{Forme :} \cle{if} + past perfect, \cle{would have} + participe passé. Exprime une situation passée qui n'a pas eu lieu.

\eng{If I had watched the debate, I would have understood the issue.} « Si j'avais regardé le débat, j'aurais compris la question. »
\end{propriete}

\begin{propriete}[Mixed conditionals — conditionnels mixtes]
Combinaison fréquente : passé $\rightarrow$ présent. \eng{If I had studied journalism (past), I would be a reporter now (present).} « Si j'avais étudié le journalisme, je serais reporter maintenant. »
\end{propriete}

\begin{attention}[Pas de would dans la subordonnée]
Le \eng{would} (\eng{would have}) se met dans la proposition principale, jamais après \eng{if} : on dit \eng{If I had money, I would buy a radio}, jamais \eng{If I would have money}.
\end{attention}

\courssub{Wish — exprimer le regret ou le souhait}

\begin{propriete}[Wish + past simple / past perfect / would]
\begin{itemize}
\item \cle{wish} + \textbf{past simple} : regret sur le \textbf{présent} / futur. \eng{I wish I spoke English fluently.} « J'aimerais parler anglais couramment » (mais ce n'est pas le cas).
\item \cle{wish} + \textbf{past perfect} : regret sur le \textbf{passé}. \eng{I wish I had watched the debate last night.} « J'aurais aimé regarder le débat hier soir » (mais je ne l'ai pas fait).
\item \cle{wish} + \textbf{would} + infinitif : agacement ou désir de changement chez autrui. \eng{I wish you would stop shouting.} « J'aimerais que tu arrêtes de crier. »
\end{itemize}
\end{propriete}

\begin{attention}[Wish vs Hope]
\eng{Wish} + past = irréel (regret). \eng{Hope} + present = réel possible : \eng{I hope you pass your exam.} (pas \eng{I wish you pass}).
\end{attention}

\courssub{Gerund or infinitive — gérondif ou infinitif}

\begin{propriete}[Verbes suivis du gérondif ou de l'infinitif (sens inchangé)]
\begin{center}
\begin{tabularx}{\linewidth}{|X|X|}
\hline
\rowcolor{popblueL} Verbe + \textbf{-ing} (gérondif) & Verbe + \textbf{to} + infinitif \\
\hline
\eng{enjoy watching TV} & \eng{want to learn English} \\
\hline
\eng{suggest watching the debate} & \eng{decide to buy a newspaper} \\
\hline
\eng{avoid reading fake news} & \eng{hope to become a journalist} \\
\hline
\eng{finish writing the article} & \eng{promise to call back} \\
\hline
\eng{mind answering questions} & \eng{agree to give an interview} \\
\hline
\end{tabularx}
\end{center}
\end{propriete}

\begin{propriete}[Verbes à double construction — changement de sens]
\begin{center}
\begin{tabularx}{\linewidth}{|l|X|X|}
\hline
\rowcolor{popblueL} Verbe & \textbf{-ing} & \textbf{to + infinitif} \\
\hline
\eng{stop} & \eng{stop reading} = arrêter de lire & \eng{stop to read} = s'arrêter pour lire \\
\hline
\eng{remember} & \eng{remember posting} = se souvenir avoir posté & \eng{remember to post} = ne pas oublier de poster \\
\hline
\eng{forget} & \eng{forget meeting him} = oublier l'avoir rencontré & \eng{forget to meet him} = oublier de le rencontrer \\
\hline
\eng{try} & \eng{try clicking} = essayer (tester) de cliquer & \eng{try to click} = tenter de cliquer (effort) \\
\hline
\eng{regret} & \eng{regret telling him} = regretter de lui avoir dit & \eng{regret to inform you} = regretter de devoir vous informer (formel) \\
\hline
\end{tabularx}
\end{center}
\end{propriete}

\begin{attention}[Suggest : le piège classique]
Après \eng{suggest}, jamais de \eng{to} ni de \eng{that} + indicatif : on dit \eng{He suggested watching the debate.} « Il a suggéré de regarder le débat », et non \eng{He suggested to watch}. La structure \eng{suggest that} + subjonctif (\eng{He suggested that we watch}) est possible mais le gérondif est privilégié aux transformations du Bac.
\end{attention}

\courssub{Sentence transformation — méthode et entraînement type Bac}

\begin{methode}[Réussir les transformations de phrases]
\begin{enumerate}
\item \textbf{Identifier} la structure demandée : passif ? discours rapporté ? conditionnel ? comparatif ?
\item \textbf{Repérer} ce qui change : le sujet, la forme du verbe, l'ordre des mots, la négation.
\item \textbf{Ne changer que} ce qui doit changer : garder le même temps de référence et le même sens.
\item \textbf{Relire} : vérifier l'accord (\eng{is/are}, \eng{was/were}), l'ordre sujet + verbe dans les questions rapportées, et l'orthographe des irréguliers.
\end{enumerate}
\end{methode}

\begin{exoresolu}[Transformations type Bac]
\textbf{Énoncé.} Rewrite each sentence as indicated.

1. \eng{The BBC interviewed the journalist.} (passive voice)

2. \eng{“Where do you get your news?” she asked me.} (reported speech)

3. \eng{“Don't share that video,” he told me.} (reported speech)

4. \eng{I don't have a smartphone, so I can't follow the news online.} (use \eng{if} + conditional)
\tcblower
\textbf{Solution détaillée.}

1. Le complément d'objet \eng{the journalist} devient sujet ; \eng{interviewed} (past simple) devient \eng{was} + participe passé : \eng{The journalist was interviewed by the BBC.}

2. Question rapportée : mot interrogatif conservé, ordre sujet + verbe, recul du temps : \eng{She asked me where I got my news.} (et non \eng{where did I get}).

3. Ordre négatif rapporté : \eng{He told me not to share that video.}

4. Hypothèse irréelle au présent (Second conditional) : \eng{If I had a smartphone, I would follow the news online.}
\end{exoresolu}

\begin{exercice}[À vous de jouer]
Transformez les phrases suivantes :

1. \eng{People spread fake news very quickly.} (passive voice)

2. \eng{“Do you listen to the radio?” the reporter asked the students.} (reported speech)

3. \eng{“Read the article carefully,” the teacher said to us.} (reported speech)

4. \eng{I regret that I didn't watch the presidential debate.} (use \eng{wish})

5. \eng{Television is influential. The internet is more influential.} (comparative, one sentence)

6. \eng{If the journalist checks the facts, the article will be reliable.} (Third conditional — start with \eng{If the journalist had checked...})

7. \eng{He forgot to lock the door.} (change to gerund with change of meaning)
\end{exercice}

\begin{corrige}[Exercice — transformations]
1. \eng{Fake news is spread very quickly.} (agent \eng{by people} omis car générique) — 2. \eng{The reporter asked the students if they listened to the radio.} — 3. \eng{The teacher told us to read the article carefully.} — 4. \eng{I wish I had watched the presidential debate.} — 5. \eng{The internet is more influential than television.} — 6. \eng{If the journalist had checked the facts, the article would have been reliable.} — 7. \eng{He forgot locking the door.} (il a oublié qu'il l'avait verrouillée / il a oublié l'action passée).
\end{corrige}

\begin{aretenir}[L'essentiel des structures]
\begin{itemize}
\item Passif : \eng{be} + participe passé ; l'agent avec \eng{by} ; réservé aux verbes transitifs.
\item Question rapportée : \eng{asked if/whether} ou mot interrogatif + sujet + verbe, sans inversion, pas de \eng{?}.
\item Ordre rapporté : \eng{told / advised / warned} + me + \eng{(not) to} + infinitif.
\item Comparatif : \eng{-er than} / \eng{more ... than} / \eng{less ... than} ; Superlatif : \eng{the -est} / \eng{the most} / \eng{the least} ; Irréguliers : \eng{good/better/best}, \eng{bad/worse/worst}, \eng{far/further/furthest}.
\item Conditionnels : 2\textsuperscript{e} (\eng{if} + past, \eng{would} + V) ; 3\textsuperscript{e} (\eng{if} + past perfect, \eng{would have} + V-ed) ; Mixtes.
\item Regret / Souhait : \eng{wish} + past simple (présent), \eng{wish} + past perfect (passé), \eng{wish} + \eng{would} (agacement).
\item Gérondif après \eng{enjoy, suggest, avoid, finish, mind, stop, remember, forget, try, regret} (selon sens) ; Infinitif après \eng{want, decide, hope, promise, agree, stop, remember, forget, try, regret} (selon sens).
\end{itemize}
\end{aretenir}

\coursec{Vocabulary consolidation — Media and Communication}

Après avoir consolidé les structures grammaticales essentielles dans la section précédente, il est temps de faire le point sur le vocabulaire. Au baccalauréat, une copie riche et précise se reconnaît d'abord à la qualité de son lexique : un candidat qui écrit \eng{The journalist published an article in a daily newspaper} impressionnera toujours plus le correcteur que celui qui écrit \eng{The man wrote something in a paper}. Cette section révise et approfondit le champ lexical des médias et de la communication, thème central du chapitre 5 et sujet fréquent à l'examen.

\courssub{Discovering the vocabulary in context}

Lisons d'abord ce court texte de presse. Observez les mots en gras : ils appartiennent au champ lexical des médias.

\begin{textebox}[Reading text — “The changing face of news in Cameroon”]
Ten years ago, most Cameroonians got their news from the radio or from a daily \textbf{newspaper} bought at a kiosk in Yaoundé or Douala. Today, the situation has changed dramatically. \textbf{Journalists} still write \textbf{articles} and \textbf{editors} still check every \textbf{issue} before it is \textbf{published}, but many readers no longer buy the printed paper: they visit the newspaper's \textbf{website} instead.

Television remains powerful too. Every evening, millions of viewers watch the \textbf{news bulletin} on their favourite \textbf{channel}, presented by a well-known \textbf{presenter}. When something important happens — an election result, a football victory, a flood in the Far North — the station may interrupt its programmes for \textbf{breaking news} and send a team to make a \textbf{live report} from the scene.

But the biggest revolution is digital. Young people \textbf{post} photos and videos on \textbf{social networks}, \textbf{share} them with their \textbf{followers}, and \textbf{download} podcasts to listen to later. A single video can \textbf{go viral} in a few hours and reach a huge \textbf{audience} without any newspaper or television station being involved.

This new world has a dark side: false information spreads as fast as true information. That is why journalists and citizens must always \textbf{check the facts} and rely on \textbf{reliable sources}. \textbf{Freedom of the press} is precious, but it comes with responsibility.
\end{textebox}

\textbf{Glossaire :} \emph{daily} (adj., quotidien) ; \emph{issue} (n., numéro d'un journal) ; \emph{scene} (n., lieu des faits) ; \emph{to spread} (v., se propager) ; \emph{dark side} (n., côté sombre).

\textbf{Questions de compréhension :}
\begin{enumerate}
\item How did most Cameroonians get their news ten years ago?
\item What do many readers do instead of buying the printed paper?
\item What is the “dark side” of the digital revolution, according to the text?
\end{enumerate}

\courssub{The three families of media vocabulary}

Le vocabulaire des médias s'organise naturellement en trois familles. La carte mentale suivante les résume ; mémorisez-la, elle vous servira de plan mental pour toute composition sur ce thème.

\begin{popfigure}
\begin{tikzpicture}[mindmap, grow cyclic, every node/.style={concept}, root concept/.append style={concept color=popblue, font=\Large\bfseries},
level 1/.append style={level distance=3.6cm, sibling angle=120},
level 2/.append style={level distance=2.4cm, sibling angle=45, concept color=popblueL}]
\node[root concept]{Media}
  child[concept color=poporange] { node {Print}
    child { node {newspaper} }
    child { node {headline} }
    child { node {article} }
    child { node {editor} }
  }
  child[concept color=popgreen] { node {Broadcast}
    child { node {channel} }
    child { node {presenter} }
    child { node {live report} }
    child { node {audience} }
  }
  child[concept color=poppurple] { node {Digital}
    child { node {website} }
    child { node {to post} }
    child { node {follower} }
    child { node {to go viral} }
  };
\end{tikzpicture}
\legende{Carte mentale du champ lexical des médias : trois branches, quatre mots-clés chacune.}
\end{popfigure}

\textbf{1. La presse écrite (Print).}

\begin{definition}[Print media — le vocabulaire de la presse]
\begin{itemize}
\item \cle{newspaper} (n.) : un journal. \eng{I read the newspaper every morning.} « Je lis le journal chaque matin. »
\item \cle{headline} (n.) : le titre (en gros caractères). \eng{The headline announced the election results.} « Le titre annonçait les résultats des élections. »
\item \cle{article} (n.) : un article. \eng{She wrote an article about cocoa farmers.} « Elle a écrit un article sur les planteurs de cacao. »
\item \cle{editor} (n.) : le rédacteur en chef ; \cle{journalist} (n.) : le journaliste.
\item \cle{to publish} (v.) : publier ; \cle{issue} (n.) : un numéro (d'un journal ou magazine). \eng{The magazine publishes a new issue every month.} « Le magazine publie un nouveau numéro chaque mois. »
\end{itemize}
\end{definition}

\textbf{2. L'audiovisuel (Broadcast).}

\begin{definition}[Broadcast media — radio et télévision]
\begin{itemize}
\item \cle{to broadcast} (v.) : diffuser. Attention, verbe irrégulier : broadcast — broadcast — broadcast. \eng{CRTV broadcasts the news in English and French.} « La CRTV diffuse les informations en anglais et en français. »
\item \cle{channel} (n.) : une chaîne (de télévision) ; \cle{news bulletin} (n.) : un bulletin d'informations, un journal télévisé.
\item \cle{presenter} (n.) : le présentateur / la présentatrice. \eng{The presenter reads the news at 8 p.m.} « Le présentateur lit les informations à 20 heures. »
\item \cle{live report} (n.) : un reportage en direct ; \cle{audience} (n.) : le public, l'audience. \eng{The programme attracts a large audience.} « L'émission attire un large public. »
\end{itemize}
\end{definition}

\textbf{3. Le numérique (Digital).}

\begin{definition}[Digital media — les médias numériques]
\begin{itemize}
\item \cle{social network} (n.) : un réseau social ; \cle{to post} (v.) : publier, poster ; \cle{to share} (v.) : partager.
\item \cle{follower} (n.) : un abonné (sur un réseau social). \eng{She has ten thousand followers.} « Elle a dix mille abonnés. »
\item \cle{to go viral} (expr.) : devenir viral. \eng{His song went viral overnight.} « Sa chanson est devenue virale en une nuit. »
\item \cle{website} (n.) : un site web ; \cle{to download} (v.) : télécharger (vers son appareil). À l'inverse, \cle{to upload} = envoyer un fichier vers Internet.
\end{itemize}
\end{definition}

\courssub{Key collocations and idiomatic expressions}

À l'examen, les collocations (mots qui vont naturellement ensemble) font la différence entre un anglais correct et un anglais naturel.

\begin{center}
\begin{tabularx}{\linewidth}{|X|X|}\hline
\rowcolor{popblueL} \textbf{Collocation} & \textbf{Traduction et exemple} \\ \hline
\cle{breaking news} & une information de dernière minute. \eng{We interrupt this programme for breaking news.} \\ \hline
\cle{reliable source} & une source fiable. \eng{A good journalist always uses reliable sources.} \\ \hline
\cle{mass media} & les médias de masse. \eng{Television is part of the mass media.} \\ \hline
\cle{freedom of the press} & la liberté de la presse. \eng{Freedom of the press is essential in a democracy.} \\ \hline
\cle{to check the facts} & vérifier les faits. \eng{Always check the facts before sharing a post.} \\ \hline
\end{tabularx}
\end{center}

\begin{definition}[Mass media]
\eng{Mass media} means \emph{all the means of communication that reach large numbers of people} — newspapers, radio, television and the Internet. En français : les médias de masse. Exemple : \eng{The mass media play a major role in shaping public opinion.} « Les médias de masse jouent un rôle majeur dans la formation de l'opinion publique. »
\end{definition}

Trois expressions idiomatiques à réutiliser dans vos compositions :

\begin{itemize}
\item \cle{to be in the news} : faire l'actualité. \eng{Cameroon was in the news after the football match.} « Le Cameroun a fait l'actualité après le match de football. »
\item \cle{word of mouth} : le bouche-à-oreille. \eng{The restaurant became famous by word of mouth.} « Le restaurant est devenu célèbre par le bouche-à-oreille. »
\item \cle{to get the message across} : faire passer le message. \eng{The campaign got its message across to young people.} « La campagne a fait passer son message auprès des jeunes. »
\end{itemize}

\begin{exemplebox}[Exemples traduits]
\begin{itemize}
\item \eng{Journalists must check their sources.} « Les journalistes doivent vérifier leurs sources. »
\item \eng{The video went viral in a few hours.} « La vidéo est devenue virale en quelques heures. »
\item \eng{The editor refused to publish the article.} « Le rédacteur en chef a refusé de publier l'article. »
\item \eng{Millions of people watched the live report from Buea.} « Des millions de personnes ont regardé le reportage en direct de Buea. »
\end{itemize}
\end{exemplebox}

\begin{savaistu}[“Media” : un pluriel latin]
Le mot \eng{media} est en réalité le pluriel du latin \eng{medium} (un moyen de communication). C'est pourquoi la presse britannique écrit souvent \eng{the media are} au pluriel : \eng{The media are responsible for informing the public.} Dans l'usage courant et américain, on entend aussi \eng{the media is}. À l'écrit, préférez le pluriel : c'est la forme la plus soignée. Même phénomène avec \eng{data}, pluriel de \eng{datum}.
\end{savaistu}

\courssub{Word families : building your vocabulary}

Connaître les familles de mots permet de multiplier son lexique sans effort et de réussir les exercices de transformation de mots fréquents au Bac.

\begin{center}
\begin{tabularx}{\linewidth}{|X|X|X|}\hline
\rowcolor{popblueL} \textbf{Verbe} & \textbf{Nom} & \textbf{Adjectif} \\ \hline
to inform (informer) & information (information) & informative (instructif) \\ \hline
to communicate (communiquer) & communication (communication) & communicative (communicatif) \\ \hline
to publish (publier) & publisher (éditeur) / publication (publication) & published (publié) \\ \hline
to broadcast (diffuser) & broadcast (émission) & broadcast / live (en direct) \\ \hline
to edit (éditer, corriger) & editor (rédacteur en chef) / edition (édition) & editorial (éditorial) \\ \hline
to report (rapporter, signaler) & reporter (journaliste-reporter) / report (reportage) & reported (signalé) \\ \hline
\end{tabularx}
\end{center}

\begin{exemplebox}[Une même famille, trois emplois]
\begin{itemize}
\item \eng{The radio informs the population.} (verbe) « La radio informe la population. »
\item \eng{This information is very useful.} (nom) « Cette information est très utile. »
\item \eng{The documentary was very informative.} (adjectif) « Le documentaire était très instructif. »
\end{itemize}
\end{exemplebox}

\begin{attention}[“Information” est indénombrable]
En français, on dit « des informations », au pluriel. En anglais, \eng{information} est \textbf{indénombrable} : jamais de \emph{-s}, jamais de \emph{an}.
\begin{itemize}
\item Faux : \emph{I need some informations.}
\item Correct : \eng{I need some information.} « J'ai besoin d'informations. »
\item Pour compter : \eng{a piece of information} « une information » ; \eng{two pieces of information} « deux informations ».
\end{itemize}
Même piège avec \eng{news} (singulier malgré le \emph{-s} : \eng{The news is good.}) et \eng{advice} (indénombrable : \eng{a piece of advice}).
\end{attention}

\courssub{False friends : les pièges des francophones}

\begin{attention}[Faux amis à ne plus confondre]
\begin{itemize}
\item \cle{actually} = \textbf{en fait, en réalité} (et non « actuellement »). \eng{He looks young, but he is actually fifty.} « Il paraît jeune, mais il a en fait cinquante ans. » Pour dire « actuellement », utilisez \eng{currently} ou \eng{at present}.
\item \cle{to assist} = \textbf{assister à} (être présent), et non « aider ». \eng{Hundreds of people assisted at the press conference.} Pour dire « aider », utilisez \eng{to help}.
\item \cle{a magazine} = \textbf{un magazine} (revue), et non « un magasin ». « Un magasin » se dit \eng{a shop} (GB) ou \eng{a store} (US).
\item \cle{sensible} = \textbf{sensé, raisonnable} ; \cle{sensitive} = \textbf{sensible} (émotif). \eng{It was a sensible decision.} « C'était une décision sensée. » / \eng{She is very sensitive to criticism.} « Elle est très sensible à la critique. »
\end{itemize}
\end{attention}

\begin{center}
\begin{tabularx}{\linewidth}{|X|X|X|}\hline
\rowcolor{popblueL} \textbf{Français} & \textbf{English} & \textbf{Remarque} \\ \hline
en fait, en réalité & actually & ne signifie pas « actuellement » \\ \hline
actuellement & currently, at present & faux ami avec \eng{actually} \\ \hline
aider & to help & \eng{to assist} = assister à \\ \hline
un magasin & a shop, a store & \eng{a magazine} = une revue \\ \hline
sensé, raisonnable & sensible & \eng{sensitive} = sensible (émotif) \\ \hline
\end{tabularx}
\end{center}

\courssub{Method : enriching your essay vocabulary}

\begin{methode}[Remplacer les mots « pauvres » par des mots « riches »]
Pour gagner des points en expression écrite, remplacez systématiquement les verbes et adjectifs trop simples par des équivalents plus précis :
\begin{enumerate}
\item Repérez dans votre brouillon les mots \eng{say} et \eng{important}.
\item Remplacez \eng{say} selon le sens : \cle{state} (déclarer, formel), \cle{claim} (affirmer, prétendre), \cle{report} (rapporter), \cle{announce} (annoncer).
\item Remplacez \eng{important} par : \cle{crucial}, \cle{vital}, \cle{significant}.
\item Vérifiez la collocation : \eng{to announce the results}, \eng{a crucial issue}, \eng{to claim that...}.
\end{enumerate}
Exemple de transformation : \eng{The minister said that press freedom is important.} devient \eng{The minister stated that press freedom is vital.} « Le ministre a déclaré que la liberté de la presse est vitale. »
\end{methode}

\courssub{Guided practice}

\begin{textebox}[Fill in the blanks — mini-texte à trous]
Last night, the BBC \_\_\_\_\_ (broadcast) a \_\_\_\_\_ (live) report from Lagos. The \_\_\_\_\_ (headline) was about internet access in West Africa. The story quickly went \_\_\_\_\_ (viral) on social media.
\end{textebox}

\begin{exoresolu}[Exercice résolu — le mini-texte à trous]
Complétez le texte ci-dessus avec la forme correcte des mots entre parenthèses.
\tcblower
\textbf{Solution détaillée :}
\begin{enumerate}
\item \eng{broadcast} : le verbe \eng{to broadcast} est irrégulier et \textbf{invariable} (broadcast — broadcast — broadcast). Le marqueur \eng{last night} impose le prétérit, donc \eng{broadcast} (et non \emph{broadcasted}).
\item \eng{live} : déjà un adjectif signifiant « en direct » ; il ne change pas devant le nom \eng{report}.
\item \eng{headline} : l'article défini \eng{the} et l'accord singulier du verbe \eng{was} imposent le singulier.
\item \eng{viral} : expression figée \eng{to go viral} ; l'adjectif reste invariable.
\end{enumerate}
Texte complété : \eng{Last night, the BBC broadcast a live report from Lagos. The headline was about internet access in West Africa. The story quickly went viral on social media.} « Hier soir, la BBC a diffusé un reportage en direct de Lagos. Le titre portait sur l'accès à Internet en Afrique de l'Ouest. L'histoire est rapidement devenue virale sur les réseaux sociaux. »
\end{exoresolu}

\begin{exercice}[Practice — exam training]
\textbf{A. Familles de mots.} Complétez avec la forme correcte du mot entre parenthèses.
\begin{enumerate}
\item The \_\_\_\_\_ of the article took two days. (publish)
\item This documentary is very \_\_\_\_\_. (inform)
\item Good \_\_\_\_\_ is essential between journalists and their sources. (communicate)
\end{enumerate}
\textbf{B. Faux amis.} Traduisez en anglais.
\begin{enumerate}
\item Il a en fait quarante ans.
\item C'est une décision très sensée.
\item Elle a acheté un magazine au kiosque.
\end{enumerate}
\textbf{C. Collocations.} Complétez avec : \eng{breaking}, \eng{reliable}, \eng{mass}, \eng{check}.
\begin{enumerate}
\item Always \_\_\_\_\_ the facts before you share a post.
\item The radio interrupted its music for \_\_\_\_\_ news.
\item A journalist must use \_\_\_\_\_ sources.
\item Television and radio are part of the \_\_\_\_\_ media.
\end{enumerate}
\end{exercice}

\begin{corrige}[Exercice — Practice]
\textbf{A.} 1. \eng{publication} (nom : « la publication ») ; 2. \eng{informative} (adjectif : « instructif ») ; 3. \eng{communication} (nom indénombrable : « la communication »).

\textbf{B.} 1. \eng{He is actually forty.} (\eng{actually} = en fait) ; 2. \eng{It is a very sensible decision.} (\eng{sensible} = sensé) ; 3. \eng{She bought a magazine at the kiosk.} (\eng{a magazine} = une revue).

\textbf{C.} 1. \eng{check} ; 2. \eng{breaking} ; 3. \eng{reliable} ; 4. \eng{mass}.
\end{corrige}

\begin{aretenir}[L'essentiel de la section]
\begin{itemize}
\item Trois familles de médias : \cle{Print} (newspaper, headline, article, editor), \cle{Broadcast} (to broadcast, channel, presenter, live report, audience), \cle{Digital} (to post, to share, follower, to go viral, website, to download).
\item Collocations à réutiliser : \eng{breaking news}, \eng{reliable source}, \eng{mass media}, \eng{freedom of the press}, \eng{to check the facts}.
\item \eng{Information} est indénombrable : \eng{a piece of information}, jamais \emph{informations}.
\item Faux amis : \eng{actually} = en fait ; \eng{to assist} = assister à ; \eng{a magazine} = une revue ; \eng{sensible} = sensé, \eng{sensitive} = sensible.
\item Enrichissez vos copies : \eng{state, claim, report, announce} au lieu de \eng{say} ; \eng{crucial, vital, significant} au lieu de \eng{important}.
\end{itemize}
\end{aretenir}

Votre lexique des médias est désormais consolidé. La section suivante, \emph{Writing — composition de type Bac}, vous montrera comment mobiliser ce vocabulaire dans une production écrite complète, de l'analyse du sujet à la relecture finale.

\coursec{Writing — composition de type Bac}

\courssub{From vocabulary to composition : la transition}

Dans la section précédente, nous avons consolidé le vocabulaire des médias et de la communication : \eng{social media}, \eng{fake news}, \eng{to broadcast}, \eng{to mislead}, \eng{audience}, \eng{headline}... Ce lexique ne doit pas rester une simple liste : il doit \cle{servir} votre production écrite. Au Baccalauréat, l'épreuve d'anglais se termine presque toujours par une \cle{composition} (essay) d'environ 180 à 220 mots. Cette section vous donne la méthode complète pour la réussir : structure, connecteurs, plan en cinq étapes, modèle commenté et grille d'évaluation.

\courssub{La structure de la composition : observer d'abord}

Lisez attentivement le modèle suivant, rédigé sur un sujet typique du Bac. Ne cherchez pas encore à comprendre chaque mot : repérez seulement les \cle{parties} du texte.

\begin{textebox}[Model essay — “Do social media do more harm than good to young people?”]
Nowadays, almost every young Cameroonian owns a smartphone and spends hours on Facebook, WhatsApp or TikTok. But do social media do more harm than good to young people?

Firstly, social media allow young Cameroonians to access information quickly. For instance, students can follow educational pages and share lessons. Moreover, these platforms help families to stay in touch: a student in Yaoundé can chat every day with his uncle in Douala or his cousin in London.

However, these platforms also spread fake news, which can mislead users who do not check the facts. On the other hand, many teenagers become addicted to their screens. As a result, many students waste time online instead of preparing their lessons, and their marks go down.

To conclude, social media are useful tools for information and communication, but they become dangerous when young people use them without control. In my view, schools and parents should teach teenagers how to use these platforms wisely.
\end{textebox}

\textbf{Glossaire}\par
\begin{itemize}
\item \eng{harm} (n.) : le mal, le tort
\item \eng{to stay in touch} : rester en contact
\item \eng{to mislead} : induire en erreur
\item \eng{addicted} (adj.) : accro, dépendant
\item \eng{marks} (n. pl.) : les notes
\item \eng{wisely} (adv.) : sagement, avec discernement
\end{itemize}

\textbf{Questions d'observation}\par
\begin{enumerate}
\item Combien de paragraphes compte le texte ?
\item Quelle phrase du premier paragraphe pose le problème ?
\item Quels mots introduisent les arguments favorables ? Et les arguments défavorables ?
\item Où l'auteur donne-t-il son opinion personnelle ?
\end{enumerate}

Vous l'avez constaté : le texte suit un plan en \cle{quatre mouvements} — introduction, arguments pour, arguments contre, conclusion. C'est exactement le plan attendu à l'examen.

\begin{propriete}[La structure canonique de la composition]
Une composition de type Bac comprend trois grandes parties :
\begin{itemize}
\item \textbf{Introduction} (2-3 phrases) : une phrase d'\cle{accroche} (fait ou question) + présentation du sujet + \cle{problématique} (la question à laquelle le texte va répondre).
\item \textbf{Développement} (2-3 paragraphes) : chaque paragraphe défend une idée, illustrée par un exemple.
\item \textbf{Conclusion} (2-3 phrases) : bilan de la réflexion + \cle{ouverture} (opinion personnelle, conseil, perspective).
\end{itemize}
Règle d'or : \textbf{un paragraphe = une idée principale + un exemple + un connecteur} vers le paragraphe suivant.
\end{propriete}

\begin{popfigure}
\begin{tikzpicture}[node distance=0.55cm and 1.1cm,
box/.style={draw=popblue, thick, rounded corners, fill=popblueL, align=center, text width=4.6cm, minimum height=1.1cm, font=\small},
arr/.style={-{Stealth[length=3mm]}, very thick, poporange}]
\node[box, align=center] (intro){\textbf{Introduction}\\ accroche + sujet\\ + problématique};
\node[box, below=of intro, align=center] (body1){\textbf{Body 1 — arguments for}\\ \emph{Firstly, Moreover...}\\ idée + exemple};
\node[box, below=of body1, align=center] (body2){\textbf{Body 2 — arguments against}\\ \emph{However, On the other hand...}\\ idée + exemple};
\node[box, below=of body2, align=center] (concl){\textbf{Conclusion}\\ \emph{To conclude, In my view...}\\ bilan + ouverture};
\draw[arr] (intro) -- (body1);
\draw[arr] (body1) -- (body2);
\draw[arr] (body2) -- (concl);
\end{tikzpicture}
\legende{Schéma de la composition : Introduction $\rightarrow$ Body 1 (arguments for) $\rightarrow$ Body 2 (arguments against) $\rightarrow$ Conclusion.}
\end{popfigure}

\courssub{Les connecteurs : les articulations du texte}

Les \cle{connecteurs} (linking words) sont les panneaux indicateurs de votre raisonnement. Sans eux, le correcteur lit une suite de phrases isolées ; avec eux, il suit une démonstration. Voici le tableau à connaître par cœur :

\begin{center}
\begin{tabularx}{\linewidth}{|X|X|X|}
\hline
\rowcolor{popblueL} English & Français & Emploi \\
\hline
Firstly / First of all & Premièrement / D'abord & introduire le premier argument \\
\hline
Moreover / Furthermore & De plus / En outre & ajouter un argument du même côté \\
\hline
However / Nevertheless & Cependant / Néanmoins & introduire une opposition \\
\hline
On the other hand & D'un autre côté & présenter le point de vue opposé \\
\hline
For instance / For example & Par exemple & introduire un exemple \\
\hline
As a result / Consequently & Par conséquent & exprimer une conséquence \\
\hline
In my opinion / In my view & À mon avis & donner son opinion personnelle \\
\hline
To conclude / In conclusion & Pour conclure / En conclusion & ouvrir la conclusion \\
\hline
\end{tabularx}
\end{center}

\begin{exemplebox}[Exemples traduits]
\eng{Firstly, social media allow young Cameroonians to access information quickly.} « Premièrement, les réseaux sociaux permettent aux jeunes Camerounais d'accéder rapidement à l'information. »

\eng{Moreover, these platforms help families to stay in touch.} « De plus, ces plateformes aident les familles à rester en contact. »

\eng{As a result, many students waste time online.} « Par conséquent, beaucoup d'élèves perdent du temps en ligne. »

\eng{On the other hand, many teenagers become addicted to their screens.} « D'un autre côté, beaucoup d'adolescents deviennent dépendants de leurs écrans. »

\eng{To conclude, social media are useful tools, but they must be used wisely.} « Pour conclure, les réseaux sociaux sont des outils utiles, mais ils doivent être utilisés avec discernement. »
\end{exemplebox}

\begin{attention}[Ne dites jamais “according to me” !]
Les francophones traduisent souvent « selon moi » par \eng{according to me} : c'est une \textbf{erreur}. En anglais, \eng{according to} sert uniquement à citer une source extérieure : \eng{according to the teacher}, \eng{according to the newspaper}. Pour donner votre opinion, dites \eng{in my opinion}, \eng{in my view} ou \eng{I think that}. Autres calques à bannir : \eng{actually} ne veut pas dire « actuellement » (dites \eng{currently}) ; \eng{to assist} ne veut pas dire « assister à » (dites \eng{to attend}).
\end{attention}

\courssub{Le modèle commenté : repérage des éléments}

Reprenons le texte modèle et annotons-le comme le ferait un correcteur.

\begin{itemize}
\item \textbf{Accroche} : \eng{Nowadays, almost every young Cameroonian owns a smartphone...} — un fait concret et actuel qui capte l'attention.
\item \textbf{Problématique} : \eng{But do social media do more harm than good to young people?} — une question directe qui annonce le débat pour/contre.
\item \textbf{Arguments pour} (Body 1) : accès rapide à l'information (\eng{Firstly}) + lien familial (\eng{Moreover}), avec deux exemples camerounais : les pages éducatives des élèves, la famille répartie entre Yaoundé, Douala et Londres.
\item \textbf{Arguments contre} (Body 2) : les fausses informations (\eng{However}) + l'addiction et la baisse des notes (\eng{On the other hand}, \eng{As a result}).
\item \textbf{Conclusion} : bilan équilibré (\eng{To conclude}) + ouverture sous forme de conseil (\eng{In my view, schools and parents should teach...}).
\item \textbf{Connecteurs repérés} : \eng{Firstly}, \eng{For instance}, \eng{Moreover}, \eng{However}, \eng{On the other hand}, \eng{As a result}, \eng{To conclude}, \eng{In my view} — huit connecteurs pour environ 200 mots : c'est la densité idéale.
\end{itemize}

Remarquez que les exemples sont \cle{locaux et concrets} (Yaoundé, Douala, les élèves, les notes) : un correcteur valorise toujours les exemples tirés du vécu de l'élève plutôt que des généralités vagues.

\courssub{Le plan de rédaction en cinq étapes}

\begin{methode}[Les 5 étapes de la composition]
\begin{enumerate}
\item \textbf{Lire le sujet} (2 min) : soulignez les mots-clés (\eng{social media}, \eng{young people}, \eng{harm}, \eng{good}) et identifiez le type de sujet : débat d'opinion (\eng{Do you agree?}), avantages/inconvénients, ou problème/solutions. Tout votre texte doit répondre à cette question précise.
\item \textbf{Brainstormer} (3 min) : au brouillon, notez en vrac 3 arguments pour et 3 arguments contre, avec un exemple pour chacun. Gardez les deux meilleurs de chaque côté.
\item \textbf{Planifier} (2 min) : répartissez vos idées dans le schéma Introduction / Body 1 / Body 2 / Conclusion et choisissez un connecteur pour chaque paragraphe.
\item \textbf{Rédiger} (25-30 min) : écrivez au propre, une idée par paragraphe, en visant 180 à 220 mots. Phrases simples et correctes plutôt que phrases longues et risquées.
\item \textbf{Relire} (5 min) : vérifiez l'accord sujet-verbe (\eng{media are}, pas \eng{media is}), la troisième personne du singulier (\eng{he wastes}), les temps, l'orthographe et la présence d'une conclusion.
\end{enumerate}
\end{methode}

\begin{methode}[Comment rédiger l'introduction]
L'introduction se construit en trois phrases :
\begin{enumerate}
\item \textbf{Phrase d'accroche} : un fait général ou une question qui accroche le lecteur. \eng{Nowadays, almost every young Cameroonian owns a smartphone.}
\item \textbf{Présentation du sujet} : reformulez le sujet avec vos propres mots. \eng{Social media have become part of the daily life of teenagers.}
\item \textbf{Problématique} : posez la question que votre texte va traiter. \eng{But do these platforms do more harm than good?}
\end{enumerate}
Ne donnez jamais votre opinion dans l'introduction : gardez-la pour la conclusion.
\end{methode}

\courssub{Les erreurs à éviter absolument}

\begin{attention}[Les quatre fautes qui coûtent cher au Bac]
\begin{itemize}
\item \textbf{Le hors-sujet} : si le sujet porte sur les jeunes, ne rédigez pas un texte général sur Internet. Relisez chaque paragraphe en vous demandant : « Est-ce que cela répond à la question ? »
\item \textbf{Les paragraphes non articulés} : une suite de phrases sans connecteurs est illisible. Chaque paragraphe doit s'ouvrir sur un connecteur.
\item \textbf{Les calques du français} : \eng{according to me}, \eng{I am agree} (dites \eng{I agree}), \eng{the medias} (dites \eng{the media} : \eng{media} est déjà le pluriel de \eng{medium}), \eng{informations} (indénombrable : dites \eng{information}).
\item \textbf{L'absence de conclusion} : même si le temps manque, écrivez toujours deux phrases commençant par \eng{To conclude}. Une composition sans conclusion est une composition inachevée, donc pénalisée.
\end{itemize}
\end{attention}

\begin{exoresolu}[Entraînement guidé : compléter un paragraphe]
Énoncé — Complétez ce paragraphe avec les connecteurs suivants : \eng{Moreover}, \eng{For instance}, \eng{As a result}.

\eng{Social media can improve students' English. \ldots, many pupils in Buea watch videos in English every day. \ldots, they learn new words and expressions. \ldots, their results in English tests often improve.}
\tcblower
Solution détaillée — \eng{\textbf{For instance}, many pupils in Buea watch videos in English every day.} (on introduit un exemple concret). \eng{\textbf{Moreover}, they learn new words and expressions.} (on ajoute un argument du même côté). \eng{\textbf{As a result}, their results in English tests often improve.} (on exprime la conséquence). Méthode : demandez-vous quelle relation logique relie chaque phrase à la précédente — exemple, addition ou conséquence — puis choisissez le connecteur correspondant.
\end{exoresolu}

\courssub{La grille d'évaluation : comment vous serez noté}

\begin{center}
\begin{tabularx}{\linewidth}{|l|X|X|}
\hline
\rowcolor{popblueL} Critère & Points & Ce que le correcteur attend \\
\hline
Contenu & 5 & des arguments pertinents qui répondent au sujet, des exemples concrets \\
\hline
Organisation & 5 & introduction, développement en paragraphes, conclusion ; connecteurs variés \\
\hline
Vocabulaire & 5 & lexique riche et précis (ici : médias et communication), pas de répétitions \\
\hline
Grammaire et orthographe & 5 & accords corrects, temps maîtrisés, orthographe soignée \\
\hline
\end{tabularx}
\end{center}

\begin{savaistu}[La longueur idéale au Bac]
Au Baccalauréat, une composition claire de 180 à 220 mots bien structurée vaut mieux qu'un texte long et désorganisé. Le correcteur note la \cle{qualité} du raisonnement, pas la quantité d'encre. Un texte de 200 mots avec une problématique nette, deux paragraphes articulés et une conclusion obtient presque toujours une meilleure note qu'un pavé de 400 mots sans plan. Écrire moins, mais écrire mieux : c'est la stratégie gagnante.
\end{savaistu}

\begin{exercice}[À vous de jouer]
Rédigez une composition de 180 à 220 mots sur le sujet suivant : \eng{“Does television still play an important role in the life of young Cameroonians?”} Suivez le plan en cinq étapes : lisez le sujet, notez trois arguments pour et trois contre au brouillon, planifiez vos paragraphes, rédigez avec au moins six connecteurs différents, puis relisez-vous avec la grille d'évaluation.
\end{exercice}

\begin{corrige}[Exercice — éléments de réponse]
Introduction attendue : accroche sur la place de la télévision dans les foyers camerounais + problématique (\eng{But does television still matter to young people?}). Arguments pour possibles : les informations du soir regardées en famille, les matchs de football des Lions Indomptables, les programmes éducatifs. Arguments contre possibles : la concurrence des smartphones et d'Internet, les programmes peu adaptés aux jeunes, le temps perdu devant l'écran. Conclusion attendue : bilan (\eng{To conclude, television remains important for family life...}) + ouverture (\eng{In my view, young people should choose their programmes carefully.}). Connecteurs attendus : \eng{Firstly}, \eng{Moreover}, \eng{However}, \eng{On the other hand}, \eng{As a result}, \eng{To conclude}. Vérifiez : \eng{television plays} (s à la 3e personne), \eng{information} sans s, pas de \eng{according to me}.
\end{corrige}

\coursec{Exam strategy and final training}

Après avoir révisé la composition de type Bac dans la section précédente, il reste une dernière étape, peut-être la plus importante : apprendre à \cle{gérer l'épreuve elle-même}. Beaucoup d'élèves connaissent leurs leçons mais perdent des points par mauvaise organisation, par précipitation ou par manque de relecture. Cette dernière partie vous donne une méthode complète pour le jour J, un entraînement blanc guidé et une grille d'auto-évaluation.

\courssub{Managing your time during the exam}

Une épreuve d'anglais de deux heures se découpe en sections bien identifiées. La première compétence de l'examinateur, c'est vous : celui qui sait répartir son temps marque déjà des points.

\begin{definition}[Time management]
\eng{Time management} is the ability to divide the available time between the different parts of an exam paper and to respect that plan. En français : la gestion du temps, c'est-à-dire la répartition planifiée du temps entre les sections de l'épreuve.
\end{definition}

Voici une répartition indicative pour une épreuve de deux heures :

\begin{center}
\begin{tabularx}{\linewidth}{|X|l|X|}
\hline
\rowcolor{popblueL} \textbf{Section} & \textbf{Temps} & \textbf{Conseil} \\
\hline
Reading comprehension & 30 min & Lire le texte deux fois, souligner les mots-clés \\
\hline
Grammar & 25 min & Relire chaque phrase après transformation \\
\hline
Vocabulary & 15 min & Vérifier le contexte avant de choisir le mot \\
\hline
Composition (writing) & 40 min & 10 min de plan, 25 min de rédaction, 5 min de relecture \\
\hline
Final proofreading & 10 min & Relecture générale de toute la copie \\
\hline
\end{tabularx}
\end{center}

\begin{methode}[Méthode pour le jour J]
\begin{enumerate}
\item \textbf{Lire le sujet en entier} avant d'écrire quoi que ce soit : cela évite les mauvaises surprises et permet au cerveau de commencer à travailler sur la composition pendant que vous traitez la grammaire.
\item \textbf{Souligner les mots-clés des consignes} : \eng{answer in complete sentences}, \eng{justify with the text}, \eng{write 150--200 words}. Une consigne mal lue, c'est une question ratée même avec une bonne réponse.
\item \textbf{Commencer par la partie la mieux maîtrisée} : si le vocabulaire est votre point fort, traitez-le d'abord. Vous gagnez en confiance et en rapidité.
\item \textbf{Ne jamais laisser de case vide} : une réponse imparfaite peut rapporter des points ; une case vide n'en rapporte jamais.
\item \textbf{Réserver au moins dix minutes de relecture} à la fin, quoi qu'il arrive.
\end{enumerate}
\end{methode}

\begin{popfigure}
\begin{tikzpicture}[
node distance=0.55cm,
box/.style={rectangle, rounded corners, draw=popblue, fill=popblueL, text width=5.2cm, align=center, minimum height=1cm, font=\small},
arr/.style={-{Stealth[length=3mm]}, thick, popdark}]
\node[box] (a) {\textbf{Step 1} — Read all the instructions};
\node[box, below=of a] (b) {\textbf{Step 2} — Spot the easy questions};
\node[box, below=of b] (c) {\textbf{Step 3} — Work section by section};
\node[box, below=of c, draw=poporange, fill=poporangeL] (d) {\textbf{Step 4} — Keep 10 minutes for proofreading};
\draw[arr] (a) -- (b);
\draw[arr] (b) -- (c);
\draw[arr] (c) -- (d);
\end{tikzpicture}
\legende{Exam strategy : la démarche gagnante en quatre étapes.}
\end{popfigure}

\courssub{Proofreading : the art of checking your copy}

\begin{definition}[Proofreading]
\eng{Proofreading} is the final careful reading of a text to correct mistakes before handing it in. En français : la relecture corrective. Ce n'est pas relire « pour voir » : c'est relire avec une \cle{liste de contrôle} précise.
\end{definition}

La relecture efficace suit un ordre précis. Voici la check-list du candidat sérieux :

\begin{center}
\begin{tabularx}{\linewidth}{|X|X|X|}
\hline
\rowcolor{popblueL} \textbf{Point à vérifier} & \textbf{Erreur typique} & \textbf{Correction} \\
\hline
3e personne du singulier (-s) & \eng{She watch TV} & \eng{She watches TV} \\
\hline
Concordance des temps & \eng{He said he is tired} & \eng{He said he was tired} \\
\hline
Accord sujet-verbe & \eng{The news are good} & \eng{The news is good} \\
\hline
Orthographe des mots du cours & \eng{advertisment} & \eng{advertisement} \\
\hline
A / an devant voyelle & \eng{a interview} & \eng{an interview} \\
\hline
Majuscules (I, jours, mois) & \eng{i watch tv on monday} & \eng{I watch TV on Monday} \\
\hline
\end{tabularx}
\end{center}

\begin{attention}[Le -s de la 3e personne : l'erreur n° 1 des francophones]
En français, le « s » de la troisième personne du singulier ne s'entend pas (\emph{il regarde}) ; en anglais, il est \textbf{obligatoire} à l'écrit : \eng{He watches the news every evening.} (« Il regarde les informations chaque soir. ») À la relecture, passez votre doigt sous chaque verbe au present simple et demandez-vous : « le sujet est-il he, she ou it ? » Si oui, le verbe porte -s ou -es. Attention aussi au négatif et à l'interrogatif : avec \eng{does}, le -s migre sur l'auxiliaire : \eng{Does she watch TV?} et jamais \eng{Does she watches TV?}
\end{attention}

\begin{exemplebox}[Observer avant de rédiger]
Observez cette phrase modèle, typique du chapitre Media and Communication :

\eng{The news was broadcast yesterday evening.} « Les informations ont été diffusées hier soir. »

Trois points à vérifier en relecture : (1) \eng{news} est singulier en anglais, donc \eng{was} et non \eng{were} ; (2) \eng{broadcast} est un verbe irrégulier invariable : \eng{broadcast -- broadcast -- broadcast} ; (3) \eng{yesterday evening} est un repère daté qui impose le past simple, jamais le present perfect. Comparaison avec le français : « les informations » est pluriel en français, singulier en anglais — piège classique.
\end{exemplebox}

\begin{savaistu}[Les examinateurs ne sont pas vos ennemis]
Les correcteurs du Bac et du Probatoire appliquent une grille de notation \cle{positive} : ils cherchent ce que vous savez, pas seulement ce que vous ignorez. Une réponse pertinente avec une petite erreur de langue reçoit presque toujours une partie des points. Morale : tentez toujours une réponse, même approximative. Un candidat qui écrit \eng{The journalist interview the minister} gagne des points pour l'idée ; celui qui laisse la case vide n'en gagne aucun.
\end{savaistu}

\begin{attention}[Justifier ses réponses de compréhension]
Une réponse de compréhension non justifiée par le texte \textbf{perd la moitié des points}. Quand la consigne dit \eng{justify your answer with the text}, vous devez citer ou paraphraser le passage qui prouve votre réponse. Exemple : \eng{Is the journalist optimistic? — No, he is not. He says that “many young people no longer trust the media”.} Sans la citation, la réponse \eng{No, he is not} ne vaut que la moitié des points.
\end{attention}

\courssub{Mock exam : a 30-minute training paper}

\begin{textebox}[Reading text — “The new face of news in Cameroon”]
Not so long ago, most Cameroonians got their news from two sources: the radio in the morning and the television news at eight o'clock in the evening. Today, the picture has changed completely. In Yaoundé and Douala, young people check their phones before they even get out of bed. News travels through WhatsApp groups, Facebook pages and short videos shared by friends.

This revolution has brought real advantages. Information is faster and cheaper than ever before. A student in Bamenda can follow events in Buea in real time, and citizen journalists can report stories that big newspapers ignore. During local events, ordinary people film what happens and share it within minutes.

However, there is a dark side. False information spreads as quickly as true information, and many readers cannot tell the difference. “We receive dozens of messages every day, and we do not always know what to believe,” explains a secondary school student in Douala. Teachers now encourage their students to check sources, compare several websites and ask who wrote the article and why.

Traditional media are fighting back. Radio stations have created their own applications, and newspapers now publish online editions. The future will probably belong to those who can combine the speed of social media with the seriousness of professional journalism. As one editor in Yaoundé puts it: “Fast news is good, but true news is better.”

\textbf{Glossary} — \emph{to get out of bed} (v.) : sortir du lit ; \emph{to spread} (v., \emph{spread -- spread -- spread}) : se répandre ; \emph{to tell the difference} (exp.) : faire la différence ; \emph{source} (n.) : source ; \emph{to fight back} (v.) : riposter, réagir ; \emph{editor} (n.) : rédacteur en chef ; \emph{to put it} (exp.) : le dire, s'exprimer ainsi.
\end{textebox}

\begin{exercice}[Mini-épreuve blanche — 30 minutes]
\textbf{A. Reading comprehension (10 min).} Answer in complete sentences and justify with the text.
\begin{enumerate}
\item How did most Cameroonians get their news in the past?
\item Give two advantages of the new media mentioned in the text.
\item What problem does the student in Douala mention?
\item Is the editor in Yaoundé optimistic about the future? Justify.
\end{enumerate}
\textbf{B. Grammar (8 min).}
\begin{enumerate}
\item Put into the passive: \eng{Journalists broadcast the news every evening.}
\item Put into reported speech: \eng{“I check my phone every morning,” the student said.}
\item Complete with the right tense: \eng{If young people (check) their sources, they (not/believe) false news.}
\end{enumerate}
\textbf{C. Vocabulary (5 min).} Find in the text or from your lessons: (1) the opposite of \eng{false information}; (2) a synonym of \eng{newspaper article}; (3) the person who reports news for a newspaper.
\textbf{D. Writing (7 min).} In five to eight sentences, answer: \eng{Do you prefer getting news from social media or from traditional media? Why?}
\end{exercice}

\begin{exoresolu}[Correction guidée — sections A et B]
Répondez aux questions A1 et B1--B3 de la mini-épreuve ci-dessus, en rédigeant des réponses complètes et justifiées.
\tcblower
\textbf{A1.} \eng{In the past, most Cameroonians got their news from the radio in the morning and the television news at eight o'clock in the evening.} Justification : le texte dit \eng{“most Cameroonians got their news from two sources: the radio... and the television news”}. Notez le past simple \eng{got} (verbe irrégulier \eng{get -- got -- got}) car le repère \eng{in the past} est daté.

\textbf{B1.} \eng{The news is broadcast every evening (by journalists).} Transformation au passif : le complément d'objet \eng{the news} devient sujet ; auxiliaire \eng{be} au présent (\eng{is}) + participe passé \eng{broadcast} (invariable). Erreur à éviter : \eng{is broadcasted} n'existe pas.

\textbf{B2.} \eng{The student said (that) he checked his phone every morning.} Discours indirect : \eng{said} au passé entraîne le recul du temps (\eng{check} devient \eng{checked}) et le changement de pronom (\eng{I} devient \eng{he}, \eng{my} devient \eng{his}).

\textbf{B3.} \eng{If young people check their sources, they will not believe false news.} Structure du premier conditionnel : \eng{if} + present simple, futur dans la proposition principale. Piège : jamais \eng{will} dans la proposition introduite par \eng{if}.
\end{exoresolu}

\begin{corrige}[Sections C et D]
\textbf{C.} (1) \eng{true information} ; (2) \eng{a report / a story / an article} ; (3) \eng{a journalist} (ou \eng{a reporter}).
\textbf{D.} Modèle de réponse courte : \eng{I prefer getting news from traditional media because professional journalists check their sources before publishing. Social media are faster, but false information spreads very quickly there. In my opinion, the best solution is to use both: social media for speed and traditional media for reliability.} (« Je préfère m'informer par les médias traditionnels parce que les journalistes professionnels vérifient leurs sources avant de publier. Les réseaux sociaux sont plus rapides, mais les fausses informations s'y répandent très vite. À mon avis, la meilleure solution est d'utiliser les deux. ») Barème indicatif : 1 point pour la position claire, 1 point pour l'argument, 1 point pour la correction grammaticale.
\end{corrige}

\courssub{Self-evaluation : what do I still need to revise?}

Terminez l'année par un bilan honnête. Cochez chaque ligne : c'est votre plan de révision personnel avant l'examen.

\begin{center}
\begin{tabularx}{\linewidth}{|X|c|c|}
\hline
\rowcolor{popblueL} \textbf{Point de l'année} & \textbf{Je sais} & \textbf{Je dois revoir} \\
\hline
Present simple / present continuous & & \\
\hline
Past simple / present perfect & & \\
\hline
Futur (will / going to) et conditionnel & & \\
\hline
Voix passive (\eng{The news is broadcast}) & & \\
\hline
Discours indirect (reported speech) & & \\
\hline
Comparatifs et superlatifs & & \\
\hline
Vocabulaire des médias (press, broadcast, headline...) & & \\
\hline
Rédiger une composition structurée (plan, paragraphes) & & \\
\hline
Justifier une réponse par le texte & & \\
\hline
\end{tabularx}
\end{center}

\begin{experience}[Correction collective et débriefing]
En classe entière : (1) un élève écrit sa réponse de compréhension au tableau ; (2) la classe vérifie deux choses — la réponse est-elle juste, et est-elle \textbf{justifiée par le texte} ? (3) chaque élève remplit ensuite sa grille d'auto-évaluation et choisit \textbf{trois} points à revoir en priorité la veille de l'examen. L'enseignant relève au tableau les trois erreurs les plus fréquentes et les corrige collectivement.
\end{experience}

\begin{aretenir}[Les cinq réflexes du candidat gagnant]
\begin{itemize}
\item Lire tout le sujet et souligner les mots-clés des consignes avant de commencer.
\item Commencer par la section la mieux maîtrisée.
\item Justifier chaque réponse de compréhension par le texte.
\item Ne jamais laisser de case vide : toute réponse pertinente rapporte des points.
\item Garder dix minutes de \cle{proofreading} : -s de la 3e personne, temps, accords, orthographe.
\end{itemize}
\end{aretenir}

Vous avez désormais révisé l'ensemble du chapitre Media and Communication et consolidé les points majeurs de l'année de Terminale A. La réussite à l'examen ne dépend pas seulement de ce que vous savez, mais de la manière dont vous le montrez : méthode, temps, relecture. \eng{Good luck — and remember: fast news is good, but true news is better!}

\newpage

\coursec{Activité d'intégration}

\courssub{Objectifs de l'activité}

Cette activité d'intégration clôt le chapitre \eng{Media and Communication} et sert de répétition générale avant l'épreuve écrite du baccalauréat. À l'issue de l'activité, l'élève sera capable de :

\begin{itemize}
\item mobiliser simultanément les acquis de l'année : \cle{passif}, \cle{discours rapporté}, \cle{temps verbaux} (present perfect, past simple, futurs), \cle{connecteurs logiques} et lexique des médias ;
\item comprendre et exploiter un document d'actualité en anglais (reading) ;
\item rédiger en équipe la une d'un journal anglophone : titre accrocheur, court article au passif, interview rapportée (writing) ;
\item présenter oralement une production en deux minutes avec une prononciation claire (speaking) ;
\item évaluer la production des autres groupes à l'aide d'une grille de critères, puis corriger collectivement les erreurs récurrentes.
\end{itemize}

\courssub{Situation}

Votre classe de Terminale A a été choisie par la rédaction fictive du journal anglophone \emph{The Cameroon Herald} pour produire une édition spéciale « jeunes reporters ». Chaque groupe de quatre élèves est une équipe de journalistes. Vous disposez de la dépêche d'agence suivante, reçue ce matin à la rédaction de Yaoundé.

\begin{textebox}[Agency newsflash — The Cameroon Herald, newsroom, Yaoundé]
YAOUNDÉ, Tuesday 7 a.m. — A new community radio station, \emph{Voice of the Quarter}, was officially launched yesterday in the Briqueterie neighbourhood. The project, which is funded by local shopkeepers and a youth association, was welcomed by residents who say local news has been ignored by national channels for too long. Programmes will be broadcast in English, Pidgin and French. The station manager, Mrs. Etoundi, told reporters: “We have trained twenty young volunteers, and our first talk show will be devoted to the fight against fake news on social media.” The launch was attended by the divisional officer and was widely shared on Facebook and WhatsApp. Media experts believe that community radio can play a key role in educating the public, but they warn that journalists must check their sources before publishing any information.
\end{textebox}

\begin{definition}[Glossaire utile]
\begin{itemize}
\item \eng{newsflash} (n.) : dépêche, flash d'information ;
\item \eng{to launch} (v.) : lancer, inaugurer ;
\item \eng{to be funded by} (v. passif) : être financé par ;
\item \eng{residents} (n.) : habitants ;
\item \eng{to broadcast} (v.) : diffuser (irrégulier : \eng{broadcast -- broadcast -- broadcast}) ;
\item \eng{to attend} (v.) : assister à (attention : pas de \eng{to} après ce verbe) ;
\item \eng{fake news} (n.) : fausses informations ;
\item \eng{a source} (n.) : une source (d'information).
\end{itemize}
\end{definition}

\courssub{Consignes}

Travail en groupes de quatre. Répartissez les rôles : un \cle{rédacteur en chef} (coordinateur), deux \cle{journalistes} (rédaction), un \cle{secrétaire de rédaction} (relecture et grammaire). Durée totale : 50 minutes de préparation, puis 2 minutes de présentation orale par groupe.

\begin{enumerate}
\item \textbf{Comprendre le document (10 min).} Lisez la dépêche et répondez par écrit, en phrases complètes : \eng{What happened? Where and when? Who is Mrs. Etoundi? What do media experts warn journalists about?} Soulignez dans le texte trois verbes au passif et deux verbes au present perfect.
\item \textbf{Transformer la langue (10 min).} (a) Mettez à la forme passive : \eng{Local shopkeepers fund the project.} (b) Rapportez au discours indirect la déclaration de Mrs. Etoundi : \eng{“We have trained twenty young volunteers, and our first talk show will be devoted to the fight against fake news.”} (c) Trouvez un titre accrocheur de moins de huit mots, en style journalistique anglais (présent simple, articles souvent omis).
\item \textbf{Rédiger la une (20 min).} Produisez : un \cle{titre}, un court article de 80 à 100 mots (au moins deux phrases au passif et deux verbes au present perfect), et une \cle{interview rapportée} d'un témoin imaginaire (3 à 4 phrases au discours indirect : \eng{He said that... / He explained that... / He added that...}).
\item \textbf{Présenter à l'oral (2 min par groupe).} Chaque groupe lit sa une à la classe : le titre est annoncé avec énergie, l'article est lu clairement, l'interview est restituée. Les autres groupes notent les productions selon la grille vue en section 5 : contenu (4 pts), grammaire et temps (4 pts), vocabulaire des médias (4 pts), prononciation et aisance (4 pts), titre et mise en page (4 pts).
\item \textbf{Vote et correction collective (10 min).} La classe vote pour la meilleure une. L'enseignant relève au tableau les erreurs récurrentes entendues (passif, discours rapporté, prépositions) et la classe les corrige ensemble.
\end{enumerate}

\begin{experience}[Rôles et gestion du temps]
Minutez chaque étape au tableau. Pendant les présentations orales, imposez le silence et la prise de notes : chaque élève doit remplir la grille d'évaluation pour au moins deux groupes. Cela entraîne l'écoute active (listening) en conditions d'examen.
\end{experience}

\courssub{Ressources à mobiliser}

\begin{propriete}[Rappel 1 — Le passif]
Forme : \textbf{sujet + be (au temps voulu) + participe passé} (+ \eng{by} + agent si nécessaire). Le passif est le temps roi de la presse anglaise, car l'action compte plus que l'auteur.\\
\eng{The station was launched yesterday.} « La station a été inaugurée hier. »\\
\eng{Twenty volunteers have been trained.} « Vingt volontaires ont été formés. »
\end{propriete}

\begin{propriete}[Rappel 2 — Le discours rapporté (reported speech)]
Verbe introducteur au passé $\Rightarrow$ recul des temps : present simple $\rightarrow$ past simple ; present perfect $\rightarrow$ past perfect ; will $\rightarrow$ would. Les pronoms et les marqueurs de temps changent aussi (\eng{today} $\rightarrow$ \eng{that day}, \eng{yesterday} $\rightarrow$ \eng{the day before}).\\
\eng{“We have trained twenty volunteers,” she said.} $\rightarrow$ \eng{She said (that) they had trained twenty volunteers.}
\end{propriete}

\begin{propriete}[Rappel 3 — Titres de presse et connecteurs]
Titres anglais : présent simple, articles et auxiliaires souvent supprimés : \eng{New Radio Station Opens in Yaoundé}. Connecteurs utiles pour l'article : \eng{however, moreover, according to, as a result, finally}.
\end{propriete}

\begin{attention}[Erreurs fréquentes des francophones]
\begin{itemize}
\item \eng{information} : en anglais, \eng{information} est \textbf{indénombrable} : \eng{a piece of information}, jamais de pluriel.
\item \eng{to assist to} : faux ami ; on dit \eng{to attend an event} (« assister à ») ; \eng{to assist} signifie « aider ».
\item Passif : ne pas oublier l'auxiliaire \eng{be} : \eng{was launched}, et non \eng{was launch}.
\item Discours rapporté : ne pas laisser \eng{will} après un verbe au passé : \eng{she said the show would be...}, et non \eng{will be}.
\item Prononciation de \eng{news} : « niouz », invariable (pas de \eng{a news}).
\end{itemize}
\end{attention}

\courssub{Proposition de production et corrigé commenté}

\begin{exoresolu}[Une modèle de « The Cameroon Herald »]
Énoncé : rédigez la une complète (titre, article de 80 à 100 mots, interview rapportée) à partir de la dépêche.
\tcblower
\textbf{Production modèle :}

New Community Radio Gives Youth a Voice

A new community radio station, Voice of the Quarter, was officially launched on Monday in the Briqueterie neighbourhood of Yaoundé. The project is funded by local shopkeepers and a youth association, and twenty young volunteers have already been trained. According to the station manager, programmes will be broadcast in English, Pidgin and French. The launch was attended by the divisional officer and has been widely shared on social media. Media experts believe community radio can educate the public; however, they warn that sources must always be checked before any information is published.

Interview — A resident, Mr. Mbarga, said that local news had been ignored by national channels for too long. He explained that the new station would give young people a voice, and he added that he was proud of the volunteers.

\textbf{Commentaire des choix de langue :}
\begin{itemize}
\item \textbf{Passif} (exigé par la consigne) : \eng{was launched}, \eng{is funded}, \eng{have been trained}, \eng{was attended}, \eng{has been shared}, \eng{must be checked} — six occurrences, aux temps variés (past simple, present simple, present perfect, modal).
\item \textbf{Present perfect} : \eng{have already been trained}, \eng{has been widely shared} pour des actions récentes dont le résultat compte.
\item \textbf{Discours rapporté} : \eng{said that local news had been ignored} (present perfect $\rightarrow$ past perfect) ; \eng{explained that the new station would give} (will $\rightarrow$ would) ; trois verbes introducteurs variés (\eng{said, explained, added}).
\item \textbf{Connecteurs et lexique des médias} : \eng{according to, however, sources, broadcast, social media}.
\item \textbf{Titre} : présent simple, ton accrocheur, moins de huit mots.
\end{itemize}
\end{exoresolu}

\courssub{Bilan de l'activité}

Cette activité a permis de réinvestir, dans une tâche unique et motivante, l'essentiel du programme de l'année : compréhension d'un texte de presse, passif, discours rapporté, present perfect, connecteurs logiques et vocabulaire des médias, ainsi que l'expression orale en conditions limitées de temps — exactement comme à l'examen. L'évaluation par les pairs et la correction collective des erreurs récurrentes (oubli de l'auxiliaire \eng{be} au passif, absence de recul des temps, faux amis comme \eng{assist}) renforcent la conscience linguistique de chaque élève. Pour prolonger l'entraînement à la maison, chaque élève pourra rédiger individuellement une deuxième une sur un fait d'actualité de son quartier, en respectant la même grille de critères : c'est la meilleure préparation à l'épreuve de composition du baccalauréat.

\newpage

\coursec{Méthodes et exercices résolus}

\coursec{Lesson 61 — General revision}

\courssub{Savoir-faire 1 : Réussir la compréhension de texte (Reading comprehension)}

\begin{methode}[Réussir la compréhension de texte en 5 étapes]
\begin{enumerate}
\item \textbf{Lire d'abord les questions}, pas le texte. Souligne les mots-clés de chaque question (\eng{why, when, according to the author, in your own words...}).
\item \textbf{Lire le texte deux fois} : une lecture globale pour le sens général, une lecture détaillée pour repérer les passages qui répondent aux questions.
\item \textbf{Répondre par des phrases complètes} en anglais, en reprenant les termes de la question : \eng{The journalist uses social media because...} et non une simple phrase copiée sans lien.
\item \textbf{Ne jamais recopier mot pour mot} quand la consigne dit \eng{in your own words} : reformule avec des synonymes (\eng{young people} $\rightarrow$ \eng{teenagers} ; \eng{very important} $\rightarrow$ \eng{essential}).
\item \textbf{Pour les questions de vocabulaire} (\eng{Find in the text a word meaning...}), cherche le synonyme dans le paragraphe indiqué et vérifie que la nature du mot correspond (nom, verbe, adjectif).
\end{enumerate}
\textbf{Réflexe final :} relis chaque réponse pour vérifier l'accord sujet-verbe et le temps (le texte est souvent au présent : réponds au présent).
\end{methode}

\begin{textebox}[Texte support — Reading comprehension]
In many Cameroonian towns, young people no longer buy newspapers. They prefer to read the news on their phones, through social media such as Facebook and WhatsApp. This change worries some teachers, who fear that students believe everything they read online. “Fake news spreads faster than the truth,” says a headmaster in Buea. However, digital media also has advantages: information is free, fast and available everywhere, even in remote villages.
\end{textebox}

\begin{exoresolu}[Reading comprehension — type Bac]
\textbf{Questions.}
\begin{enumerate}
\item According to the text, how do young Cameroonians read the news today? (2 marks)
\item Why are some teachers worried? Answer in your own words. (2 marks)
\item Find in the text a word or phrase meaning: (a) \eng{false information} ; (b) \eng{far from towns}. (2 marks)
\item Give two advantages of digital media mentioned in the text. (2 marks)
\end{enumerate}
\tcblower
\textbf{Raisonnement.} Question 1 : le premier paragraphe donne directement la réponse (\eng{on their phones, through social media}). Question 2 : la consigne impose une reformulation ; le texte dit \eng{students believe everything they read online}, il faut le dire autrement. Question 3 : on cherche un synonyme dans le texte, en respectant la nature du mot. Question 4 : le dernier paragraphe énumère \eng{free, fast, available everywhere}.

\textbf{Réponses rédigées.}
\begin{enumerate}
\item \eng{Today, young Cameroonians read the news on their phones, mainly through social media such as Facebook and WhatsApp.} (Aujourd'hui, les jeunes Camerounais lisent les informations sur leur téléphone, principalement via les réseaux sociaux.)
\item \eng{Some teachers are worried because students may accept as true everything they find on the Internet, without checking it.} (Reformulation : \eng{believe} $\rightarrow$ \eng{accept as true} ; \eng{read online} $\rightarrow$ \eng{find on the Internet}.)
\item (a) \eng{fake news} ; (b) \eng{remote} (adjectif signifiant « éloigné, isolé »).
\item \eng{Digital media is free and fast, and information is available everywhere, even in remote villages.}
\end{enumerate}
\textbf{Conclusion.} Chaque réponse est une phrase complète, au présent comme le texte, avec un verbe conjugué et sans recopie inutile : c'est ce que le correcteur attend.
\end{exoresolu}

\courssub{Savoir-faire 2 : Maîtriser la synthèse des temps verbaux}

\begin{propriete}[Synthèse des temps verbaux — Tableau récapitulatif]
\begin{center}
\begin{tabularx}{\linewidth}{|X|X|X|X|X|}
\hline
\rowcolor{popblueL}
\textbf{Tense} & \textbf{Form} & \textbf{Main Use} & \textbf{Key Markers} & \textbf{Example} \\
\hline
Present Simple & \eng{V / V-s (3rd pers.)} & Habits, general truths & \eng{every day, usually, often} & \eng{She reads the news daily.} \\
\hline
Present Continuous & \eng{am/is/are + V-ing} & Action now, temporary situation & \eng{now, at the moment, Look!} & \eng{He is interviewing the mayor.} \\
\hline
Simple Past & \eng{V-ed / 2nd column} & Finished past action & \eng{yesterday, last week, ago} & \eng{They broadcast the match yesterday.} \\
\hline
Present Perfect & \eng{have/has + V-en/3rd col.} & Past action linked to present (experience, duration, result) & \eng{since, for, already, just, ever, never} & \eng{I have never visited the USA.} \\
\hline
Past Perfect & \eng{had + V-en/3rd col.} & Action before another past action & \eng{when, after, before, by the time} & \eng{The programme had started when we arrived.} \\
\hline
Future (\eng{will}) & \eng{will + BV} & Prediction, spontaneous decision & \eng{tomorrow, next year, in 2030} & \eng{CRTV will broadcast the final.} \\
\hline
Future (\eng{be going to}) & \eng{am/is/are going to + BV} & Intention, plan, evidence & \eng{tomorrow, soon, later} & \eng{She is going to write an article.} \\
\hline
Conditional 1 (Real) & \eng{If + Present, will + BV} & Real possibility in future & \eng{If + present} & \eng{If he invests, people will be informed.} \\
\hline
Conditional 2 (Unreal present) & \eng{If + Past, would + BV} & Hypothetical present/future & \eng{If + past} & \eng{If I had time, I would read more.} \\
\hline
Conditional 3 (Unreal past) & \eng{If + Past Perfect, would have + V-en} & Regret, hypothesis about past & \eng{If + had + V-en} & \eng{If she had checked, she wouldn't have shared it.} \\
\hline
\end{tabularx}
\end{center}
\textbf{Règle d'or :} Le \cle{present perfect} ne s'emploie \emph{jamais} avec un marqueur de temps passé précis (\eng{yesterday, last week, in 2020, ago}). On utilise le \cle{simple past} : \eng{I saw him yesterday} (et non \eng{*I have seen him yesterday}).
\end{propriete}

\begin{methode}[Choisir le bon temps en 4 questions]
Face à un verbe à conjuguer, pose-toi ces questions dans l'ordre :
\begin{enumerate}
\item \textbf{Y a-t-il un marqueur de temps ?} \eng{now, at the moment} $\rightarrow$ present continuous ; \eng{yesterday, last week, ago} $\rightarrow$ simple past ; \eng{since, for, already, just, ever, never} $\rightarrow$ present perfect ; \eng{tomorrow, next year} $\rightarrow$ futur (\eng{will} ou \eng{be going to}).
\item \textbf{L'action est-elle habituelle ou générale ?} $\rightarrow$ present simple (\eng{She reads the news every morning.}).
\item \textbf{Deux actions passées ?} La plus ancienne au past perfect (\eng{had + participe passé}), la plus récente au simple past : \eng{When the journalist arrived, the minister had already left.}
\item \textbf{Phrase conditionnelle ou hypothétique ?} \eng{If + present, will + BV} (réel) ; \eng{If + past, would + BV} (irréel du présent) ; \eng{If + past perfect, would have + participe} (irréel du passé).
\end{enumerate}
\end{methode}

\begin{exoresolu}[Grammar — Put the verbs in the correct tense]
\textbf{Énoncé.} Complete with the correct tense of the verbs in brackets.
\begin{enumerate}
\item Every evening, my father (watch) the news on CRTV.
\item Look! The reporter (interview) the mayor of Douala.
\item I (never / visit) the United States, but I would like to go there one day.
\item When we arrived at the studio, the programme (already / start).
\item If the government invests in local radio, more villagers (be) informed.
\item She (work) for this newspaper since 2019.
\end{enumerate}
\tcblower
\textbf{Raisonnement et réponses.}
\begin{enumerate}
\item \eng{Every evening} = habitude $\rightarrow$ present simple, 3e personne : \eng{watches}. (Chaque soir, mon père regarde les informations sur CRTV.)
\item \eng{Look!} = action en cours $\rightarrow$ present continuous : \eng{is interviewing}. (Regarde ! Le journaliste interviewe le maire de Douala.)
\item \eng{Never} + expérience de vie sans date $\rightarrow$ present perfect : \eng{have never visited}. (Je n'ai jamais visité les États-Unis.)
\item Action antérieure à \eng{arrived} $\rightarrow$ past perfect : \eng{had already started}. (Quand nous sommes arrivés au studio, l'émission avait déjà commencé.)
\item Condition réelle : \eng{If + present} $\rightarrow$ futur avec \eng{will} : \eng{will be}. (Si le gouvernement investit dans la radio locale, plus de villageois seront informés.)
\item \eng{Since 2019} = durée jusqu'à maintenant $\rightarrow$ present perfect : \eng{has worked} (ou \eng{has been working}). (Elle travaille pour ce journal depuis 2019.)
\end{enumerate}
\textbf{Conclusion.} Le marqueur temporel est presque toujours la clé : entraîne-toi à le repérer avant de conjuguer.
\end{exoresolu}

\courssub{Savoir-faire 3 : Réussir les transformations de phrases (structures clés)}

\begin{methode}[Transformer une phrase sans changer le sens]
Les transformations classiques au Bac portent sur quatre structures :
\begin{enumerate}
\item \textbf{Discours direct $\rightarrow$ discours indirect} : recule les temps (\eng{am} $\rightarrow$ \eng{was} ; \eng{will} $\rightarrow$ \eng{would} ; \eng{have seen} $\rightarrow$ \eng{had seen}), adapte les pronoms et les marqueurs (\eng{now} $\rightarrow$ \eng{then} ; \eng{today} $\rightarrow$ \eng{that day} ; \eng{here} $\rightarrow$ \eng{there}).
\item \textbf{Actif $\rightarrow$ passif} : complément d'objet devient sujet ; verbe = \eng{be} (au temps du verbe actif) + participe passé ; agent introduit par \eng{by} (souvent omis).
\item \textbf{Comparatifs et superlatifs} : adjectif court + \eng{-er / the -est} ; adjectif long : \eng{more / the most + adjectif}.
\item \textbf{Reformulations imposées} (\eng{Rewrite using...}) : respecte le mot donné (\eng{too...to}, \eng{enough to}, \eng{unless} = \eng{if...not}, \eng{wish + past} pour un regret du présent).
\end{enumerate}
\textbf{Réflexe :} après la transformation, vérifie que le sens est identique et que la phrase est grammaticalement complète.
\end{methode}

\begin{exoresolu}[Grammar — Sentence transformation]
\textbf{Énoncé.} Rewrite each sentence as indicated, without changing the meaning.
\begin{enumerate}
\item \eng{The journalist said: “I will publish the article tomorrow.”} (Reported speech)
\item \eng{Millions of people watch this TV programme.} (Passive voice)
\item \eng{The exam was so difficult that many students failed.} (Rewrite using \eng{too...to})
\item \eng{If you don't read newspapers, you won't be well informed.} (Rewrite using \eng{unless})
\end{enumerate}
\tcblower
\textbf{Raisonnement et réponses.}
\begin{enumerate}
\item Recul des temps : \eng{will} $\rightarrow$ \eng{would} ; \eng{tomorrow} $\rightarrow$ \eng{the next day} ; \eng{I} $\rightarrow$ \eng{he/she}. Réponse : \eng{The journalist said (that) he would publish the article the next day.} (Le journaliste a dit qu'il publierait l'article le lendemain.)
\item Objet \eng{this TV programme} devient sujet ; \eng{watch} (present simple) $\rightarrow$ \eng{is watched} : \eng{This TV programme is watched by millions of people.} (Cette émission est regardée par des millions de personnes.)
\item \eng{So difficult that... failed} = trop difficile pour réussir : \eng{The exam was too difficult for many students to pass.} (L'examen était trop difficile pour que beaucoup d'élèves réussissent.) Structure : \eng{too + adjectif + for + quelqu'un + to + BV}.
\item \eng{Unless} = \eng{if...not} : \eng{Unless you read newspapers, you won't be well informed.} (À moins de lire les journaux, tu ne seras pas bien informé.)
\end{enumerate}
\textbf{Conclusion.} Dans une transformation, un seul élément change ; tout le reste de la phrase doit rester stable et correct.
\end{exoresolu}

\courssub{Savoir-faire 4 : Consolider le vocabulaire des médias et de la communication}

\begin{methode}[Mémoriser et réutiliser le lexique thématique]
\begin{enumerate}
\item \textbf{Apprends le vocabulaire par familles} : la presse (\eng{a newspaper, a headline, an article, a reporter, an editor}), la télévision et la radio (\eng{a channel, a programme, a broadcast, a presenter, the audience}), Internet (\eng{a website, social media, to post, to share, a follower, fake news}).
\item \textbf{Apprends les collocations}, pas les mots isolés : \eng{to watch the news} (regarder les informations), \eng{to read a newspaper}, \eng{to broadcast a programme} (diffuser une émission), \eng{to check a source} (vérifier une source), \eng{to spread information} (répandre une information).
\item \textbf{Méfie-toi des faux amis} : \eng{actually} = en fait (pas « actuellement » $\rightarrow$ \eng{currently}) ; \eng{a magazine} = un magazine ; \eng{sensible} = raisonnable (pas « sensible » $\rightarrow$ \eng{sensitive}) ; \eng{an information} est incorrect (\eng{information} est indénombrable $\rightarrow$ \eng{a piece of information}).
\item \textbf{Réutilise le lexique dans tes productions} : au Bac, employer correctement trois ou quatre mots du thème (\eng{social media, fake news, to broadcast}) valorise ta copie.
\end{enumerate}
\end{methode}

\begin{popfigure}
\begin{tikzpicture}[
  mindmap,
  grow cyclic,
  every node/.style={concept, execute at begin node=\hskip0pt},
  concept color=popblue!80,
  level 1/.append style={level distance=4.5cm, sibling angle=120, concept color=popblue!60},
  level 2/.append style={level distance=3cm, sibling angle=60, concept color=poporange!70},
  text=white,
  font=\small\bfseries
]
\node {Media \& Communication}
  child [concept color=popgreen!70] {node {Press}
    child {node {newspaper}}
    child {node {headline}}
    child {node {article}}
    child {node {reporter}}
    child {node {editor}}
  }
  child [concept color=poppurple!70] {node {TV / Radio}
    child {node {channel}}
    child {node {programme}}
    child {node {broadcast}}
    child {node {presenter}}
    child {node {audience}}
  }
  child [concept color=poppink!70] {node {Digital}
    child {node {website}}
    child {node {social media}}
    child {node {to post / share}}
    child {node {follower}}
    child {node {fake news}}
  };
\end{tikzpicture}
\legende{Carte mentale : Vocabulaire des médias et de la communication (niveau Terminale A)}
\end{popfigure}

\begin{exoresolu}[Vocabulary — Media and Communication]
\textbf{Énoncé.} Complete the sentences with the correct word from the box: \eng{headline — broadcast — source — fake news — audience — reporter}
\begin{enumerate}
\item The \ldots\ldots asked the minister three difficult questions during the press conference.
\item You should always check the \ldots\ldots of information before sharing it on WhatsApp.
\item CRTV will \ldots\ldots the football match live on Saturday evening.
\item The \ldots\ldots of the newspaper was: “New school year begins in September.”
\item Many people believed the story, but it was only \ldots\ldots.
\item The programme is very popular: it attracts a large \ldots\ldots every evening.
\end{enumerate}
\tcblower
\textbf{Raisonnement et réponses.}
\begin{enumerate}
\item Celui qui pose des questions lors d'une conférence de presse est le journaliste : \eng{reporter}. (Le journaliste a posé trois questions difficiles au ministre.)
\item On vérifie l'origine d'une information : \eng{source}. (Tu devrais toujours vérifier la source d'une information avant de la partager.)
\item Diffuser un match en direct : \eng{broadcast} (verbe irrégulier : \eng{broadcast – broadcast – broadcast}). (CRTV diffusera le match en direct samedi soir.)
\item Le gros titre d'un journal : \eng{headline}. (Le titre du journal était : « La rentrée scolaire commence en septembre ».)
\item Une histoire fausse que beaucoup ont crue : \eng{fake news} (invariable). (Beaucoup de gens ont cru l'histoire, mais c'était seulement une fausse information.)
\item Le public d'une émission : \eng{audience}. (L'émission attire un large public chaque soir.)
\end{enumerate}
\textbf{Conclusion.} Pour choisir le bon mot, appuie-toi sur le contexte de la phrase et sur la collocation attendue (\eng{to broadcast a match}, \eng{a large audience}).
\end{exoresolu}

\begin{savaistu}[Le paysage médiatique au Cameroun]
Le Cameroun possède un paysage médiatique bilingue (français/anglais) unique en Afrique.
\begin{itemize}
\item \textbf{Public :} CRTV (Cameroon Radio Television) diffuse en français et en anglais (radio, TV, web). C'est le média de référence national.
\item \textbf{Privé :} De nombreuses chaînes TV (Canal 2 International, STV, Equinoxe TV) et radios privées (Radio Balafon, Nostalgie, Sweet FM) animent le débat public.
\item \textbf{Presse écrite :} Quotidiens comme \emph{Cameroon Tribune} (bilingue, public), \emph{The Guardian Post} (anglophone), \emph{Mutations}, \emph{Le Jour} (francophones).
\item \textbf{Numérique :} Explosion des médias en ligne (\emph{Actu Cameroun}, \emph{Journal du Cameroun}) et usage massif de WhatsApp/Facebook pour l'info, d'où le défi des \eng{fake news} et la loi sur la cybercriminalité (2010).
\end{itemize}
Au Bac, citer \eng{CRTV}, \eng{bilingualism} ou \eng{cybercrime law} montre une culture générale ancrée dans ton contexte.
\end{savaistu}

\courssub{Savoir-faire 5 : Rédiger une composition de type Bac (Writing)}

\begin{methode}[Rédiger une composition en 5 étapes]
\begin{enumerate}
\item \textbf{Analyse le sujet} (2 min) : quel type de texte ? (lettre, article, essai d'opinion). Qui écrit ? À qui ? Sur quel thème ? Combien de mots ?
\item \textbf{Fais un plan simple} (3 min) : introduction (poser le sujet), deux ou trois paragraphes d'idées (une idée = un paragraphe + un exemple), conclusion (ton avis ou une ouverture).
\item \textbf{Rédige} avec des connecteurs variés : \eng{First of all, Moreover, However, For example, In my opinion, To conclude}. Utilise les temps étudiés : present simple pour les généralités, present perfect pour les bilans, conditionnel pour les hypothèses.
\item \textbf{Soigne les phrases} : courtes, sujet + verbe + complément, dans cet ordre. Une phrase simple correcte vaut mieux qu'une phrase longue fautive.
\item \textbf{Relis} (3 min) : accords sujet-verbe, \eng{-s} à la 3e personne, temps, orthographe, ponctuation, nombre de mots respecté.
\end{enumerate}
\textbf{Structure d'opinion à retenir :} \eng{In my opinion, social media is useful because... However, we must be careful with fake news. To conclude, I believe that...}
\end{methode}

\begin{experience}[Expression orale — Débat / Role-play]
\textbf{Consigne.} En binômes, préparez un dialogue de 2 minutes sur le thème : \eng{“Should students be allowed to use smartphones in class for educational purposes?”}
\begin{enumerate}
\item \textbf{Student A (Pour) :} utilise \eng{In my opinion...}, \eng{For instance...}, \eng{Moreover...}, \eng{access to information}, \eng{digital skills}, \eng{online dictionaries}.
\item \textbf{Student B (Contre) :} utilise \eng{I disagree...}, \eng{On the one hand...}, \eng{However...}, \eng{distraction}, \eng{cheating}, \eng{cyberbullying}, \eng{addiction}.
\item \textbf{Conclusion commune :} \eng{To sum up, we agree that... but...}
\end{enumerate}
\textbf{Objectif :} mobiliser le vocabulaire du chapitre (\eng{source, fake news, social media, broadcast, check}) et les connecteurs logiques à l'oral.
\end{experience}

\begin{exoresolu}[Writing — Essay type Bac]
\textbf{Énoncé.} \emph{Write an essay of about 150 words on the following topic: “Social media does more harm than good to young people.” Do you agree?}
\tcblower
\textbf{Raisonnement.} Sujet d'opinion : il faut une position claire, deux arguments avec exemples, une nuance (\eng{however}) et une conclusion personnelle. Temps dominant : present simple. Connecteurs obligatoires pour la cohérence.

\textbf{Modèle de réponse (environ 150 mots).}

\eng{Nowadays, almost every young Cameroonian uses social media such as Facebook, WhatsApp and TikTok. In my opinion, these platforms do more good than harm, if they are used wisely.}

\eng{First of all, social media helps students to learn: they can watch lessons, read the news and practise English every day. Moreover, it allows young people to communicate quickly with their families, even in remote villages.}

\eng{However, social media can also be dangerous. Fake news spreads very fast, and some students spend too much time online instead of studying. For example, a student who chats all night cannot concentrate in class the next day.}

\eng{To conclude, I believe that social media is a powerful tool. Young people should use it to learn and to communicate, but they must always check their sources and limit their time online.}

\textbf{Conclusion.} La copie respecte le plan introduction / arguments / nuance / conclusion, emploie le lexique du chapitre (\eng{social media, fake news, sources}) et des connecteurs variés : c'est exactement ce que le correcteur récompense.
\end{exoresolu}

\courssub{Savoir-faire 6 : Gérer l'épreuve le jour de l'examen}

\begin{methode}[Stratégie d'examen : les réflexes gagnants]
\begin{enumerate}
\item \textbf{Répartis ton temps} : lis l'ensemble du sujet d'abord ; commence par les exercices que tu maîtrises ; garde au moins 15 minutes pour la composition et 5 minutes pour la relecture.
\item \textbf{Respecte les consignes} : nombre de mots, \eng{in your own words}, \eng{rewrite using...}. Une réponse hors consigne perd des points même si l'anglais est correct.
\item \textbf{Ne laisse aucune question vide} : une réponse approchante peut rapporter un demi-point ; une case vide rapporte zéro.
\item \textbf{En grammaire}, justifie mentalement chaque réponse par un marqueur (\eng{since} $\rightarrow$ present perfect) avant d'écrire.
\item \textbf{Relis en te posant trois questions} : le verbe est-il au bon temps ? Le sujet et le verbe s'accordent-ils ? L'ordre des mots est-il anglais (sujet – verbe – complément) ?
\end{enumerate}
\end{methode}

\begin{exoresolu}[Entraînement final — mini-épreuve complète]
\textbf{Énoncé.} (a) Complete: \eng{If I (have) more time, I would read more books.} (b) Put into the passive: \eng{The students wrote the article.} (c) Find the mistake and correct it: \eng{She is living in Yaoundé since 2020.} (d) Translate into English: « Les réseaux sociaux sont très populaires chez les jeunes. »
\tcblower
\textbf{Raisonnement et réponses.}
\begin{enumerate}
\item[(a)] Condition irréelle du présent : \eng{If + past, would + BV} $\rightarrow$ \eng{If I had more time, I would read more books.} (Si j'avais plus de temps, je lirais plus de livres.)
\item[(b)] \eng{Wrote} = simple past $\rightarrow$ passif : \eng{was/were + participe passé}. Sujet singulier : \eng{The article was written by the students.} (L'article a été écrit par les élèves.)
\item[(c)] \eng{Since 2020} impose le present perfect, pas le present continuous : \eng{She has lived in Yaoundé since 2020.} (Elle vit à Yaoundé depuis 2020.) Le verbe \eng{live} au present perfect exprime une situation qui dure.
\item[(d)] \eng{Social media is very popular with young people.} Attention : \eng{social media} se construit souvent au singulier dans ce type de phrase ; \eng{popular with} (et non \eng{popular for}) ; \eng{young people} (et non \eng{the youngs}).
\end{enumerate}
\textbf{Conclusion.} Quatre points de grammaire différents en quatre questions : à l'examen, chaque question teste une structure précise — identifie-la avant de répondre.
\end{exoresolu}

\begin{attention}[Les 5 erreurs qui coûtent des points au Bac]
\begin{enumerate}
\item \textbf{Le present perfect avec une date passée.} Ne dis jamais \eng{I have seen him yesterday} : dis \eng{I saw him yesterday}. Le present perfect réclame \eng{since, for, already, just, ever, never} — jamais un moment précis du passé.
\item \textbf{L'oubli du \eng{-s} à la 3e personne du present simple.} \eng{She watch} est faux : \eng{She watches}. C'est la faute la plus fréquente des copies francophones.
\item \textbf{Les faux amis.} \eng{Actually} signifie « en fait » (« actuellement » = \eng{currently}) ; \eng{to assist} signifie « aider » (« assister à » = \eng{to attend}) ; \eng{an information} est incorrect : \eng{information} est indénombrable, dis \eng{a piece of information}.
\item \textbf{Les calques du français et l'ordre des mots.} On dit \eng{I am 16 years old} (jamais \eng{I have 16 years}) ; \eng{to listen to the radio} (préposition \eng{to} obligatoire) ; l'adjectif se place \emph{avant} le nom : \eng{a black car}, jamais \eng{a car black}.
\item \textbf{Les fautes d'orthographe de mots fréquents.} \eng{wich} $\rightarrow$ \eng{which} ; \eng{because} (pas \eng{because of} devant une phrase) ; \eng{their / there / they're} ne sont pas interchangeables ; \eng{its} (possessif) et \eng{it's} (= \eng{it is}) non plus.
\end{enumerate}
\textbf{Réflexe final :} consacre toujours les cinq dernières minutes à traquer ces cinq erreurs dans ta copie.
\end{attention}

\newpage

\coursec{Exercices}

\courssub{Je vérifie mes connaissances}

\begin{exercice}[QCM : grammaire et vocabulaire]
Pour chaque question, choisis la bonne réponse (a, b ou c).
\begin{enumerate}
\item The journalist said that the story \ldots\ the previous day.
\begin{enumerate}
\item[a)] publishes \item[b)] had been published \item[c)] has been published
\end{enumerate}
\item I have been working for this radio station \ldots\ 2019.
\begin{enumerate}
\item[a)] for \item[b)] since \item[c)] during
\end{enumerate}
\item The news \ldots\ broadcast every evening at 8 p.m.
\begin{enumerate}
\item[a)] is \item[b)] are \item[c)] have been
\end{enumerate}
\item A person who presents the news on television is a \ldots
\begin{enumerate}
\item[a)] headline \item[b)] newsreader \item[c)] viewer
\end{enumerate}
\item She asked me \ldots\ I had read the article.
\begin{enumerate}
\item[a)] that \item[b)] if \item[c)] what
\end{enumerate}
\end{enumerate}
\end{exercice}

\begin{exercice}[Vrai ou faux, à justifier]
Dis si chaque affirmation est \eng{true} ou \eng{false}, puis justifie ta réponse par une règle ou un exemple.
\begin{enumerate}
\item \eng{The word “information” can be used in the plural in English.}
\item \eng{In reported speech, “will” usually becomes “would”.}
\item \eng{“Actually” means “actuellement” in English.}
\item \eng{We use the present perfect with a finished past time expression like “yesterday”.}
\item \eng{“The media” is followed by a plural verb.}
\end{enumerate}
\end{exercice}

\begin{exercice}[Vocabulaire : Media and Communication]
Complète chaque phrase avec le mot anglais qui convient, choisi dans la liste suivante : \eng{headline}, \eng{source}, \eng{broadcast}, \eng{audience}, \eng{editor}.
\begin{enumerate}
\item The \ldots\ of the newspaper decided not to publish the photo.
\item A good journalist always checks his or her \ldots\ before writing an article.
\item The programme attracted an \ldots\ of two million viewers.
\item The match will be \ldots\ live on national television tonight.
\item “Youths take over social media” was the main \ldots\ of the Daily Post yesterday.
\end{enumerate}
\end{exercice}

\begin{exercice}[Conjugaison à compléter]
Mets le verbe entre parenthèses au temps qui convient.
\begin{enumerate}
\item By the time we arrived at the studio, the interview \ldots\ (already / start).
\item Look! The presenter \ldots\ (read) the headlines right now.
\item They \ldots\ (never / see) such a large crowd at a press conference before.
\item While the reporter \ldots\ (interview) the minister, the camera suddenly stopped working.
\item Next year, our school \ldots\ (launch) its own radio programme.
\end{enumerate}
\end{exercice}

\courssub{Je m'entraîne}

\begin{exercice}[Traduction]
Traduis les phrases suivantes en anglais.
\begin{enumerate}
\item On m'a dit que le journal était publié chaque matin à Douala.
\item Le présentateur a demandé aux téléspectateurs de rester calmes.
\item Je regarde cette chaîne depuis cinq ans.
\item Les fausses informations se répandent très vite sur les réseaux sociaux.
\item Il m'a suggéré de lire l'article avant le débat.
\end{enumerate}
\end{exercice}

\begin{exercice}[Transformations : voix passive et discours rapporté]
\begin{enumerate}
\item Mets les phrases suivantes à la voix passive :
\begin{enumerate}
\item[a)] The editor rejected the young journalist's article.
\item[b)] Millions of people watched the presidential debate.
\item[c)] Hackers have attacked the newspaper's website.
\end{enumerate}
\item Mets les phrases suivantes au discours rapporté (commence par \eng{He/She said that\ldots} ou utilise le verbe introducteur approprié) :
\begin{enumerate}
\item[a)] The minister said: “We will improve internet access in rural areas.”
\item[b)] The teacher asked the students: “Have you checked your sources?”
\item[c)] The presenter said: “Do not share unverified information.”
\end{enumerate}
\end{enumerate}
\end{exercice}

\begin{exercice}[Texte à trous : choix du bon temps]
Conjugue les dix verbes entre parenthèses au temps qui convient (past simple, past continuous, present perfect ou futur).
\begin{textebox}[Election night on the radio]
Last night, thousands of listeners (1) \ldots\ (follow) the election results on the national radio. While the presenters (2) \ldots\ (announce) the first figures, people (3) \ldots\ (gather) in front of their radios in Yaoundé and Bamenda. The electoral commission (4) \ldots\ (publish) partial results at 9 p.m., and since then, the tension (5) \ldots\ (grow) steadily. So far, no official winner (6) \ldots\ (be) declared. A journalist who (7) \ldots\ (work) at the station for ten years said he (8) \ldots\ (never / see) such excitement. Tomorrow morning, the commission (9) \ldots\ (hold) a press conference, and the radio (10) \ldots\ (broadcast) it live.
\end{textebox}
\end{exercice}

\begin{exercice}[Article à compléter et correction d'erreurs]
\begin{enumerate}
\item Complète l'article suivant avec dix mots de la banque : \eng{headline}, \eng{broadcast}, \eng{reliable}, \eng{fake news}, \eng{to share}, \eng{audience}, \eng{editor}, \eng{source}, \eng{viral}, \eng{issue}.
\begin{textebox}[A lesson in responsible journalism]
Last week, a photo went (1) \ldots\ on social media in Douala. Many users rushed (2) \ldots\ it without checking the (3) \ldots\ . The next day, the (4) \ldots\ of a local newspaper explained that the picture was (5) \ldots\ . He reminded his (6) \ldots\ that only (7) \ldots\ information should be (8) \ldots\ . The newspaper's front-page (9) \ldots\ the following morning was: “Think before you click!” The (10) \ldots\ was discussed in many classrooms across the city.
\end{textebox}
\item Repère et corrige l'erreur dans chacune des phrases suivantes :
\begin{enumerate}
\item[a)] She gave me many useful informations about the press.
\item[b)] I study journalism since three years.
\item[c)] Actually, more young people prefer social media to television.
\item[d)] He suggested me to read the article carefully.
\item[e)] The reporter asked me where did I get the information.
\end{enumerate}
\end{enumerate}
\end{exercice}

\courssub{Je me prépare au Bac}

\begin{exercice}[Compréhension écrite type Bac]
Lis le texte suivant, puis réponds aux questions.

\begin{textebox}[Young Nigerians and the internet]
In Lagos, Nigeria's largest city, the smartphone has become an essential tool for young people. In cybercafés, schools and even on crowded buses, teenagers scroll through their screens to chat, learn and do business. For many of them, the internet is not a luxury but a gateway to the world.

Take Chiamaka, a seventeen-year-old student. Every evening, after finishing her homework, she watches educational videos to prepare for her exams. “The internet helps me understand lessons that were not clear in class,” she explains. Her younger brother Emeka uses his phone differently: he has created a small online shop where he sells phone accessories to his classmates.

However, this digital revolution has a dark side. Teachers complain that some students spend more time on entertainment than on their studies. Parents worry about fake news and online fraud, which have trapped many inexperienced users. Moreover, not everyone can afford a smartphone or a stable connection: in rural areas, many young people remain cut off from these opportunities.

The Nigerian government has promised to improve internet access in schools and to teach digital literacy. Experts believe that if young people learn to use the web responsibly, the internet could become one of the country's greatest strengths. The challenge, they say, is not to reject technology, but to master it.
\end{textebox}

\begin{enumerate}
\item \textbf{True or false?} Justify each answer with a quotation from the text. (5 marks)
\begin{enumerate}
\item[a)] Chiamaka uses the internet only for entertainment.
\item[b)] Emeka earns money thanks to his phone.
\item[c)] All young Nigerians have access to a good internet connection.
\item[d)] Parents are completely satisfied with their children's internet use.
\item[e)] Experts think technology should be rejected.
\end{enumerate}
\item \textbf{Choose the correct answer.} (3 marks)
\begin{enumerate}
\item[a)] The internet is described as a “gateway to the world” because it: (i) is very expensive (ii) opens young people to new opportunities (iii) replaces schools.
\item[b)] The government's promise concerns: (i) banning smartphones (ii) building more cybercafés (iii) improving internet access in schools.
\item[c)] The main idea of the text is that the internet: (i) is only dangerous (ii) has both advantages and risks (iii) should be forbidden to teenagers.
\end{enumerate}
\item \textbf{Vocabulary in context.} Find in the text words or expressions meaning: (3 marks)
\begin{enumerate}
\item[a)] to move through content on a screen (paragraph 1)
\item[b)] cheating people to get their money (paragraph 3)
\item[c)] the ability to use digital tools well (paragraph 4)
\end{enumerate}
\item \textbf{Language work.} (4 marks)
\begin{enumerate}
\item[a)] Put into the passive voice: “Teachers complain that some students spend more time on entertainment.” (transform: \eng{More time\ldots})
\item[b)] Report the sentence: Chiamaka said: “The internet helps me understand my lessons.”
\end{enumerate}
\end{enumerate}
\end{exercice}

\begin{exercice}[Traduction type Bac]
Traduis le passage suivant en anglais. (5 marks)

\begin{textebox}[Texte à traduire]
Les médias jouent un rôle important dans nos vies. Chaque matin, des millions de Camerounais écoutent la radio ou lisent les journaux pour s'informer. Cependant, tout le monde ne vérifie pas les informations avant de les partager. Les fausses nouvelles peuvent créer la peur et la confusion dans la population. C'est pourquoi les journalistes doivent toujours vérifier leurs sources, et les citoyens doivent apprendre à réfléchir avant de cliquer.
\end{textebox}
\end{exercice}

\begin{exercice}[Composition type Bac]
Rédige une composition d'environ \textbf{200 mots} sur le sujet suivant. Respecte le plan imposé : introduction (présentation du sujet et problématique), deux arguments personnels, un contre-argument nuancé, conclusion avec ta position claire. (10 marks)

\begin{textebox}[Essay topic]
“Television is still more powerful than social media.” Do you agree?
\end{textebox}

Avant de rédiger, note tes idées dans le tableau de préparation suivant :

\begin{center}
\begin{tabularx}{\linewidth}{|X|X|}
\hline
\rowcolor{popblueL} \textbf{Étape du plan} & \textbf{Mes idées (en anglais)} \\
\hline
Introduction : présenter le sujet et la problématique & \\
\hline
Argument 1 (ex. portée, crédibilité de la télévision) & \\
\hline
Argument 2 (ex. rapidité, interactivité des réseaux sociaux) & \\
\hline
Contre-argument et réponse & \\
\hline
Conclusion : ta position personnelle & \\
\hline
\end{tabularx}
\end{center}
\end{exercice}

\newpage

\coursec{Corrigés des exercices}

\begin{corrige}[Exercice 1 — QCM : grammaire et vocabulaire]
\begin{enumerate}
\item \textbf{b) had been published.} Discours rapporté avec verbe introducteur au passé (\eng{said}) : le past perfect passif (\eng{had been + participe passé}) exprime une action antérieure au moment de la parole. \eng{The journalist said that the story had been published the previous day.} « Le journaliste a dit que l'article avait été publié la veille. »
\item \textbf{b) since.} \eng{Since + point de départ} (2019) ; \eng{for} s'emploie avec une durée (\eng{for five years}) ; \eng{during} ne se combine pas avec le present perfect continu.
\item \textbf{a) is.} \eng{News} est un nom \textbf{singulier indénombrable} malgré son « s » final : \eng{The news is broadcast every evening at 8 p.m.} « Les informations sont diffusées chaque soir à 20 h. »
\item \textbf{b) newsreader.} Un \eng{newsreader} (ou \eng{news anchor}) est un présentateur de journal télévisé. Un \eng{headline} est un titre de article ; un \eng{viewer} est un téléspectateur.
\item \textbf{b) if.} Question fermée (oui/non) rapportée : \eng{asked + if/whether + sujet + verbe}. \eng{She asked me if I had read the article.} « Elle m'a demandé si j'avais lu l'article. »
\end{enumerate}
\textbf{Points de vigilance :} \eng{news}, \eng{information}, \eng{advice} sont indénombrables ; ne jamais écrire \eng{informations} ni \eng{an advice}.
\end{corrige}

\begin{corrige}[Exercice 2 — Vrai ou faux, à justifier]
\begin{enumerate}
\item \textbf{False.} \eng{Information} est indénombrable en anglais : il n'a pas de pluriel. On dit \eng{some information}, \eng{a piece of information}. Exemple : \eng{She gave me useful information.} (et non \eng{informations}).
\item \textbf{True.} Au discours rapporté, avec un verbe introducteur au passé, \eng{will} devient \eng{would} : \eng{He said: “I will come.”} $\rightarrow$ \eng{He said (that) he would come.}
\item \textbf{False.} Faux ami : \eng{actually} signifie « en fait, en réalité ». « Actuellement » se traduit par \eng{currently} ou \eng{at present}. Exemple : \eng{He actually lives in Buea.} « En fait, il vit à Buea. »
\item \textbf{False.} Le present perfect est incompatible avec un repère de temps passé révolu (\eng{yesterday}, \eng{last week}, \eng{in 2010}). On utilise le past simple : \eng{I watched the news yesterday.} (et non \eng{I have watched}).
\item \textbf{True.} \eng{The media} est le pluriel de \eng{medium} et prend un verbe au pluriel : \eng{The media play an important role.} « Les médias jouent un rôle important. »
\end{enumerate}
\end{corrige}

\begin{corrige}[Exercice 3 — Vocabulaire : Media and Communication]
\begin{enumerate}
\item \textbf{editor} — \eng{The editor of the newspaper decided not to publish the photo.} « Le rédacteur en chef du journal a décidé de ne pas publier la photo. »
\item \textbf{source} — \eng{A good journalist always checks his or her source before writing an article.} « Un bon journaliste vérifie toujours sa source avant d'écrire un article. »
\item \textbf{audience} — \eng{The programme attracted an audience of two million viewers.} « L'émission a attiré un public de deux millions de téléspectateurs. »
\item \textbf{broadcast} — \eng{The match will be broadcast live on national television tonight.} « Le match sera diffusé en direct à la télévision nationale ce soir. » (passif : \eng{will be + participe passé} ; participe passé de \eng{broadcast} = \eng{broadcast}.)
\item \textbf{headline} — \eng{“Youths take over social media” was the main headline of the Daily Post yesterday.} « “Les jeunes s'emparent des réseaux sociaux” était le gros titre du Daily Post hier. »
\end{enumerate}
\end{corrige}

\begin{corrige}[Exercice 4 — Conjugaison à compléter]
\begin{enumerate}
\item \textbf{had already started.} Past perfect : action antérieure à une autre action passée (\eng{by the time we arrived}). « Quand nous sommes arrivés au studio, l'interview avait déjà commencé. »
\item \textbf{is reading.} Present continuous (\eng{Look!}, \eng{right now}) : action en cours au moment de la parole. « Regarde ! Le présentateur lit les gros titres en ce moment même. »
\item \textbf{have never seen.} Present perfect avec \eng{before} : expérience vécue jusqu'à aujourd'hui. « Ils n'ont jamais vu une telle foule à une conférence de presse. »
\item \textbf{was interviewing.} Past continuous après \eng{while} : action longue interrompue par une action brève au past simple (\eng{stopped}). « Pendant que le journaliste interviewait le ministre, la caméra s'est soudain arrêtée de fonctionner. »
\item \textbf{will launch.} Futur simple avec \eng{next year}. « L'année prochaine, notre école lancera sa propre émission de radio. »
\end{enumerate}
\textbf{Points de vigilance :} ne pas confondre \eng{for + durée} et \eng{since + point de départ} ; le past perfect exige le participe passé (\eng{had started}, pas \eng{had start}).
\end{corrige}

\begin{corrige}[Exercice 5 — Traduction]
\begin{enumerate}
\item \eng{I was told that the newspaper was published every morning in Douala.} (Passif : \eng{I was told} ; au discours rapporté, le présent passif « est publié » recule au past simple passif \eng{was published}.)
\item \eng{The newsreader asked the viewers to stay calm.} (Structure : \eng{ask + quelqu'un + to + infinitif} ; \eng{newsreader} ou \eng{presenter}.)
\item \eng{I have been watching this channel for five years.} (Present perfect continu : action commencée dans le passé et qui continue ; \eng{for five years}.)
\item \eng{Fake news spreads very quickly on social media.} (\eng{Fake news} est singulier indénombrable $\rightarrow$ \eng{spreads} ; \eng{social media} sans article.)
\item \eng{He suggested that I (should) read the article before the debate.} ou \eng{He suggested reading the article before the debate.} (Jamais \eng{He suggested me to read} : \eng{suggest} ne se construit pas avec un objet indirect + infinitif.)
\end{enumerate}
\textbf{Points de vigilance :} « se répandent » peut donner l'impression d'un pluriel, mais \eng{news} est singulier ; « depuis cinq ans » = \eng{for five years} avec un present perfect, pas un présent simple.
\end{corrige}

\begin{corrige}[Exercice 6 — Transformations : voix passive et discours rapporté]
\textbf{1. Voix passive}
\begin{enumerate}
\item[a)] \eng{The young journalist's article was rejected by the editor.} (past simple passif : \eng{was + rejected}.)
\item[b)] \eng{The presidential debate was watched by millions of people.} (Le complément d'objet \eng{the presidential debate} devient sujet.)
\item[c)] \eng{The newspaper's website has been attacked by hackers.} (present perfect passif : \eng{has been + participe passé}.)
\end{enumerate}
\textbf{2. Discours rapporté}
\begin{enumerate}
\item[a)] \eng{The minister said (that) they would improve internet access in rural areas.} (\eng{will} $\rightarrow$ \eng{would} ; \eng{we} $\rightarrow$ \eng{they}.)
\item[b)] \eng{The teacher asked the students if they had checked their sources.} (Question fermée $\rightarrow$ \eng{if} ; present perfect $\rightarrow$ past perfect ; ordre sujet--verbe, sans point d'interrogation.)
\item[c)] \eng{The presenter told the audience not to share unverified information.} (Impératif négatif rapporté : \eng{tell + quelqu'un + not to + infinitif}.)
\end{enumerate}
\textbf{Points de vigilance :} dans une question rapportée, on ne fait jamais d'inversion : \eng{He asked me where I got the information} (et non \eng{where did I get}).
\end{corrige}

\begin{corrige}[Exercice 7 — Texte à trous : choix du bon temps]
\begin{enumerate}
\item \textbf{followed} — past simple (\eng{last night}) : « ont suivi ».
\item \textbf{were announcing} — past continuous après \eng{while} : « annonçaient ».
\item \textbf{gathered} — past simple : action ponctuelle : « se sont rassemblés ».
\item \textbf{published} — past simple avec repère précis (\eng{at 9 p.m.}) : « a publié ».
\item \textbf{has grown} — present perfect avec \eng{since then} : « n'a cessé de croître ».
\item \textbf{has been} — present perfect avec \eng{so far} : « jusqu'ici, aucun vainqueur officiel n'a été déclaré ».
\item \textbf{had worked} — past perfect : antériorité par rapport à \eng{said} : « qui travaillait à la station depuis dix ans » (on accepte aussi \eng{had been working}).
\item \textbf{had never seen} — past perfect au discours rapporté : « il n'avait jamais vu ».
\item \textbf{will hold} — futur (\eng{tomorrow morning}) : « tiendra ».
\item \textbf{will broadcast} — futur : « la diffusera en direct ».
\end{enumerate}
\textbf{Méthode :} repérer d'abord les marqueurs temporels (\eng{last night}, \eng{while}, \eng{since then}, \eng{so far}, \eng{tomorrow}), puis choisir le temps associé.
\end{corrige}

\begin{corrige}[Exercice 8 — Article à compléter et correction d'erreurs]
\textbf{1. Article à compléter}
\begin{enumerate}
\item \textbf{viral} (\eng{to go viral} = devenir viral)
\item \textbf{to share} (\eng{to rush to do something})
\item \textbf{source} (vérifier la source)
\item \textbf{editor} (le rédacteur en chef)
\item \textbf{fake news} (la photo était une fausse information)
\item \textbf{audience} (son public)
\item \textbf{reliable} (des informations fiables)
\item \textbf{broadcast} (passif : \eng{should be broadcast} = devraient être diffusées)
\item \textbf{headline} (le gros titre de une)
\item \textbf{issue} (la question/le sujet a été débattu)
\end{enumerate}
\textbf{2. Correction d'erreurs}
\begin{enumerate}
\item[a)] \eng{informations} $\rightarrow$ \eng{information} (indénombrable) : \eng{She gave me much useful information about the press.} (On peut aussi dire \eng{a lot of useful information} ; \eng{many} ne s'emploie pas avec un indénombrable.)
\item[b)] \eng{I study} $\rightarrow$ \eng{I have been studying / I have studied} : \eng{I have been studying journalism for three years.} (present perfect + \eng{for} ; \eng{since} s'emploie avec un point de départ : \eng{since 2022}.)
\item[c)] \eng{Actually} $\rightarrow$ \eng{Currently / At present} : \eng{Currently, more young people prefer social media to television.} (faux ami.)
\item[d)] \eng{He suggested me to read} $\rightarrow$ \eng{He suggested (that) I (should) read} ou \eng{He suggested reading} : \eng{He suggested that I read the article carefully.}
\item[e)] \eng{where did I get} $\rightarrow$ \eng{where I got} : \eng{The reporter asked me where I got the information.} (pas d'inversion dans la question indirecte.)
\end{enumerate}
\end{corrige}

\begin{corrige}[Exercice 9 — Compréhension écrite type Bac]
\textbf{1. True or false (5 marks — 1 par réponse justifiée)}
\begin{enumerate}
\item[a)] \textbf{False.} \eng{“Every evening, after finishing her homework, she watches educational videos to prepare for her exams.”} (Elle utilise internet pour ses études, pas seulement pour se divertir.)
\item[b)] \textbf{True.} \eng{“He has created a small online shop where he sells phone accessories to his classmates.”}
\item[c)] \textbf{False.} \eng{“Not everyone can afford a smartphone or a stable connection: in rural areas, many young people remain cut off from these opportunities.”}
\item[d)] \textbf{False.} \eng{“Parents worry about fake news and online fraud.”}
\item[e)] \textbf{False.} \eng{“The challenge, they say, is not to reject technology, but to master it.”}
\end{enumerate}
\textbf{2. Choose the correct answer (3 marks — 1 par item)}
\begin{enumerate}
\item[a)] \textbf{(ii)} opens young people to new opportunities.
\item[b)] \textbf{(iii)} improving internet access in schools.
\item[c)] \textbf{(ii)} has both advantages and risks.
\end{enumerate}
\textbf{3. Vocabulary in context (3 marks — 1 par item)}
\begin{enumerate}
\item[a)] \eng{to scroll through} (« faire défiler »).
\item[b)] \eng{online fraud} (« fraude en ligne »).
\item[c)] \eng{digital literacy} (« culture numérique »).
\end{enumerate}
\textbf{4. Language work (4 marks — 2 par item)}
\begin{enumerate}
\item[a)] \eng{More time is spent on entertainment than on studies (by some students).} (present simple passif : \eng{is + spent} ; participe passé irrégulier de \eng{spend} = \eng{spent}.)
\item[b)] \eng{Chiamaka said (that) the internet helped her understand her lessons.} (\eng{helps} $\rightarrow$ \eng{helped} ; \eng{me} $\rightarrow$ \eng{her}.)
\end{enumerate}
\textbf{Points de vigilance :} toute réponse \eng{true/false} non justifiée par une citation perd la moitié des points ; recopier la citation entre guillemets, sans la modifier.
\end{corrige}

\begin{corrige}[Exercice 10 — Traduction type Bac]
\textbf{Traduction modèle :}
\end{corrige}

\begin{exemplebox}[Version corrigée]
The media play an important role in our lives. Every morning, millions of Cameroonians listen to the radio or read newspapers to get informed. However, not everyone checks information before sharing it. Fake news can create fear and confusion among the population. That is why journalists must always check their sources, and citizens must learn to think before they click.
\end{exemplebox}
\textbf{Barème indicatif (5 marks) :} 1 point par phrase correctement traduite (sens + grammaire) ; demi-point retiré par erreur grave (temps, accord, faux ami).
\textbf{Points de vigilance :}
\begin{itemize}
\item \eng{The media play} (pluriel), pas \eng{plays}.
\item « pour s'informer » : \eng{to get informed} ou \eng{to keep informed} (pas \eng{to inform themselves}, lourd).
\item « Les fausses nouvelles » : \eng{fake news} (singulier : \eng{Fake news can create\ldots}).
\item « dans la population » : \eng{among the population}, pas \eng{in the population}.
\item « C'est pourquoi » : \eng{That is why} ou \eng{This is why} (pas \eng{It is why}).
\end{itemize}


\begin{corrige}[Exercice 11 — Composition type Bac]
\textbf{Barème indicatif (10 marks) :} respect du plan et de la longueur (2) ; pertinence des arguments et exemples (3) ; correction grammaticale (3) ; vocabulaire et connecteurs (2).
\textbf{Exemple de tableau de préparation rempli :}
\begin{center}
\begin{tabularx}{\linewidth}{|X|X|}
\hline
\rowcolor{popblueL} \textbf{Étape du plan} & \textbf{Idées} \\
\hline
Introduction & Both TV and social media shape opinions; which is more powerful today? \\
\hline
Argument 1 & TV reaches everyone, even in villages; it is more reliable because journalists check facts. \\
\hline
Argument 2 & Social media are faster and interactive: young people share and comment instantly. \\
\hline
Contre-argument & Social media spread fake news, but television can also be biased. \\
\hline
Conclusion & Television remains more powerful for trust, but social media win for speed. \\
\hline
\end{tabularx}
\end{center}
\textbf{Composition modèle (environ 200 mots) :}
\end{corrige}

\begin{exemplebox}[Model essay]
Nowadays, both television and social media play a major role in shaping public opinion. But which of the two is more powerful? This question deserves careful thought.

On the one hand, television still reaches almost every household, from Yaoundé to the smallest village. It is also more reliable: professional journalists check their sources before broadcasting the news. During elections, for example, millions of viewers trust the national channel for official results.

On the other hand, social media are faster and more interactive. Young people can share a video, comment on it and debate with friends within seconds. A single post can go viral and influence thousands of users in one evening.

Admittedly, social media sometimes spread fake news, which weakens their credibility. However, television is not perfect either: it can be biased or controlled by powerful interests.

In conclusion, I believe that television is still more powerful in terms of trust and credibility, while social media dominate in speed and interaction. The ideal citizen should therefore use both wisely, checking every piece of information before believing it.
\end{exemplebox}
\textbf{Points de vigilance :} utiliser des connecteurs (\eng{on the one hand}, \eng{admittedly}, \eng{however}, \eng{in conclusion}) ; éviter les faux amis (\eng{actually}, \eng{to assist} $\neq$ assister à) ; accorder \eng{social media} avec un verbe au pluriel ; rester dans la limite des 200 mots.


\newpage

\coursec{Fiche bilan}

\begin{popfigure}
\begin{tikzpicture}[mindmap, grow cyclic, every node/.style={concept}, concept color=popblue!60,
level 1/.append style={level distance=3.4cm, sibling angle=60},
level 2/.append style={level distance=2.2cm}]
\node[concept color=popblue!70, text=white, align=center, font=\bfseries] {General\\Revision\\(Tle A)}
  child[concept color=poporange!70] { node[align=center, text=white] {Reading\\comprehension} }
  child[concept color=popgreen!70] { node[align=center, text=white] {Verb tenses\\(synthesis)} }
  child[concept color=poppurple!70] { node[align=center, text=white] {Key structures\\(reported speech,\\passive, relatives)} }
  child[concept color=poppink!70] { node[align=center, text=white] {Vocabulary\\Media \&\\Communication} }
  child[concept color=popgold!80] { node[align=center, text=white] {Writing\\(Bac composition)} }
  child[concept color=popblue!50] { node[align=center, text=white] {Exam strategy\\and training} };
\end{tikzpicture}
\legende{Carte mentale de la leçon : les six axes de la révision générale de fin d'année.}
\end{popfigure}

Cette leçon de révision générale clôt le chapitre \emph{Media and Communication} et l'année de Terminale A. Objectif : mobiliser l'ensemble des acquis (temps verbaux, structures, vocabulaire) et s'entraîner dans les conditions de l'épreuve d'anglais du Baccalauréat : compréhension de l'écrit, grammaire, vocabulaire et composition.

\begin{propriete}[L'essentiel à retenir — General revision]
\textbf{1. Lexique clé — Media and Communication}
\begin{itemize}\setlength{\itemsep}{1pt}
  \item \eng{the press / a newspaper / a headline} : la presse / un journal / un titre (de journal)
  \item \eng{a journalist / a reporter / an editor} : un journaliste / un reporter / un rédacteur en chef
  \item \eng{to broadcast / a broadcast} : diffuser / une émission (verbe irrégulier : \eng{broadcast -- broadcast -- broadcast})
  \item \eng{social media / a post / to share / to follow} : les réseaux sociaux / une publication / partager / suivre
  \item \eng{fake news / misinformation / a source / reliable} : les fausses informations / la désinformation / une source / fiable (\eng{news} et \eng{information} sont indénombrables : \eng{The news is good.})
  \item \eng{an advertisement / to advertise} : une publicité / faire de la publicité
  \item \eng{an interview / to interview / an interviewee} : une interview / interviewer / la personne interviewée
  \item \eng{censorship / freedom of the press / bias / biased} : la censure / la liberté de la presse / un parti pris / partial
\end{itemize}

\textbf{2. Règles de grammaire à connaître par cœur}
\begin{center}
\begin{tabularx}{\linewidth}{|l|X|X|}\hline
\rowcolor{popblueL} \textbf{Point} & \textbf{Forme} & \textbf{Exemple} \\ \hline
Present simple & base verbale (+ s à la 3\textsuperscript{e} personne du singulier) & \eng{She reads the news every day.} \\ \hline
Present continuous & be + V-ing & \eng{He is watching TV now.} \\ \hline
Past simple & V-ed / 2\textsuperscript{e} colonne des verbes irréguliers & \eng{They broadcast the match yesterday.} \\ \hline
Past continuous & was/were + V-ing & \eng{I was listening when the phone rang.} \\ \hline
Present perfect & have/has + participe passé & \eng{She has just published an article.} \\ \hline
Future & will + BV ; be going to + BV & \eng{The station will close next year.} \\ \hline
Passive voice & be (au temps voulu) + participe passé & \eng{The news is read at 8 p.m.} \\ \hline
Reported speech & recul des temps + adaptation des pronoms & \eng{He said (that) he was tired.} \\ \hline
Relatives & who, which, whose, where, that & \eng{The reporter who interviewed me was kind.} \\ \hline
Conditionnel (type 2) & if + past simple, would + BV & \eng{If I had a phone, I would call you.} \\ \hline
Modaux & can, must, should, may + BV & \eng{You should check your sources.} \\ \hline
Comparatifs et superlatifs & -er / more ... than ; the -est / the most & \eng{Radio is cheaper than TV.} \\ \hline
\end{tabularx}
\end{center}

\textbf{3. Méthodes express}
\begin{itemize}\setlength{\itemsep}{1pt}
  \item \textbf{Compréhension de texte} : lire les questions d'abord, repérer les mots-clés, souligner les passages utiles, répondre en phrases complètes avec ses propres mots.
  \item \textbf{Grammaire} : identifier le temps ou la structure demandée, vérifier l'accord sujet-verbe, relire chaque phrase.
  \item \textbf{Vocabulaire} : compléter avec le mot exact (nature et sens), attention aux faux amis : \eng{actually} = en fait (pas « actuellement »), \eng{a magazine} = un magazine (pas « un magasin », qui se dit \eng{a shop}).
  \item \textbf{Composition} : respecter le nombre de mots, faire un plan simple (introduction, 2 ou 3 paragraphes, conclusion), utiliser des connecteurs (\eng{first, moreover, however, in conclusion}), laisser 5 minutes pour relire.
\end{itemize}
\end{propriete}

\begin{attention}[Pièges fréquents des francophones à l'examen]
\begin{itemize}\setlength{\itemsep}{2pt}
  \item \eng{news} est singulier : \eng{The news is surprising.} (« Les nouvelles sont surprenantes. ») — jamais \emph{the news are}.
  \item \eng{information} est indénombrable : \eng{a piece of information} pour « une information ».
  \item Pas de futur après \eng{if} dans une hypothèse : \eng{If it rains, we will stay at home.} (et non \emph{if it will rain}).
  \item Au discours indirect, recul obligatoire des temps : \eng{“I am busy,” she said.} $\rightarrow$ \eng{She said (that) she was busy.}
  \item \eng{actually} signifie « en fait, en réalité » ; « actuellement » se traduit par \eng{currently} ou \eng{at present}.
  \item Ordre des mots : l'adjectif précède le nom (\eng{a reliable source}), et l'adverbe de fréquence se place avant le verbe principal (\eng{She always checks her sources.}).
\end{itemize}
\end{attention}

\begin{aretenir}[Checklist avant le Bac]
Coche chaque point que tu maîtrises vraiment :
\begin{itemize}\setlength{\itemsep}{2pt}
  \item[$\square$] Je connais les formes et les emplois de tous les temps verbaux principaux.
  \item[$\square$] Je sais conjuguer les verbes irréguliers les plus fréquents (\eng{go -- went -- gone}, \eng{write -- wrote -- written}, \eng{speak -- spoke -- spoken}...).
  \item[$\square$] Je sais transformer une phrase active en phrase passive.
  \item[$\square$] Je sais rapporter une parole au discours indirect avec le recul des temps.
  \item[$\square$] Je choisis correctement le pronom relatif (\eng{who, which, whose, where}).
  \item[$\square$] Je connais les modaux et leur sens (obligation, conseil, capacité, permission).
  \item[$\square$] Je me méfie des faux amis et je vérifie l'ordre des mots anglais.
  \item[$\square$] Je maîtrise le vocabulaire des médias et de la communication.
  \item[$\square$] Je sais rédiger une composition structurée avec des connecteurs logiques.
  \item[$\square$] Je gère mon temps : je lis d'abord toutes les consignes et je réserve du temps pour relire.
  \item[$\square$] Je réponds aux questions de compréhension en phrases complètes, sans recopier le texte.
  \item[$\square$] Je soigne la présentation : écriture lisible, paragraphes clairs, orthographe vérifiée.
\end{itemize}
\end{aretenir}

\findecours

\end{document}
